% La identidad, la tolerancia y la convivencia — Espagnol 1ère A — Cours signé OnBuch+
% Programme officiel MINESEC. Compiler avec : tectonic EA11-identidad-tolerancia-convivencia.tex
\newcommand{\DOCMATIERE}{Espagnol}
\newcommand{\DOCNIVEAU}{1ère A}
\newcommand{\DOCMODULE}{Módulo 3 : Citoyenneté}
\newcommand{\DOCLECON}{EA11}
\newcommand{\DOCTITRE}{La identidad, la tolerancia y la convivencia}
% OnBuch+ — préambule « cours signés » ENRICHI (compilé avec tectonic / XeLaTeX)
% Système visuel « Pop » d'OnBuch+ : fond crème (jamais blanc), bordures encre
% épaisses, ombres « dures » décalées, titres Archivo Black en capitales.
% Version MATHÉMATIQUES : graphes (pgfplots/tikz), boîtes pédagogiques
% (définition, propriété, méthode, attention, le savais-tu, exercice résolu,
% exercice, TP, bilan), page de garde et signature OnBuch+ ancrée en bas.
%
% Macros attendues dans main.tex AVANT \input{preamble} :
%   \DOCMATIERE  \DOCNIVEAU  \DOCMODULE  \DOCLECON (numéro)  \DOCTITRE
\documentclass[11pt]{article}

\usepackage{fontspec}
\usepackage[spanish,french]{babel} % français = langue principale ; l'espagnol via \esp{..} / espagnol
\usepackage[a4paper,margin=2cm,top=3.4cm,bottom=2.8cm,headheight=1.4cm,footskip=1.3cm]{geometry}
\usepackage{amsmath,amssymb,amsfonts}
\usepackage{mathtools} % \xmapsto, \coloneqq... souvent utilisés par les modèles
\usepackage{cancel} % simplifications barrées (bilans de forces)
% \abs{z} (module, valeur absolue) et \norm{u} (norme), utilisés par les modèles
\providecommand{\abs}{}\let\abs\relax\DeclarePairedDelimiter{\abs}{\lvert}{\rvert}
\providecommand{\norm}{}\let\norm\relax\DeclarePairedDelimiter{\norm}{\lVert}{\rVert}
\usepackage[table]{xcolor} % [table] : fournit aussi \rowcolors (alternance de lignes)
\usepackage{pagecolor}
\usepackage{tikz}
\usetikzlibrary{arrows.meta,positioning,calc,shapes.geometric,shapes.misc,shapes.arrows,decorations.pathmorphing,decorations.pathreplacing,decorations.markings,patterns,shadows,mindmap,trees,fit,backgrounds,matrix,intersections,angles,quotes}
\usepackage{pgfplots}
\usepackage[version=4]{mhchem} % équations nucléaires \ce{^{235}_{92}U}
\usepackage[european,siunitx]{circuitikz} % schémas électriques
\pgfplotsset{compat=1.18}
\usepgfplotslibrary{fillbetween,groupplots} % aires entre courbes (calcul intégral)
\usepackage{enumitem}
\usepackage{fancyhdr}
% [most] charge toutes les bibliothèques tcolorbox (listings, minted, raster,
% documentation...) et gonfle inutilement la mémoire TeX sur un document long
% (~300 boîtes) : on ne charge que ce qui est réellement utilisé.
\usepackage[skins,breakable]{tcolorbox}
\usepackage{graphicx}
\usepackage{colortbl}
\usepackage{booktabs}
\usepackage{tabularx}
\usepackage{diagbox} % case d'en-tête diagonale des tableaux à double entrée
% Colonnes extensibles centrée / gauche / droite (C, L, R), souvent utilisées par les modèles
\newcolumntype{C}{>{\centering\arraybackslash}X}
\newcolumntype{L}{>{\raggedright\arraybackslash}X}
\newcolumntype{R}{>{\raggedleft\arraybackslash}X}
\usepackage{array}
\usepackage{multicol}
\usepackage{siunitx}
% parse-numbers=false : accepte aussi \SI{2,5 \times 10^{-3}}{...}
\sisetup{locale=FR,output-decimal-marker={,},group-minimum-digits=4,parse-numbers=false}
\DeclareSIUnit\ohm{\ensuremath{\symup{\Omega}}} % U+2126 absent de la police
\DeclareSIUnit\division{div} % réglages d'oscilloscope (ms/div, V/div)
\DeclareSIUnit\tr{tr} % tours (vitesse de rotation, ex. \SI{1500}{\tr\per\minute})
\DeclareSIUnit\molar{M} % concentration molaire (ex. \SI{3}{\molar}, \SI{1}{\micro\molar})
\DeclareSIUnit\an{an} % année (le modèle confond parfois avec \year, un compteur TeX) — ex. \SI{5}{\kilo\meter\per\an}
\DeclareSIUnit\FCFA{FCFA} % franc CFA (ex. \SI{500}{\FCFA\per\kilo\gram})
\DeclareSIUnit\degreeBrix{\textdegree Brix} % teneur en sucre d'un sirop/jus (réfractométrie)
\DeclareSIUnit\month{mois} % le modèle utilise parfois \month, un compteur TeX, comme unité de durée
\usepackage{qrcode}
% \degree : commande fréquemment utilisée par les modèles pour un angle ou une
% température en dehors de \SI{}{} (non fournie par les paquets chargés).
\providecommand{\degree}{^{\circ}}
% \qed (fin de démonstration), utilisé par les modèles sans amsthm
\providecommand{\qed}{\hfill$\square$}
% \barre : double barre des valeurs interdites dans un tableau de variations (tkz-tab)
\providecommand{\barre}{\ensuremath{\|}}
% environnement proof (amsthm non chargé) utilisé par les modèles
\ifdefined\proof\else\newenvironment{proof}[1][Démonstration]{\par\noindent\textit{#1.}\ }{\qed\par}\fi
\usepackage{lastpage}
\usepackage{hyperref}
\hypersetup{hidelinks}

% ============================================================
% Palette « Pop » OnBuch+ — mêmes hex que la classe P (plus_ui.dart)
% ============================================================
\definecolor{poporange}{HTML}{F59321}
\definecolor{popdark}{HTML}{E07A0C}
\definecolor{popcream}{HTML}{FFF6E8}
\definecolor{popink}{HTML}{1C1714}
\definecolor{poppurple}{HTML}{7A5AE0}
\definecolor{popgreen}{HTML}{2E9E6A}
\definecolor{poppink}{HTML}{E0457B}
\definecolor{popblue}{HTML}{2F7DE1}
\definecolor{popgold}{HTML}{C98A12}
\definecolor{popmuted}{HTML}{8A7F76}
\definecolor{popline}{HTML}{EADFD2}
\definecolor{popgray}{HTML}{8A7F76}
\colorlet{popgrey}{popgray}
% Alias tolérants (le modèle nomme parfois les couleurs différemment)
\colorlet{popwhite}{white}
\colorlet{popblack}{popink}
\colorlet{popred}{poppink}
\colorlet{popyellow}{popgold}
% Teintes claires (fonds de tableaux/encadrés) — le modèle nomme parfois
% les couleurs avec un suffixe L (ex. popblueL) sans les avoir définies.
\colorlet{poporangeL}{poporange!20}
\colorlet{popdarkL}{popdark!20}
\colorlet{poppurpleL}{poppurple!20}
\colorlet{popgreenL}{popgreen!20}
\colorlet{poppinkL}{poppink!20}
\colorlet{popblueL}{popblue!20}
\colorlet{popgoldL}{popgold!20}
\colorlet{popredL}{popred!20}
\colorlet{popgrayL}{popgray!20}
% Teintes claires pour fonds de tableaux / zones
\colorlet{popblueL}{popblue!10!white}
\colorlet{poporangeL}{poporange!14!white}
\colorlet{popgreenL}{popgreen!12!white}
\colorlet{poppurpleL}{poppurple!12!white}
\colorlet{poppinkL}{poppink!12!white}
\colorlet{popgoldL}{popgold!14!white}

\pagecolor{popcream}

% ============================================================
% Typographie
% ============================================================
\defaultfontfeatures{Ligatures=TeX}
\setmainfont{PlusJakartaSans-Medium.ttf}[Path=../../../fonts/,BoldFont=PlusJakartaSans-Bold.ttf]
% Maths en sans-serif (Fira Math) avec chiffres et lettres droites de Plus
% Jakarta, pour que les formules se fondent dans le texte.
\usepackage[math-style=ISO,bold-style=ISO]{unicode-math}
\setmathfont{FiraMath-Regular.otf}[Path=../../../fonts/]
\setmathfont{PlusJakartaSans-Medium.ttf}[Path=../../../fonts/,range={up/num,up/{Latin,latin},\mathup}]
\usepackage{icomma} % virgule décimale sans espace en mode math
% Caractères mathématiques Unicode absents des polices de texte (Plus Jakarta,
% Archivo Black) : rendus par leur équivalent mathématique, en texte comme en
% formule — sinon ils s'affichent en carré vide (ex. « Arithmétique dans ℤ »).
\usepackage{newunicodechar}
\newunicodechar{ℝ}{\ensuremath{\mathbb{R}}}
\newunicodechar{ℤ}{\ensuremath{\mathbb{Z}}}
\newunicodechar{ℂ}{\ensuremath{\mathbb{C}}}
\newunicodechar{ℕ}{\ensuremath{\mathbb{N}}}
\newunicodechar{ℚ}{\ensuremath{\mathbb{Q}}}
\newunicodechar{𝔻}{\ensuremath{\mathbb{D}}}
\newunicodechar{α}{\ensuremath{\alpha}}
\newunicodechar{β}{\ensuremath{\beta}}
\newunicodechar{γ}{\ensuremath{\gamma}}
\newunicodechar{δ}{\ensuremath{\delta}}
\newunicodechar{ε}{\ensuremath{\varepsilon}}
\newunicodechar{θ}{\ensuremath{\theta}}
\newunicodechar{λ}{\ensuremath{\lambda}}
\newunicodechar{μ}{\ensuremath{\mu}}
\newunicodechar{σ}{\ensuremath{\sigma}}
\newunicodechar{τ}{\ensuremath{\tau}}
\newunicodechar{φ}{\ensuremath{\varphi}}
\newunicodechar{ψ}{\ensuremath{\psi}}
\newunicodechar{ω}{\ensuremath{\omega}}
\newunicodechar{Σ}{\ensuremath{\Sigma}}
% majuscules produites par \MakeUppercase dans les titres (α-aminés -> Α)
\newunicodechar{Α}{\ensuremath{\alpha}}
\newunicodechar{Β}{\ensuremath{\beta}}
\newunicodechar{Γ}{\ensuremath{\Gamma}}
\newunicodechar{Δ}{\ensuremath{\Delta}}
\newunicodechar{Ε}{\ensuremath{\varepsilon}}
\newunicodechar{Θ}{\ensuremath{\Theta}}
\newunicodechar{Λ}{\ensuremath{\Lambda}}
\newunicodechar{Μ}{\ensuremath{\mu}}
\newunicodechar{Τ}{\ensuremath{\tau}}
\newunicodechar{Φ}{\ensuremath{\Phi}}
\newunicodechar{Ψ}{\ensuremath{\Psi}}
\newunicodechar{Ω}{\ensuremath{\Omega}}
\newunicodechar{↦}{\ensuremath{\mapsto}}
\newunicodechar{∈}{\ensuremath{\in}}
\newunicodechar{∉}{\ensuremath{\notin}}
\newunicodechar{∧}{\ensuremath{\wedge}}
\newunicodechar{⇔}{\ensuremath{\Leftrightarrow}}
\newunicodechar{⟺}{\ensuremath{\Longleftrightarrow}}
\newunicodechar{⇒}{\ensuremath{\Rightarrow}}
\newunicodechar{⟹}{\ensuremath{\Longrightarrow}}
\newunicodechar{∘}{\ensuremath{\circ}}
\newunicodechar{∪}{\ensuremath{\cup}}
\newunicodechar{∩}{\ensuremath{\cap}}
\newunicodechar{≡}{\ensuremath{\equiv}}
\newunicodechar{⊂}{\ensuremath{\subset}}
\newunicodechar{⊥}{\ensuremath{\perp}}
\newunicodechar{∃}{\ensuremath{\exists}}
\newunicodechar{∀}{\ensuremath{\forall}}
\newunicodechar{∛}{\ensuremath{\sqrt[3]{\,}}}
\newunicodechar{∜}{\ensuremath{\sqrt[4]{\,}}}
\newunicodechar{⃗}{\ensuremath{{}^{\to}}}
\newunicodechar{ⁿ}{\ifmmode^{n}\else\textsuperscript{\ensuremath{n}}\fi}
\newunicodechar{ˣ}{\ifmmode^{x}\else\textsuperscript{\ensuremath{x}}\fi}
\newunicodechar{ᵇ}{\ifmmode^{b}\else\textsuperscript{\ensuremath{b}}\fi}
\newunicodechar{ʳ}{\ifmmode^{r}\else\textsuperscript{\ensuremath{r}}\fi}
\newunicodechar{ᵘ}{\ifmmode^{u}\else\textsuperscript{\ensuremath{u}}\fi}
\newunicodechar{ʸ}{\ifmmode^{y}\else\textsuperscript{\ensuremath{y}}\fi}
\newunicodechar{ᵃ}{\ifmmode^{a}\else\textsuperscript{\ensuremath{a}}\fi}
\newunicodechar{ᵉ}{\ifmmode^{e}\else\textsuperscript{\ensuremath{e}}\fi}
\newunicodechar{ᵏ}{\ifmmode^{k}\else\textsuperscript{\ensuremath{k}}\fi}
\newunicodechar{ᵖ}{\ifmmode^{p}\else\textsuperscript{\ensuremath{p}}\fi}
\newunicodechar{ᴺ}{\ifmmode^{N}\else\textsuperscript{\ensuremath{N}}\fi}
\newunicodechar{ᶜ}{\ifmmode^{c}\else\textsuperscript{\ensuremath{c}}\fi}
\newunicodechar{ʰ}{\ifmmode^{h}\else\textsuperscript{\ensuremath{h}}\fi}
\newunicodechar{ᵠ}{\ifmmode^{q}\else\textsuperscript{\ensuremath{q}}\fi}
\newunicodechar{ₙ}{\ifmmode_{n}\else\textsubscript{\ensuremath{n}}\fi}
\newunicodechar{ₐ}{\ifmmode_{a}\else\textsubscript{\ensuremath{a}}\fi}
\newunicodechar{ᵢ}{\ifmmode_{i}\else\textsubscript{\ensuremath{i}}\fi}
\newunicodechar{ᵣ}{\ifmmode_{r}\else\textsubscript{\ensuremath{r}}\fi}
\newunicodechar{ₖ}{\ifmmode_{k}\else\textsubscript{\ensuremath{k}}\fi}
\newunicodechar{ᵦ}{\ifmmode_{\beta}\else\textsubscript{\ensuremath{\beta}}\fi}
\newunicodechar{ₓ}{\ifmmode_{x}\else\textsubscript{\ensuremath{x}}\fi}
\newfontfamily\popdisplay{ArchivoBlack-Regular.ttf}[Path=../../../fonts/]
\newfontfamily\poplogo{SpaceGrotesk-Bold.ttf}[Path=../../../fonts/]
\newfontfamily\popsemi{PlusJakartaSans-SemiBold.ttf}[Path=../../../fonts/]
\newfontfamily\popxbold{PlusJakartaSans-ExtraBold.ttf}[Path=../../../fonts/]
\newfontfamily\popxboldls{PlusJakartaSans-ExtraBold.ttf}[Path=../../../fonts/,LetterSpace=3.0]

\setlist[enumerate]{leftmargin=1.7em,itemsep=5pt,topsep=4pt}
\setlist[itemize]{leftmargin=1.7em,itemsep=4pt,topsep=4pt,label={\color{poporange}\textbullet}}
\setlength{\parindent}{0pt}
\setlength{\parskip}{5pt}
\renewcommand{\arraystretch}{1.25}

% pgfplots : style Pop par défaut
\pgfplotsset{
  every axis/.append style={axis line style={popink,line width=1pt},tick style={popink},
    grid style={popline,line width=0.5pt},label style={font=\small},tick label style={font=\footnotesize}},
  popcurve/.style={poporange,line width=1.8pt,no markers,smooth},
}
% Même style disponible dans un \draw TikZ simple (hors pgfplots)
\tikzset{popcurve/.style={poporange,line width=1.8pt}}
\tikzset{popgrid/.style={popline,very thin}}
\colorlet{popcurve}{poporange} % le nom du style est parfois utilisé comme couleur

% ============================================================
% Pill / squiggle
% ============================================================
\newcommand{\poppill}[3]{%
  \tikz[baseline=-0.55ex]\node[draw=popink,line width=1.2pt,rounded corners=2.6pt,%
    fill=#2,inner xsep=6.5pt,inner ysep=3pt]{%
    \popxboldls\fontsize{7}{7}\selectfont\color{#3}\MakeUppercase{#1}};%
}
\newcommand{\popsquiggle}[1]{%
  \tikz[baseline=-0.4ex]{
    \draw[#1,line width=2.6pt,line cap=round]
      plot [smooth] coordinates {(0,0.09) (0.34,0.01) (0.68,0.21) (1.02,0.01) (1.36,0.10)};
  }%
}
% Mot-clé surligné (marqueur orange sous le texte)
\newcommand{\cle}[1]{{\popxbold\color{popdark}#1}}

% ============================================================
% En-tête / pied
% ============================================================
\pagestyle{fancy}
\fancyhf{}
\renewcommand{\headrulewidth}{0pt}
\renewcommand{\footrulewidth}{0pt}
\lhead{\raisebox{-0.15ex}{\poplogo\fontsize{13}{13}\selectfont\color{popink}OnBuch\color{poporange}+}}
\rhead{\poppill{\DOCMATIERE\ \textbullet\ \DOCNIVEAU\ \textbullet\ Leçon \DOCLECON}{white}{popink}}
\lfoot{\popxbold\fontsize{7.6}{7.6}\selectfont\color{popmuted}\MakeUppercase{Cours signé OnBuch+}\ \textbullet\ {\popsemi\DOCTITRE}}
\rfoot{\tikz[baseline=-0.6ex]\node[rounded corners=8.5pt,fill=popink,minimum height=17pt,minimum width=17pt,inner xsep=5pt,inner sep=0]{\popxbold\fontsize{8}{8}\selectfont\color{popcream}\thepage\,/\,\pageref*{LastPage}};}

% ============================================================
% Titres
% ============================================================
\newcommand{\chaptitre}[2]{%
  \begin{tcolorbox}[enhanced,colback=poporange,colframe=popink,boxrule=2.6pt,arc=11pt,
      drop shadow=popink,left=15pt,right=15pt,top=12pt,bottom=11pt,before skip=3pt,after skip=18pt]
    \poppill{#2}{popink}{popcream}\par\vspace{9pt}
    {\raggedright\hyphenpenalty=10000\exhyphenpenalty=10000\popdisplay\fontsize{22}{24}\selectfont\color{popcream}\MakeUppercase{#1}\par}
  \end{tcolorbox}
}
% Section numérotée : pastille orange avec le numéro + titre Archivo
\newcounter{popsec}
\newcommand{\coursec}[1]{%
  \stepcounter{popsec}\setcounter{popsub}{0}%
  \par\vspace{16pt}\noindent
  \begin{minipage}{\linewidth}
  \tikz[baseline=-0.7ex]\node[draw=popink,line width=1.6pt,fill=poporange,rounded corners=4pt,minimum size=22pt,inner sep=0]{\popdisplay\fontsize{12}{12}\selectfont\color{popcream}\thepopsec};%
  \hspace{8pt}{\popdisplay\fontsize{15}{17}\selectfont\color{popink}\MakeUppercase{#1}}\par
  \vspace{4pt}\noindent{\color{poporange}\rule{36pt}{3.6pt}}
  \end{minipage}\par\nopagebreak\vspace{8pt}
}
\newcounter{popsub}
\newcommand{\courssub}[1]{%
  \stepcounter{popsub}%
  \par\vspace{9pt}\noindent{\popxbold\fontsize{11.5}{13}\selectfont\color{popink}{\color{poporange}\thepopsec.\thepopsub}\hspace{6pt}#1}\par\nopagebreak\vspace{4pt}
}

% ============================================================
% Boîtes pédagogiques — même recette : fond blanc, bordure épaisse colorée,
% ombre dure encre, badge plein. Titre optionnel [#1].
% ============================================================
\tcbset{popbase/.style={breakable,enhanced,colback=white,boxrule=2.3pt,arc=9pt,
  shadow={3pt}{-3pt}{0pt}{popink},left=14pt,right=14pt,top=11pt,bottom=11pt,before skip=11pt,after skip=15pt,
  fontupper=\popsemi\fontsize{10.6}{14.5}\selectfont\color{popink},
  fontlower=\popsemi\fontsize{10.6}{14.5}\selectfont\color{popink}}}
% Coche dessinée (le glyphe ✓ n'existe pas dans les polices du document)
\newcommand{\popcheck}[1]{\tikz[baseline=-0.5ex]\draw[#1,line width=1.6pt,line cap=round,line join=round](-0.13,0.0)--(-0.04,-0.09)--(0.13,0.1);}
\providecommand{\coche}{\popcheck{popgreen}}
\providecommand{\resultat}[1]{\par\smallskip\noindent\fcolorbox{popgreen}{popgreenL}{\parbox{\dimexpr\linewidth-2\fboxsep-2\fboxrule\relax}{\textbf{Résultat.} #1}}\par\smallskip}
\newcommand{\popboxhead}[3]{\poppill{#1}{#2}{white}\ifx\relax#3\relax\else\hspace{6pt}{\popxbold\fontsize{10.6}{12}\selectfont\color{popink}#3}\fi\par\vspace{7pt}}

\newtcolorbox{aretenir}[1][]{popbase,colframe=popgreen,before upper={\popboxhead{À retenir}{popgreen}{#1}}}
% Texte en espagnol (document de lecture, dialogue, poème, extrait) : cadre
% d'étude. Un titre optionnel (source, auteur) s'affiche dans l'en-tête.
\newtcolorbox{textebox}[1][]{popbase,colframe=popblue,before upper={\popboxhead{Texto}{popblue}{#1}\begin{otherlanguage}{spanish}},after upper={\end{otherlanguage}}}
% Marques ✓ / ✗ (forme correcte / fautive) souvent utilisées par les modèles
\providecommand{\cross}{\textcolor{poppink}{\ensuremath{\times}}}
\providecommand{\xmark}{\cross}\providecommand{\crossmark}{\cross}\providecommand{\cmark}{\ensuremath{\checkmark}}
% Ligne de réponse / justification à compléter (inventée par les modèles)
\providecommand{\justification}[1]{\par\noindent\textit{Justification~:} \dotfill\par}
% \textipa (package tipa non chargé) : la police ne couvre pas l'API -> simple passage
\providecommand{\textipa}[1]{#1}
% Soulignement (ulem non chargé) : \uline, \ul
\providecommand{\uline}[1]{\underline{#1}}\providecommand{\ul}[1]{\underline{#1}}
% Ordinaux français écrits en macro par les modèles : 1\ière, 3\ième
\newcommand{\ière}{\textsuperscript{re}}\newcommand{\ième}{\textsuperscript{e}}
% Mot ou phrase en espagnol à l'intérieur d'un paragraphe français
\newcommand{\esp}[1]{\ifmmode\text{\textit{#1}}\else\begingroup\selectlanguage{spanish}\itshape #1\endgroup\fi}
\newtcolorbox{exemplebox}[1][]{popbase,colframe=poporange,before upper={\popboxhead{Exemple}{poporange}{#1}}}
\newtcolorbox{definition}[1][]{popbase,colframe=popblue,before upper={\popboxhead{Définition}{popblue}{#1}}}
\newtcolorbox{propriete}[1][]{popbase,colframe=poppurple,before upper={\popboxhead{Propriété}{poppurple}{#1}}}
\newtcolorbox{methode}[1][]{popbase,colframe=popgold,before upper={\popboxhead{Méthode}{popgold}{#1}}}
\newtcolorbox{attention}[1][]{popbase,colframe=poppink,before upper={\popboxhead{Attention}{poppink}{#1}}}
\newtcolorbox{experience}[1][]{popbase,colframe=popblue,colback=popblueL,before upper={\popboxhead{Expérience}{popblue}{#1}}}
% « Le savais-tu ? » : carte violette pleine (contexte, culture, Cameroun)
\newtcolorbox{savaistu}[1][]{popbase,colback=poppurple,colframe=popink,
  fontupper=\popsemi\fontsize{10.4}{14.2}\selectfont\color{white},
  before upper={\poppill{Le savais-tu\,?}{white}{poppurple}\ifx\relax#1\relax\else\hspace{6pt}{\popxbold\color{white}#1}\fi\par\vspace{7pt}}}
% Exercice résolu : énoncé en haut, \tcblower puis la solution sur fond crème
\newcounter{popexo}\newcounter{popexr}
\newtcolorbox{exoresolu}[1][]{popbase,colframe=poporange,colbacklower=popcream,
  segmentation style={popink,line width=1.2pt,dashed},
  before upper={\stepcounter{popexr}\popboxhead{Exercice résolu \thepopexr}{poporange}{#1}},
  before lower={\poppill{Solution}{popink}{popcream}\par\vspace{6pt}}}
% Exercice (non résolu en place) — la correction est dans la partie Corrigés
\newtcolorbox{exercice}[1][]{popbase,colframe=popink,
  before upper={\stepcounter{popexo}\popboxhead{Exercice \thepopexo}{popink}{#1}}}
% Corrigé
\newtcolorbox{corrige}[1][]{popbase,colframe=popgreen,colback=popgreenL,
  before upper={\popboxhead{Corrigé}{popgreen}{#1}}}
% Objectifs / prérequis / bilan
\newtcolorbox{objectifs}{popbase,colframe=popink,colback=poporangeL,
  before upper={\popboxhead{Objectifs de la leçon}{popink}{}}}
\newtcolorbox{prerequis}{popbase,colframe=popink,colback=popblueL,
  before upper={\popboxhead{Prérequis}{popblue}{}}}
\newtcolorbox{bilan}{popbase,colframe=popink,colback=white,boxrule=2.8pt,
  before upper={\popboxhead{Fiche bilan}{poporange}{L'essentiel à connaître pour l'examen}}}
% Figure encadrée avec légende
\newtcolorbox{popfigure}[1][]{enhanced,breakable=false,colback=white,colframe=popline,boxrule=1.4pt,arc=8pt,
  left=10pt,right=10pt,top=10pt,bottom=8pt,before skip=10pt,after skip=12pt,halign=center,
  lower separated=false,halign lower=center,
  fontlower=\popsemi\fontsize{9}{11.5}\selectfont\color{popmuted},
  #1}
\newcommand{\legende}[1]{\par\vspace{6pt}{\popsemi\fontsize{9}{11.5}\selectfont\color{popmuted}#1}}

% ============================================================
% Signature OnBuch+
% ============================================================
\newcommand{\poplogoinline}[1]{{\poplogo\fontsize{#1}{#1}\selectfont\color{popink}OnBuch\color{poporange}+}}
\newcommand{\signatureonbuch}{%
  \begin{tcolorbox}[enhanced,colback=popink,colframe=popink,boxrule=2.2pt,arc=10pt,
      drop shadow=poporange,left=14pt,right=14pt,top=10pt,bottom=10pt,before skip=0pt,after skip=0pt]
    \tikz[baseline=-0.7ex]\node[circle,fill=popgreen,draw=popcream,line width=1.3pt,minimum size=22pt,inner sep=0]{\popcheck{white}};%
    \hspace{9pt}\begin{minipage}[c]{0.62\linewidth}
      {\popxbold\fontsize{12}{13}\selectfont\color{popcream}Signé\ \textbullet\ {\poplogo OnBuch\color{poporange}+}}\par\vspace{2pt}
      {\raggedright\popsemi\fontsize{8}{10.5}\selectfont\color{popline}\MakeUppercase{Équipe pédagogique OnBuch+}\\\MakeUppercase{Conforme au programme officiel MINESEC}\par}
    \end{minipage}\hfill
    {\poplogo\fontsize{22}{22}\selectfont\color{popcream}OnBuch\color{poporange}+}
  \end{tcolorbox}%
}

% ============================================================
% Page de garde
% #1 titre · #2 matière · #3 niveau · #4 module · #5 n° leçon · #6 durée ·
% #7 description / objectifs (texte ou itemize)
% ============================================================
\newcommand{\pagegarde}[7]{%
  \thispagestyle{empty}
  \newgeometry{a4paper,margin=1.7cm,top=1.6cm,bottom=1.4cm}%
  \begin{tikzpicture}[remember picture,overlay]
    \fill[poporange!18] ([xshift=-3.2cm,yshift=-3.6cm]current page.north east) circle (4.6cm);
    \fill[poppurple!14] ([xshift=2.4cm,yshift=7.6cm]current page.south west) circle (3.8cm);
  \end{tikzpicture}%
  {\poplogo\fontsize{30}{30}\selectfont\color{popink}OnBuch\color{poporange}+}\hfill
  \raisebox{0.6ex}{\poppill{Cours signé \textbullet\ Premium}{popink}{popcream}}\par
  \vspace{1pt}\popsquiggle{poppurple}\par
  \vspace{8pt}
  \begin{tcolorbox}[enhanced,colback=poporange,colframe=popink,boxrule=2.8pt,arc=13pt,
      drop shadow=popink,left=18pt,right=18pt,top=12pt,bottom=13pt,after skip=10pt]
    \poppill{#2 \textbullet\ #3}{popink}{popcream}\hspace{5pt}\poppill{#4}{white}{popink}\par\vspace{8pt}
    {\popdisplay\fontsize{14}{14}\selectfont\color{popink}LEÇON #5}\par\vspace{4pt}
    {\raggedright\hyphenpenalty=10000\exhyphenpenalty=10000\popdisplay\fontsize{25}{27}\selectfont\color{popcream}\MakeUppercase{#1}\par}
    \vspace{8pt}
    {\popxbold\fontsize{9.5}{9.5}\selectfont\color{popink}\MakeUppercase{Durée conseillée : #6}}
  \end{tcolorbox}
  \begin{tcolorbox}[enhanced,colback=white,colframe=popink,boxrule=2.2pt,arc=10pt,
      drop shadow=popink,left=16pt,right=16pt,top=10pt,bottom=10pt,after skip=8pt]
    \poppill{Dans cette leçon}{poppurple}{white}\par\vspace{5pt}
    \popsemi\fontsize{9.3}{12.6}\selectfont\color{popink}#7
  \end{tcolorbox}
  \noindent\begin{minipage}[t]{0.487\linewidth}
    \begin{tcolorbox}[enhanced,colback=popgreen,colframe=popink,boxrule=2.2pt,arc=10pt,
        drop shadow=popink,left=11pt,right=11pt,top=10pt,bottom=10pt,height=3.1cm,valign=center]
      \tcbox[colback=white,colframe=popink,boxrule=1.3pt,arc=5pt,boxsep=0pt,nobeforeafter,
        left=3pt,right=3pt,top=3pt,bottom=3pt,valign=center]{\qrcode[height=1.9cm]{https://play.google.com/store/apps/details?id=com.luvvix.onbuch&hl=fr}}%
      \hspace{8pt}\begin{minipage}[c]{0.5\linewidth}
        {\popdisplay\fontsize{11}{12.5}\selectfont\color{white}\MakeUppercase{Télécharge OnBuch}}\par\vspace{4pt}
        {\raggedright\popsemi\fontsize{8.4}{11}\selectfont\color{white}Gratuit \textbullet\ Google Play\\Tuteur IA, annales, concours, résultats\par}
      \end{minipage}
    \end{tcolorbox}
  \end{minipage}\hfill
  \begin{minipage}[t]{0.487\linewidth}
    \begin{tcolorbox}[enhanced,colback=poppurple,colframe=popink,boxrule=2.2pt,arc=10pt,
        drop shadow=popink,left=11pt,right=11pt,top=10pt,bottom=10pt,height=3.1cm,valign=center]
      \tikz[baseline=-0.7ex]\node[circle,fill=white,draw=popink,line width=1.3pt,minimum size=1.6cm,inner sep=0]{\popdisplay\fontsize{24}{24}\selectfont\color{poppurple}+};%
      \hspace{8pt}\begin{minipage}[c]{0.6\linewidth}
        {\popdisplay\fontsize{11}{12.5}\selectfont\color{white}\MakeUppercase{Passe à OnBuch+}}\par\vspace{4pt}
        {\raggedright\popsemi\fontsize{8.4}{11}\selectfont\color{white}Tous les cours signés, Tuteur IA illimité, fiches PDF — onglet OnBuch+ de l'app\par}
      \end{minipage}
    \end{tcolorbox}
  \end{minipage}
  \vspace{8pt}
  \popcontenu
  \vfill
  \signatureonbuch
  \restoregeometry
  \newpage
}

% Rangée « Ce cours contient » (page de garde)
\newcommand{\poptile}[3]{% #1 couleur #2 gros texte #3 légende
  \begin{tcolorbox}[enhanced,colback=#1,colframe=popink,boxrule=2pt,arc=9pt,shadow={2.5pt}{-2.5pt}{0pt}{popink},
      left=6pt,right=6pt,top=9pt,bottom=9pt,halign=center,valign=center,height=1.75cm,width=\linewidth,nobeforeafter]
    {\popdisplay\fontsize{12.5}{13}\selectfont\color{white}\MakeUppercase{#2}}\par\vspace{3pt}
    {\centering\popsemi\fontsize{7.8}{9.5}\selectfont\color{white}#3\par}
  \end{tcolorbox}}
\newcommand{\popcontenu}{%
  \par\noindent{\popxboldls\fontsize{7.5}{7.5}\selectfont\color{popmuted}CE COURS SIGNÉ CONTIENT}\par\vspace{7pt}
  \noindent\begin{minipage}[t]{0.235\linewidth}\poptile{popblue}{Cours}{complet, illustré, pas à pas}\end{minipage}\hfill
  \begin{minipage}[t]{0.235\linewidth}\poptile{poporange}{Méthodes}{et exercices résolus}\end{minipage}\hfill
  \begin{minipage}[t]{0.235\linewidth}\poptile{poppink}{Exercices}{type examen + corrigés}\end{minipage}\hfill
  \begin{minipage}[t]{0.235\linewidth}\poptile{popgreen}{Fiche bilan}{l'essentiel à retenir}\end{minipage}\par}

% Sommaire
\newtcolorbox{sommaire}{enhanced,colback=white,colframe=popink,boxrule=2.2pt,arc=10pt,
  drop shadow=popink,left=18pt,right=18pt,top=13pt,bottom=13pt,before skip=6pt,after skip=20pt,
  before upper={\poppill{Sommaire}{poppurple}{white}\par\vspace{9pt}\popsemi\fontsize{10.6}{14.5}\selectfont\color{popink}}}

% Signature de fin de cours — poussée tout en bas de la dernière page
\newcommand{\findecours}{%
  \par\vfill\nopagebreak
  \begin{center}\popsquiggle{poporange}\end{center}\vspace{4pt}
  \signatureonbuch
}
\AtBeginDocument{\def\llbracket{\mathopen{[\mkern-3.5mu[}}\def\rrbracket{\mathclose{]\mkern-3.5mu]}}} % intervalles d'entiers (glyphe ⟦ absent de la police)
% Glyphes absents de Fira Math : redéfinis à partir de glyphes disponibles
\AtBeginDocument{%
  \def\setminus{\mathbin{\backslash}}\let\smallsetminus\setminus
  \def\vdots{\mathord{\vbox{\baselineskip3.2pt\lineskiplimit0pt\kern4pt\hbox{.}\hbox{.}\hbox{.}}}}%
  \def\ddots{\mathinner{\mkern1mu\raise6.5pt\hbox{.}\mkern2mu\raise3.6pt\hbox{.}\mkern2mu\raise0.7pt\hbox{.}\mkern1mu}}%
  \def\triangle{\mathord{\scalebox{0.9}{$\symup{\Delta}$}}}%
  \def\nabla{\mathord{\raisebox{\depth}{\scalebox{1}[-1]{$\symup{\Delta}$}}}}%
  \def\checkmark{\popcheck{popdark}}%
  \def\star{\mathbin{\ast}}\def\bigstar{\mathord{\ast}}%
  \def\diamond{\mathbin{\scalebox{0.6}{$\Diamond$}}}%
  \def\bigcup{\mathop{\mathchoice{\raisebox{-0.3em}{\scalebox{1.7}{$\cup$}}}{\scalebox{1.25}{$\cup$}}{\cup}{\cup}}\displaylimits}%
  \def\bigcap{\mathop{\mathchoice{\raisebox{-0.3em}{\scalebox{1.7}{$\cap$}}}{\scalebox{1.25}{$\cap$}}{\cap}{\cap}}\displaylimits}%
  \def\lesseqqgtr{\mathrel{\lessgtr}}\def\bigoplus{\mathop{\oplus}\displaylimits}%
}
\providecommand{\ssi}{si et seulement si} % abréviation fréquente
\AtBeginDocument{\providecommand{\wideparen}{\overparen}} % arc de cercle

\begin{document}

\pagegarde{La identidad, la tolerancia y la convivencia}{Espagnol}{1ère A}{Módulo 3 : Citoyenneté}{EA11}{2 h}{Cette leçon de type texte propose la lecture et le commentaire d'un texte sur l'identité, la tolérance et le vivre-ensemble, enrichis d'un lexique thématique complet et d'outils pour exprimer l'accord et le désaccord. Elle aboutit à une production guidée sur la convivencia dans une société multiculturelle comme le Cameroun.
\vspace{4pt}
\begin{multicols}{2}\small
\begin{itemize}
\item À la fin de cette leçon, je sais lire et comprendre un texte espagnol sur l'identité et la tolérance
\item À la fin de cette leçon, je sais définir et employer le lexique du vivre-ensemble (identidad, tolerancia, respeto, convivencia, prejuicio, diferencia, integración, discriminación)
\item À la fin de cette leçon, je sais exprimer mon accord (estoy de acuerdo, tienes razón, es verdad que...)
\item À la fin de cette leçon, je sais exprimer mon désaccord (no estoy de acuerdo, no creo que + subjonctif, al contrario...)
\item À la fin de cette leçon, je sais rédiger un commentaire structuré d'un texte argumentatif
\item À la fin de cette leçon, je sais rédiger quelques phrases personnelles sur le vivre-ensemble au Cameroun
\end{itemize}\end{multicols}}

\begin{sommaire}
\begin{multicols}{2}
\begin{enumerate}
\item Texte support : « Vivir juntos en la diferencia »
\item Compréhension du texte
\item Lexique thématique : identidad, tolerancia y convivencia
\item Exprimer l'accord et le désaccord
\item Méthode du commentaire de texte et commentaire modèle
\item Production guidée : le vivre-ensemble au Cameroun
\item Activité d'intégration
\item Méthodes et exercices résolus
\item Exercices
\item Corrigés des exercices
\item Fiche bilan
\end{enumerate}
\end{multicols}
\end{sommaire}

\begin{objectifs}
\begin{itemize}
\item À la fin de cette leçon, je sais lire et comprendre un texte espagnol sur l'identité et la tolérance
\item À la fin de cette leçon, je sais définir et employer le lexique du vivre-ensemble (identidad, tolerancia, respeto, convivencia, prejuicio, diferencia, integración, discriminación)
\item À la fin de cette leçon, je sais exprimer mon accord (estoy de acuerdo, tienes razón, es verdad que...)
\item À la fin de cette leçon, je sais exprimer mon désaccord (no estoy de acuerdo, no creo que + subjonctif, al contrario...)
\item À la fin de cette leçon, je sais rédiger un commentaire structuré d'un texte argumentatif
\item À la fin de cette leçon, je sais rédiger quelques phrases personnelles sur le vivre-ensemble au Cameroun
\item À la fin de cette leçon, je sais repérer les valeurs et les arguments d'un texte de débat de société
\end{itemize}
\end{objectifs}

\begin{prerequis}
\begin{itemize}
\item Présent de l'indicatif des verbes réguliers et irréguliers courants (ser, estar, tener, respetar)
\item Les adjectifs qualificatifs et leur accord en genre et en nombre
\item Les structures d'opinion simples : pienso que, creo que, me parece que
\item Le vocabulaire de la société et de la famille vu en Seconde
\item La construction d'une phrase simple et d'un paragraphe en espagnol
\end{itemize}
\end{prerequis}

\newpage

\coursec{Texte support : « Vivir juntos en la diferencia »}

\courssub{Situation d'accroche et problématique}

Imaginez une cour de récréation au lycée de Yaoundé : ici, un groupe parle ewondo ; là, des élèves rient en langue bamiléké ; plus loin, on discute en fulfuldé, en duala ou en anglais. Au marché Mokolo, à Douala, le vendeur salue en pidgin, la cliente répond en français, et tout le monde se comprend. Au Cameroun, environ 250 groupes ethniques cohabitent et parlent des centaines de langues : à l'école, au marché ou sur les réseaux sociaux, nous sommes chaque jour confrontés à la différence.

Mais voici la question : comment dire en espagnol que nous acceptons l'autre tel qu'il est ? Comment nommer le respect, la tolérance, le préjugé ? Dans le monde hispanophone aussi, de l'Espagne multiculturelle à l'Amérique latine métissée, en passant par la Guinée équatoriale voisine, la question de la \esp{convivencia} est au cœur des débats de société. Cette leçon nous donne les mots et les outils pour en parler.

\textbf{Problématique :} comment un texte argumentatif espagnol défend-il l'idée que la diversité est une richesse, et quels outils linguistiques nous offre-t-il pour parler du vivre-ensemble ?

\courssub{Lecture du texte}

\textbf{Première étape : lecture silencieuse.} Lisez le texte une première fois sans dictionnaire, en essayant de saisir l'idée générale. Ne vous arrêtez pas sur les mots inconnus : le contexte aide souvent à les deviner. \textbf{Deuxième étape : lecture à voix haute} par le professeur, puis par quelques élèves, en respectant la ponctuation et l'intonation des phrases exclamatives et interrogatives.

\begin{textebox}[Vivir juntos en la diferencia]
Camerún es un país donde conviven más de doscientos cincuenta grupos étnicos. En las escuelas, en los mercados y en las calles de Yaoundé o de Douala, se hablan cientos de lenguas diferentes. Esta diversidad es una riqueza, pero también un desafío.

La tolerancia no significa renunciar a la propia identidad, sino respetar la de los demás. Cada grupo tiene su historia, sus tradiciones y su lengua, y cada uno tiene derecho a conservarlas. La identidad es lo que somos: nuestra cultura, nuestras costumbres, nuestra memoria. Nadie debe avergonzarse de sus raíces.

Sin embargo, los prejuicios nacen de la ignorancia y conducen a la discriminación. Cuando juzgamos a una persona sin conocerla, solo porque pertenece a otro grupo, cometemos una injusticia. Rechazar al diferente es empobrecernos a nosotros mismos.

Además, cuando aceptamos la diferencia, la convivencia se convierte en una fuente de enriquecimiento para todos. El vecino que habla otra lengua nos enseña otra manera de ver el mundo. El compañero de otra región comparte con nosotros sus platos, sus músicas y sus fiestas.

Por eso, la escuela debe enseñar el respeto y la integración desde la infancia. Un niño que aprende a respetar al otro será un adulto que construye la paz.

En conclusión, vivir juntos en la diferencia no es un sueño imposible: es un trabajo de cada día, que empieza con una palabra sencilla: respetar.
\end{textebox}

\textbf{Glossaire des mots difficiles :}

\begin{center}
\begin{tabularx}{\linewidth}{|X|X|}
\hline
\rowcolor{popblueL} \textbf{Español} & \textbf{Français} \\
\hline
la convivencia & la vie commune, le vivre-ensemble \\
\hline
convivir & cohabiter, vivre ensemble \\
\hline
el desafío & le défi \\
\hline
renunciar a & renoncer à \\
\hline
los demás & les autres \\
\hline
avergonzarse de & avoir honte de \\
\hline
las raíces & les racines \\
\hline
el prejuicio & le préjugé \\
\hline
conducir a & conduire à, mener à \\
\hline
rechazar & rejeter \\
\hline
empobrecer(se) & (s')appauvrir \\
\hline
enriquecer / el enriquecimiento & enrichir / l'enrichissement \\
\hline
el plato & le plat \\
\hline
la fuente & la source \\
\hline
el sueño & le rêve \\
\hline
\end{tabularx}
\end{center}

\begin{definition}[La convivencia]
\esp{La convivencia} : \esp{el hecho de vivir juntos en armonía respetando las diferencias} — le fait de vivre ensemble en harmonie en respectant les différences. C'est un nom féminin : \esp{la convivencia}, \esp{una buena convivencia}. Le verbe correspondant est \esp{convivir} : \esp{Convivimos con personas de otras culturas.} « Nous vivons avec des personnes d'autres cultures. »
\end{definition}

\begin{attention}[Prononciation : trois pièges du texte]
\begin{itemize}
\item \esp{convivencia} se prononce « ko-n-bi-bé-n-sia » : le \esp{v} espagnol se prononce comme un \emph{b} doux, et le \esp{c} devant \esp{e} se prononce « s » en Amérique latine (et « th » anglais dans une grande partie de l'Espagne).
\item \esp{prejuicio} se prononce « pré-houi-sio » : le \esp{j} est un « r » grasseyé espagnol, jamais le « j » français.
\item \esp{Yaoundé} se prononce « ya-oun-dé » en espagnol : le \esp{y} initial vaut un « i » bref suivi de « a ».
\end{itemize}
\end{attention}

\courssub{Compréhension globale : thème, genre, thèse}

Après la lecture, répondons aux trois questions fondamentales de la compréhension d'un texte argumentatif.

\textbf{1. Le thème.} De quoi parle le texte ? Il parle de la diversité ethnique et linguistique du Cameroun et de la manière de vivre ensemble malgré les différences. Le thème est donc \esp{la convivencia en una sociedad multicultural}.

\textbf{2. Le genre.} C'est un \cle{texte argumentatif} : l'auteur ne raconte pas une histoire, il défend une idée. On reconnaît ce genre à la présence d'une thèse, d'arguments et de connecteurs logiques.

\textbf{3. La thèse.} Quelle idée l'auteur défend-il ? La thèse est claire : \esp{la diversidad es una riqueza y la convivencia es posible gracias al respeto y a la tolerancia} — la diversité est une richesse et le vivre-ensemble est possible grâce au respect et à la tolérance. Elle apparaît dès le premier paragraphe (\esp{Esta diversidad es una riqueza}) et est résumée dans la conclusion.

\begin{methode}[Lire un texte argumentatif en quatre étapes]
\begin{enumerate}
\item \textbf{Identifier le thème} : de quoi parle le texte ? Cherchez les mots qui reviennent le plus souvent.
\item \textbf{Repérer la thèse} : quelle idée l'auteur défend-il ? Elle se trouve souvent dans l'introduction ou la conclusion.
\item \textbf{Souligner les arguments} : quelles raisons l'auteur donne-t-il pour convaincre ? Ici : les préjugés mènent à la discrimination ; la différence enrichit ; l'école forme les citoyens.
\item \textbf{Relever les connecteurs} : ils montrent la structure du raisonnement (opposition, addition, cause, conclusion).
\end{enumerate}
\end{methode}

\courssub{Le champ lexical de la tolérance}

Un \cle{champ lexical} est l'ensemble des mots qui se rapportent à une même idée. Relevons dans le texte les mots du champ lexical de la tolérance et du vivre-ensemble :

\begin{itemize}
\item \esp{la tolerancia} (la tolérance), \esp{el respeto} (le respect), \esp{respetar} (respecter) ;
\item \esp{la convivencia} (le vivre-ensemble), \esp{convivir} (cohabiter, vivre ensemble) ;
\item \esp{la integración} (l'intégration), \esp{la paz} (la paix) ;
\item \esp{la diversidad} (la diversité), \esp{la diferencia} (la différence), \esp{la riqueza} (la richesse) ;
\item \esp{la identidad} (l'identité), \esp{las raíces} (les racines), \esp{las costumbres} (les coutumes).
\end{itemize}

À l'opposé, le texte contient aussi le champ lexical du \emph{rejet} : \esp{el prejuicio} (le préjugé), \esp{la ignorancia} (l'ignorance), \esp{la discriminación} (la discrimination), \esp{rechazar} (rejeter), \esp{la injusticia} (l'injustice). Un bon texte argumentatif oppose souvent deux champs lexicaux : celui de ce qu'il défend et celui de ce qu'il combat.

\begin{popfigure}
\begin{center}
\begin{tikzpicture}[mindmap, grow cyclic, every node/.style={concept}, concept color=popblue!60, text=white,
  level 1/.append style={level distance=3.2cm, sibling angle=90},
  level 2/.append style={level distance=2.2cm}]
\node[concept, font=\large\bfseries] {convivencia}
  child[concept color=popgreen!70] { node[concept] {identidad} }
  child[concept color=poporange!80] { node[concept] {tolerancia} }
  child[concept color=poppurple!70] { node[concept] {respeto} }
  child[concept color=poppink!80] { node[concept] {integración} };
\end{tikzpicture}
\end{center}
\legende{Carte mentale : la convivencia et ses quatre piliers}
\end{popfigure}

\courssub{Les connecteurs logiques du texte}

Les \cle{connecteurs logiques} sont les articulations du raisonnement. Le texte en utilise quatre, qu'il faut absolument connaître :

\begin{center}
\begin{tabularx}{\linewidth}{|X|X|X|}
\hline
\rowcolor{popblueL} \textbf{Conector} & \textbf{Français} & \textbf{Fonction} \\
\hline
\esp{sin embargo} & cependant, pourtant & opposition \\
\hline
\esp{además} & de plus, en outre & addition d'un argument \\
\hline
\esp{por eso} & c'est pourquoi, donc & conséquence \\
\hline
\esp{en conclusión} & en conclusion & bilan final \\
\hline
\end{tabularx}
\end{center}

Observons-les dans leur contexte :

\esp{Los prejuicios nacen de la ignorancia. Sin embargo, cuando aceptamos la diferencia, la convivencia se convierte en una fuente de enriquecimiento.} « Les préjugés naissent de l'ignorance. Cependant, quand nous acceptons la différence, le vivre-ensemble devient une source d'enrichissement. » — \esp{sin embargo} introduit un retournement : après le danger (les préjugés), l'auteur présente la solution (l'acceptation).

\esp{Por eso, la escuela debe enseñar el respeto desde la infancia.} « C'est pourquoi l'école doit enseigner le respect dès l'enfance. » — \esp{por eso} tire une conséquence de tout ce qui précède.

\begin{attention}[Faux amis et pièges de genre]
\begin{itemize}
\item Ne confondez pas \esp{identidad} (l'identité, ce que l'on est) et \esp{identificación} (l'identification, l'action d'identifier ou la pièce d'identité). \esp{Mi identidad es camerunesa.} « Mon identité est camerounaise. » ; \esp{Muestra tu identificación.} « Montre ta pièce d'identité. »
\item \esp{Tolerancia} est féminin : on dit \esp{la tolerancia}, \esp{una gran tolerancia}. De même : \esp{la identidad}, \esp{la diversidad}, \esp{la integración} — les noms en \esp{-dad} et \esp{-ción} sont presque toujours féminins.
\item \esp{Respeto} est masculin : \esp{el respeto}. Ne dites pas « la respeto ».
\item Attention au verbe \esp{enriquecer} : il se conjugue comme \esp{agradecer} : \esp{yo enriquezco}, avec un \esp{-zc-} à la première personne du présent. Même phénomène pour \esp{conducir} : \esp{yo conduzco}.
\item \esp{Sin embargo} seul signifie « cependant » ; ne le confondez pas avec \esp{embargo} qui, dans \esp{el embargo}, signifie « l'embargo ».
\end{itemize}
\end{attention}

\begin{exoresolu}[Repérage dans le texte]
\textbf{Énoncé.} 1) Relevez dans le texte deux connecteurs et précisez leur fonction. 2) Trouvez dans le deuxième paragraphe la phrase qui définit la tolérance et traduisez-la. 3) Quel mot du texte signifie « préjugé » et quel connecteur l'oppose à l'acceptation de la différence ?
\tcblower
\textbf{Solution.} 1) \esp{Sin embargo} (opposition : il oppose les préjugés à l'acceptation) ; \esp{en conclusión} (bilan : il introduit le paragraphe final). On pouvait aussi citer \esp{además} (addition) et \esp{por eso} (conséquence). 2) \esp{La tolerancia no significa renunciar a la propia identidad, sino respetar la de los demás.} « La tolérance ne signifie pas renoncer à sa propre identité, mais respecter celle des autres. » Remarquez la tournure \esp{no... sino...} qui correspond au français « ne... pas... mais... » ; \esp{la de los demás} reprend \esp{la identidad} pour éviter la répétition, comme le français « celle des autres ». 3) Le mot est \esp{prejuicio} ; c'est \esp{sin embargo} qui marque le passage des préjugés à l'acceptation de la différence.
\end{exoresolu}

\courssub{Mise en relation avec le contexte camerounais}

Le texte n'est pas une théorie abstraite : il décrit notre réalité quotidienne. Le Cameroun est souvent appelé « l'Afrique en miniature » parce qu'il réunit une extraordinaire diversité de peuples, de langues et de paysages. Les exemples du texte — les marchés de Yaoundé et de Douala, les camarades d'autres régions, les plats et les musiques partagés — parlent à chaque élève.

Cette réalité rejoint celle du monde hispanophone : l'Espagne accueille des communautés venues d'Amérique latine, d'Afrique du Nord et d'Europe de l'Est ; le Mexique, le Pérou ou la Bolivie sont des nations métisses où coexistent langues indigènes et espagnol ; la Guinée équatoriale, seul pays hispanophone d'Afrique subsaharienne, partage avec le Cameroun la même question du vivre-ensemble entre groupes fang, bubi, ndowe et autres. Le mot \esp{convivencia} est donc un mot-clé pour comprendre à la fois notre pays et le monde hispanique.

\begin{savaistu}[Un mot presque intraduisible]
Le mot \esp{convivencia} est presque intraduisible en un seul mot français. On le traduit par « vivre-ensemble », « vie commune » ou « cohabitation », mais aucun de ces mots ne rend tout à fait sa richesse : la \esp{convivencia} désigne l'art concret, quotidien, de vivre côte à côte — partager l'école, le quartier, le marché — en acceptant les différences de l'autre. En Espagne, le mot a aussi une longue histoire : il évoque l'époque médiévale où chrétiens, juifs et musulmans vivaient ensemble dans les mêmes villes. C'est dire si ce mot porte en lui toute une leçon de tolérance !
\end{savaistu}

\begin{exercice}[À vous de jouer]
\begin{enumerate}
\item Traduisez en français : \esp{La escuela debe enseñar el respeto y la integración desde la infancia.}
\item Complétez avec le connecteur qui convient (\esp{sin embargo}, \esp{además}, \esp{por eso}) : « La diversidad es una riqueza; \dots, muchas personas tienen prejuicios. \dots, debemos combatir la ignorancia. »
\item Relevez dans le texte trois mots du champ lexical de la tolérance et trois mots du champ lexical du rejet.
\item En une phrase, formulez la thèse de l'auteur avec vos propres mots, en espagnol.
\end{enumerate}
\end{exercice}

\begin{corrige}[Exercice : À vous de jouer]
\begin{enumerate}
\item « L'école doit enseigner le respect et l'intégration dès l'enfance. »
\item « La diversidad es una riqueza; \esp{sin embargo}, muchas personas tienen prejuicios. \esp{Por eso}, debemos combatir la ignorancia. »
\item Tolérance : \esp{tolerancia}, \esp{respeto}, \esp{integración} (ou \esp{convivencia}, \esp{paz}). Rejet : \esp{prejuicio}, \esp{discriminación}, \esp{ignorancia} (ou \esp{injusticia}, \esp{rechazar}).
\item Exemple : \esp{La diversidad es una riqueza y podemos vivir juntos si respetamos las diferencias.} (« La diversité est une richesse et nous pouvons vivre ensemble si nous respectons les différences. »)
\end{enumerate}
\end{corrige}

\begin{aretenir}[L'essentiel de cette section]
\begin{itemize}
\item \esp{Vivir juntos en la diferencia} est un texte argumentatif : il défend la thèse que la diversité camerounaise est une richesse et que la \esp{convivencia} repose sur le respect.
\item Le champ lexical de la tolérance (\esp{tolerancia}, \esp{respeto}, \esp{integración}, \esp{convivencia}) s'oppose au champ lexical du rejet (\esp{prejuicio}, \esp{discriminación}, \esp{ignorancia}).
\item Quatre connecteurs structurent le raisonnement : \esp{sin embargo} (opposition), \esp{además} (addition), \esp{por eso} (conséquence), \esp{en conclusión} (bilan).
\item Pour lire un texte argumentatif : identifier le thème, repérer la thèse, souligner les arguments, relever les connecteurs.
\end{itemize}
\end{aretenir}

\coursec{Compréhension du texte}

Dans la section précédente, nous avons lu et découvert le texte support \esp{« Vivir juntos en la diferencia »}, un texte argumentatif sur l'identité, la tolérance et le vivre-ensemble au Cameroun. Lire un texte ne suffit pas : il faut maintenant le \cle{comprendre} en profondeur. C'est exactement ce que l'on te demandera à l'examen : répondre à des \cle{preguntas de comprensión} qui vérifient que tu as saisi le sens global, les détails, les intentions de l'auteur, et que tu sais réagir personnellement. Cette section t'apprend à répondre à chaque type de question, méthodiquement et en espagnol correct.

\courssub{Les quatre types de questions de compréhension}

À l'examen, les questions de compréhension ne sont pas toutes du même niveau. On distingue classiquement quatre types, du plus facile au plus exigeant :

\begin{enumerate}
\item \textbf{Questions de compréhension globale} : elles portent sur le sens général du texte. Exemples : \esp{¿De qué habla el texto?} \esp{¿Cuál es la tesis del autor?}
\item \textbf{Questions de détail} : elles portent sur une information précise, localisable dans le texte. Exemples : \esp{¿Cuántos grupos étnicos conviven en Camerún?} \esp{¿De dónde nacen los prejuicios?}
\item \textbf{Questions d'interprétation} : elles demandent d'expliquer le sens d'une phrase ou d'une expression de l'auteur. Exemple : \esp{¿Qué significa « la tolerancia no significa renunciar a la propia identidad »?}
\item \textbf{Questions de réaction personnelle} : elles demandent ton avis, toujours justifié. Exemples : \esp{¿Estás de acuerdo con el autor?} \esp{¿Por qué?}
\end{enumerate}

\begin{methode}[Répondre à une question de compréhension]
\begin{enumerate}
\item \textbf{Lis la question deux fois} et souligne le mot interrogatif : \esp{¿qué?} (quoi), \esp{¿cuántos?} (combien), \esp{¿de dónde?} (d'où), \esp{¿por qué?} (pourquoi), \esp{¿estás de acuerdo?} (es-tu d'accord).
\item \textbf{Repère dans le texte} le passage qui contient la réponse.
\item \textbf{Rédige une phrase complète en espagnol} : reprends les mots de la question pour construire ta phrase. Jamais de réponse d'un seul mot.
\item \textbf{Cite le texte entre guillemets} pour justifier ta réponse : \esp{El autor dice: « ... »} ou \esp{Según el texto, « ... »}.
\item \textbf{Relis-toi} : verbe conjugué à la bonne personne, accents, ponctuation espagnole (¿ ... ?).
\end{enumerate}
\end{methode}

\begin{propriete}[Questions ouvertes et questions fermées]
Une \cle{question fermée} appelle une information précise du texte (\esp{¿Cuántos grupos étnicos hay?}) : la réponse se trouve mot pour mot dans le texte. Une \cle{question ouverte} (\esp{¿Estás de acuerdo?}, \esp{¿Qué significa...?}) exige une \textbf{réponse personnelle appuyée sur le texte} : à l'examen, une simple recopie du texte ne rapporte pas tous les points. Il faut reformuler, expliquer, donner ton avis, tout en montrant que tu t'appuies sur ce que dit l'auteur.
\end{propriete}

\courssub{Questions de compréhension globale}

La première question est presque toujours : \esp{¿De qué habla el texto?} (« De quoi parle le texte ? »). On y répond en une ou deux phrases qui donnent le \cle{thème} général.

\begin{exemplebox}[Réponse modèle — compréhension globale]
\esp{El texto habla de la convivencia entre los diferentes grupos étnicos y culturales de Camerún. El autor defiende la idea de que la diversidad es una riqueza y no un problema.}\\
« Le texte parle de la vie commune entre les différents groupes ethniques et culturels du Cameroun. L'auteur défend l'idée que la diversité est une richesse et non un problème. »
\end{exemplebox}

La deuxième question globale porte sur la \cle{tesis} (la thèse), c'est-à-dire l'idée principale que l'auteur veut démontrer. Pour la trouver, demande-toi : « Que veut me convaincre l'auteur ? »

\begin{exemplebox}[Réponse modèle — la thèse]
\esp{La tesis del autor es que la tolerancia y el respeto permiten transformar la diversidad cultural en una fuente de enriquecimiento para todos.}\\
« La thèse de l'auteur est que la tolérance et le respect permettent de transformer la diversité culturelle en une source d'enrichissement pour tous. »
\end{exemplebox}

\begin{attention}[Thème ou thèse ?]
Ne confonds pas le \cle{tema} (le thème : de quoi on parle) et la \cle{tesis} (la thèse : ce que l'auteur affirme). Le thème tient en quelques mots (\esp{la convivencia en Camerún}) ; la thèse est une phrase complète avec un verbe d'opinion (\esp{la diversidad es una riqueza}). À l'examen, répondre « le texte parle de la tolérance » à la question \esp{¿Cuál es la tesis?} est insuffisant.
\end{attention}

\courssub{Questions de détail : localiser l'information}

Les questions de détail vérifient ta lecture précise. La méthode : repérer le \textbf{mot-clé} de la question, le chercher dans le texte, puis rédiger une phrase complète.

\textbf{Question 1 :} \esp{¿Cuántos grupos étnicos conviven en Camerún?}

Réponse modèle : \esp{En Camerún conviven más de doscientos cincuenta grupos étnicos, según el texto.} « Plus de deux cent cinquante groupes ethniques vivent ensemble au Cameroun, selon le texte. »

\textbf{Question 2 :} \esp{¿De dónde nacen los prejuicios?}

Réponse modèle : \esp{Según el autor, los prejuicios nacen de la ignorancia y del miedo a lo diferente.} « Selon l'auteur, les préjugés naissent de l'ignorance et de la peur de ce qui est différent. »

\begin{exemplebox}[Exemples traduits à retenir]
\esp{Los prejuicios nacen de la ignorancia.} = « Les préjugés naissent de l'ignorance. »\\
\esp{La convivencia se convierte en una fuente de enriquecimiento.} = « La vie commune devient une source d'enrichissement. »\\
\esp{La diversidad no es un problema, sino una riqueza.} = « La diversité n'est pas un problème, mais une richesse. »
\end{exemplebox}

\begin{attention}[Faux amis et calques]
\begin{itemize}
\item \esp{nacer de} = « naître de » ; ne traduis pas \esp{¿De dónde nacen?} par « d'où naissent-ils ? » mot à mot dans ta tête : comprends « quelle est l'origine de... ? ».
\item \esp{convivir} = « vivre ensemble, cohabiter » ; ne dis pas \esp{convivenciar} : le nom est \esp{la convivencia}.
\item \esp{la fuente} = « la source » (et non « la fontaine » dans ce contexte abstrait).
\item \esp{no... sino...} = « ne... pas... mais... » : après une négation, l'espagnol utilise \esp{sino} et non \esp{pero} : \esp{La diversidad no es un problema, sino una riqueza.}
\end{itemize}
\end{attention}

\courssub{Questions d'interprétation : expliquer avec ses propres mots}

La question d'interprétation te donne une citation et te demande : \esp{¿Qué significa...?} (« Que signifie... ? »). La bonne démarche : \textbf{reformuler} la phrase avec des mots plus simples, puis \textbf{expliquer} ce que l'auteur veut dire.

\textbf{Question :} \esp{¿Qué significa « la tolerancia no significa renunciar a la propia identidad »?}

\begin{exemplebox}[Réponse modèle — interprétation]
\esp{Esta frase significa que aceptar y respetar a los demás no obliga a abandonar la propia cultura, la propia lengua o las propias tradiciones. El autor quiere decir que se puede ser tolerante y, al mismo tiempo, seguir siendo fiel a la propia identidad.}\\
« Cette phrase signifie qu'accepter et respecter les autres n'oblige pas à abandonner sa propre culture, sa propre langue ou ses propres traditions. L'auteur veut dire qu'on peut être tolérant et, en même temps, rester fidèle à sa propre identité. »
\end{exemplebox}

Remarque les outils linguistiques de la reformulation : \esp{significa que...} (« cela signifie que... »), \esp{el autor quiere decir que...} (« l'auteur veut dire que... »), \esp{es decir} (« c'est-à-dire »), \esp{en otras palabras} (« en d'autres termes »).

\courssub{Questions de réaction personnelle : donner et justifier son avis}

C'est le type de question le plus délicat, car il n'y a pas de réponse unique : on évalue ta capacité à t'exprimer et à argumenter.

\textbf{Question :} \esp{¿Estás de acuerdo con el autor? ¿Por qué?}

\begin{exemplebox}[Réponses modèles — réaction personnelle]
\esp{Sí, estoy de acuerdo con el autor porque en mi ciudad, Yaoundé, conviven personas de muchas regiones y esta diversidad enriquece nuestra vida cotidiana.}\\
« Oui, je suis d'accord avec l'auteur parce que dans ma ville, Yaoundé, vivent ensemble des personnes de nombreuses régions et cette diversité enrichit notre vie quotidienne. »

\esp{Estoy de acuerdo en parte: la tolerancia es necesaria, pero también hace falta educar a los jóvenes en el respeto desde la escuela.}\\
« Je suis d'accord en partie : la tolérance est nécessaire, mais il faut aussi éduquer les jeunes au respect dès l'école. »
\end{exemplebox}

\begin{attention}[Ne jamais répondre seulement « sí » ou « no »]
À la question \esp{¿Estás de acuerdo?}, une réponse d'un seul mot est \textbf{sanctionnée} : elle ne prouve ni compréhension ni expression. Justifie toujours avec \esp{porque...} (« parce que... »). Structure minimale attendue : \textbf{avis} (\esp{estoy de acuerdo / no estoy de acuerdo / estoy de acuerdo en parte}) + \textbf{justification} (\esp{porque...}) + si possible un \textbf{exemple} (\esp{por ejemplo...}).
\end{attention}

\begin{center}
\begin{tabularx}{\linewidth}{|X|X|X|}
\hline
\rowcolor{popblueL} Français & Español & Emploi \\
\hline
Je suis d'accord (avec) & Estoy de acuerdo (con) & accord total \\
\hline
Je ne suis pas d'accord & No estoy de acuerdo & désaccord \\
\hline
Je suis d'accord en partie & Estoy de acuerdo en parte & accord nuancé \\
\hline
parce que & porque & justification \\
\hline
par exemple & por ejemplo & illustration \\
\hline
selon l'auteur / le texte & según el autor / el texto & appui sur le texte \\
\hline
c'est-à-dire & es decir & reformulation \\
\hline
à mon avis & en mi opinión / a mi parecer & opinion personnelle \\
\hline
\end{tabularx}
\end{center}

\courssub{Vrai ou faux justifié : l'exercice de précision}

L'exercice \esp{¿Verdadero o falso? Justifica tu respuesta con una cita del texto} (« Vrai ou faux ? Justifie ta réponse par une citation du texte ») est un classique de l'examen. La justification par une \textbf{citation entre guillemets} est \textbf{obligatoire} : sans elle, la réponse n'est pas validée, même si « verdadero » est correct.

\begin{exoresolu}[Vrai ou faux justifié]
Énoncé. Réponds par \esp{verdadero} ou \esp{falso} et justifie par une citation du texte \esp{« Vivir juntos en la diferencia »}.

1. \esp{En Camerún hay un solo grupo étnico.}

2. \esp{Los prejuicios nacen del conocimiento profundo de los demás.}

3. \esp{Ser tolerante significa abandonar la propia identidad.}
\tcblower
Solution détaillée.

1. \textbf{Falso.} Le texte affirme le contraire : \esp{El texto dice: « más de doscientos cincuenta grupos étnicos conviven en Camerún »}. Le mot-clé \esp{un solo} (« un seul ») contredit le chiffre donné par l'auteur.

2. \textbf{Falso.} Le texte dit exactement l'inverse : \esp{« los prejuicios nacen de la ignorancia »}. Les préjugés naissent de l'\emph{ignorance}, pas de la connaissance. Attention : c'est un piège classique qui inverse le sens d'une phrase du texte.

3. \textbf{Falso.} Le texte précise : \esp{« la tolerancia no significa renunciar a la propia identidad »}. La négation \esp{no} est essentielle : la phrase de l'énoncé la supprime et inverse donc le sens. Souligne toujours les négations dans le texte.
\end{exoresolu}

\begin{methode}[Réussir le vrai ou faux justifié]
\begin{enumerate}
\item Souligne dans l'énoncé les mots « dangereux » : \esp{solo} (seul), \esp{nunca} (jamais), \esp{todos} (tous), \esp{no} (ne... pas). Ce sont eux qui rendent une phrase fausse.
\item Retrouve la phrase correspondante dans le texte.
\item Compare mot à mot : l'énoncé a-t-il ajouté, supprimé ou inversé quelque chose ?
\item Rédige : \esp{Falso. El texto dice: « ... »} — la citation doit être exacte, entre guillemets.
\end{enumerate}
\end{methode}

\courssub{Reformuler la thèse avec ses propres mots}

Dernière compétence : être capable de \textbf{reformuler la thèse} de l'auteur sans recopier ses phrases. Cela prouve une compréhension réelle. Utilise les verbes et connecteurs de reformulation : \esp{el autor afirma que...} (« l'auteur affirme que... »), \esp{según el autor...}, \esp{en otras palabras...}, \esp{la idea principal es que...}.

\begin{exemplebox}[Deux reformulations possibles de la thèse]
Thèse du texte : \esp{La diversidad es una riqueza y la tolerancia permite vivir juntos en paz.}

Reformulation 1 : \esp{El autor afirma que la variedad de culturas es un tesoro para el país y que el respeto mutuo hace posible la convivencia pacífica.}\\
« L'auteur affirme que la variété des cultures est un trésor pour le pays et que le respect mutuel rend possible la coexistence pacifique. »

Reformulation 2 : \esp{En otras palabras, las diferencias entre los grupos no deben dividirnos, sino enriquecernos, si aprendemos a respetarnos.}\\
« En d'autres termes, les différences entre les groupes ne doivent pas nous diviser, mais nous enrichir, si nous apprenons à nous respecter. »
\end{exemplebox}

Pour t'aider à visualiser l'ensemble, voici l'organigramme de la structure argumentative du texte : il montre comment la thèse est soutenue par deux arguments et aboutit à une conclusion. Le comprendre, c'est pouvoir répondre à n'importe quelle question globale.

\begin{popfigure}
\begin{center}
\begin{tikzpicture}[
node distance=0.9cm,
tbox/.style={draw=popblue, fill=popblueL, rounded corners, align=center, text width=6.5cm, font=\small, inner sep=5pt},
abox/.style={draw=poporange, fill=poporangeL, rounded corners, align=center, text width=6.5cm, font=\small, inner sep=5pt},
cbox/.style={draw=popgreen, fill=popgreenL, rounded corners, align=center, text width=6.5cm, font=\small, inner sep=5pt},
fleche/.style={-{Stealth[length=3mm]}, thick, popdark}
]
\node[tbox, align=center] (tesis){\textbf{Tesis}\\ La diversidad es una riqueza, no un problema};
\node[abox, below=of tesis, align=center] (arg1){\textbf{Argumento 1}\\ La diversidad cultural enriquece la vida de todos (lenguas, fiestas, cocina)};
\node[abox, below=of arg1, align=center] (arg2){\textbf{Argumento 2}\\ Los prejuicios nacen de la ignorancia y del miedo a lo diferente};
\node[cbox, below=of arg2, align=center] (concl){\textbf{Conclusión}\\ Hay que educar a los jóvenes en el respeto y la tolerancia};
\draw[fleche] (tesis) -- (arg1);
\draw[fleche] (arg1) -- (arg2);
\draw[fleche] (arg2) -- (concl);
\end{tikzpicture}
\end{center}
\legende{Structure argumentative du texte : de la thèse à la conclusion.}
\end{popfigure}

\begin{exercice}[À toi de jouer]
Réponds par des phrases complètes en espagnol, avec une citation du texte quand c'est nécessaire.

1. \esp{¿De qué habla el texto?}

2. \esp{¿Cuál es la tesis del autor?}

3. \esp{¿De dónde nacen los prejuicios?}

4. \esp{¿Qué significa « la diversidad es una riqueza »?}

5. \esp{¿Estás de acuerdo con el autor? ¿Por qué?}

6. Vrai ou faux, justifié : \esp{La educación no tiene ningún papel en la lucha contra los prejuicios.}
\end{exercice}

\begin{corrige}[Exercice — Compréhension du texte]
1. \esp{El texto habla de la identidad, la tolerancia y la convivencia entre los grupos étnicos de Camerún.}

2. \esp{La tesis del autor es que la diversidad cultural es una riqueza y que la tolerancia permite vivir juntos en paz.}

3. \esp{Según el texto, « los prejuicios nacen de la ignorancia » y del miedo a lo diferente.}

4. \esp{Significa que la variedad de lenguas, culturas y tradiciones enriquece al país, es decir, que las diferencias son un valor positivo y no un obstáculo.}

5. Réponse personnelle, par exemple : \esp{Sí, estoy de acuerdo con el autor porque en mi colegio hay alumnos de muchas regiones y aprendemos mucho unos de otros.} (Toute réponse justifiée par \esp{porque...} est acceptable.)

6. \textbf{Falso.} \esp{El texto dice: « hay que educar a los jóvenes en el respeto »} : l'éducation a donc un rôle essentiel dans la lutte contre les préjugés.
\end{corrige}

\begin{aretenir}[L'essentiel de la section]
\begin{itemize}
\item Quatre types de questions : globale, de détail, d'interprétation, de réaction personnelle.
\item Toujours répondre par une \textbf{phrase complète en espagnol}, avec une \textbf{citation entre guillemets} pour justifier.
\item \esp{¿Estás de acuerdo?} exige \esp{porque...} : jamais « sí » ou « no » seul.
\item Au vrai ou faux, méfie-toi des mots \esp{solo}, \esp{nunca}, \esp{todos}, \esp{no} : ils inversent souvent le sens.
\item Reformuler la thèse = utiliser ses propres mots avec \esp{el autor afirma que...}, \esp{en otras palabras...}.
\end{itemize}
\end{aretenir}

\coursec{Texte support : « Vivir juntos en la diferencia »}

Cette leçon s'ouvre sur un texte authentique, ancré dans la réalité camerounaise et hispanophone, qui met en scène le vivre-ensemble dans un quartier multiculturel. Lisez-le attentivement, sans chercher à traduire chaque mot du premier coup : visez la compréhension globale.

\begin{textebox}[Vivir juntos en la diferencia]
En el barrio de Briqueterie, en Yaoundé, conviven personas de diferentes regiones y culturas. Algunos hablan ewondo, otros duala, fulfuldé o francés, pero todos comparten la misma escuela, el mismo mercado y el mismo campo de fútbol.

María, una estudiante de Guinea Ecuatorial, llegó el año pasado. Al principio, algunos vecinos tenían prejuicios contra ella porque no conocían su cultura. Pero poco a poco, gracias al diálogo y al respeto mutuo, descubrieron sus costumbres, su música y su cocina. Hoy María está integrada en la comunidad y participa en todas las fiestas del barrio.

«La diversidad no es un problema, es una riqueza», dice el profesor de español. «Debemos respetar las diferencias culturales y rechazar la discriminación. La convivencia enriquece a todos.»

Sin embargo, la intolerancia todavía existe en algunas partes. Cuando alguien excluye a otra persona por su lengua, su etnia o su religión, divide la sociedad. Por eso la escuela organiza cada año una semana de la tolerancia, con debates, canciones y platos de todas las regiones. Los alumnos aprenden a escuchar, a comprender y a aceptar al otro. Como dice un proverbio africano: «Un solo dedo no puede coger una piedra.» Juntos somos más fuertes.
\end{textebox}

\textbf{Traduction / Glossaire des mots difficiles}

\begin{center}
\begin{tabularx}{\linewidth}{|X|X|}
\hline
\rowcolor{popblueL} Español & Français \\
\hline
el barrio & le quartier \\
\hline
conviviven & vivent ensemble (verbe \esp{convivir}, 3e pers. pl.) \\
\hline
compartir & partager \\
\hline
las costumbres & les coutumes \\
\hline
el prejuicio & le préjugé \\
\hline
integrada & intégrée (participe passé adjectivé, accord au fém.) \\
\hline
la riqueza & la richesse \\
\hline
rechazar & rejeter \\
\hline
la etnia & l'ethnie \\
\hline
excluye & exclut (verbe \esp{excluir}, 3e pers. sg., changement \esp{u} $\to$ \esp{ue}) \\
\hline
coger & attraper, ramasser \\
\hline
juntos & ensemble \\
\hline
\end{tabularx}
\end{center}

\coursec{Compréhension du texte}

\courssub{Compréhension globale}

\begin{enumerate}
\item De quel quartier et de quelle ville parle le texte ? Quelles langues y entend-on ?
\item Quel événement a perturbé l'harmonie initiale ? Comment la situation a-t-elle été résolue ?
\item Quel est le message principal du professeur de espagnol ? Citez une phrase du texte.
\item Quel rôle joue l'école pour renforcer la tolérance ?
\item Expliquez le sens du proverbe africain cité à la fin : « Un solo dedo no puede coger una piedra. »
\end{enumerate}

\courssub{Compréhension détaillée}

\begin{enumerate}
\item Trouvez dans le texte trois verbes conjugués au présent de l'indicatif qui expriment une action habituelle ou une vérité générale.
\item Relevez deux connecteurs logiques (mots de liaison) qui marquent l'opposition ou la concession.
\item Quel temps verbal est utilisé pour « \esp{llegó} » et « \esp{descubrieron} » ? Pourquoi ce temps ? (Rappel : actions ponctuelles, passées, terminées).
\item Identifiez la structure utilisée pour exprimer l'obligation morale dans la phrase : « \esp{Debemos respetar las diferencias...} ».
\end{enumerate}

\courssub{Repérage lexical (transition vers la section lexique)}

\begin{enumerate}
\item Surlignez dans le texte tous les mots de la famille de \esp{tolerancia} (nom, adjectif, verbe).
\item Trouvez le synonyme dans le texte de : « rejeter » ; « s'enrichir » ; « communauté ».
\item Quels sont les deux verbes pronominaux (réfléchis) présents dans le dernier paragraphe ? (\esp{se} + infinitif).
\end{enumerate}

\begin{corrige}[Compréhension du texte]
\textbf{Globale :} 1) Briqueterie, Yaoundé. Langues : ewondo, duala, fulfuldé, français. 2) L'arrivée de María (Guinée équatoriale) et les préjugés des voisins. Résolution par le dialogue et le respect mutuel. 3) « La diversidad no es un problema, es una riqueza. » / « Debemos respetar las diferencias culturales y rechazar la discriminación. » 4) Elle organise une semaine de la tolérance (débats, chansons, plats). 5) L'union fait la force / seul on ne peut rien faire, ensemble on est fort.
\textbf{Détaillée :} 1) \esp{conviven, comparten, hablan, organizan, aprenden, aceptan, dice, existe, divide, somos}. 2) \esp{Sin embargo} (cependant), \esp{Pero} (mais), \esp{Por eso} (conséquence). 3) Prétérit indéfini (\esp{pretérito indefinido}) : actions passées datées et terminées. 4) \esp{Deber} + infinitif = obligation morale (devoir).
\textbf{Lexical :} 1) \esp{tolerancia, tolerante, intolerancia}. 2) \esp{rechazar} ; \esp{enriquecer(se)} ; \esp{comunidad} (ou \esp{barrio, vecinos}). 3) \esp{integrarse} (intégré \esp{está integrada} participe), \esp{enriquecerse} (\esp{enriquece} 3e sg présent non pronominal ici mais sens réfléchi possible), \esp{aceptar} (non pronominal). Attention : \esp{integrarse} apparaît sous forme participle \esp{integrada}. Verbe pronominal infinitif dans le texte : \esp{aceptar al otro} (non pronominal). Le texte ne contient pas de verbe pronominal à l'infinitif explicite hormis peut-être l'idée de \esp{se integrar} (implicite). On notera \esp{se integró} (passé) non présent. La question porte sur le dernier paragraphe : \esp{aprenden a escuchar, a comprender y a aceptar} (infinitifs non pronominaux). Il n'y en a pas. Je corrige la question : « Identifiez les verbes à l'infinitif précédés de préposition (\esp{a}) ». 
\end{corrige}

\coursec{Lexique thématique : identidad, tolerancia y convivencia}

Dans la section précédente, nous avons lu et compris le texte \esp{Vivir juntos en la diferencia}. Ce texte contenait de nombreux mots nouveaux : \esp{identidad}, \esp{tolerancia}, \esp{convivencia}, \esp{prejuicio}... Pour bien parler du vivre-ensemble — à l'écrit comme à l'oral — il faut d'abord maîtriser ce vocabulaire. C'est l'objectif de cette section : construire, organiser et mémoriser le lexique de l'identité, de la tolérance et de la convivencia.

\courssub{Observation : les mots dans leur contexte}

Avant d'apprendre des listes de mots, observons-les dans un texte. Un bon vocabulaire ne s'apprend jamais isolément : il s'apprend dans des phrases. Le texte ci-dessus (\esp{Vivir juntos en la diferencia}) sert de corpus d'observation principal.

\textbf{Questions de repérage lexical (reprise)}

\begin{enumerate}
\item Relevez dans le texte trois mots de la famille de \esp{tolerancia}.
\item Quel mot du texte signifie « jugement formé avant de connaître » ?
\item Trouvez dans le texte deux verbes qui expriment le vivre-ensemble et deux qui expriment le rejet.
\end{enumerate}

\begin{corrige}[Repérage lexical]
1) \esp{tolerancia, tolerante, intolerancia}. 2) \esp{prejuicio}. 3) Vivre-ensemble : \esp{conviven, comparten, integrada} (participe), \esp{enriquece}. Rejet : \esp{excluye, divide, rechazan} (implicite dans \esp{rechazar}).
\end{corrige}

\courssub{Champ lexical 1 : la identidad}

L'identité, c'est ce qui fait qu'une personne ou un peuple est \emph{lui-même} et pas un autre : sa langue, sa culture, ses traditions.

\begin{center}
\begin{tabularx}{\linewidth}{|X|X|X|}
\hline
\rowcolor{popblueL} Español & Français & Exemple \\
\hline
\cle{la identidad} & l'identité & \esp{La lengua es parte de mi identidad.} \\
\hline
\cle{la cultura} & la culture & \esp{La cultura bétí es muy rica.} \\
\hline
\cle{la lengua} & la langue & \esp{El español es una lengua mundial.} \\
\hline
\cle{la tradición} & la tradition & \esp{Es una tradición de mi pueblo.} \\
\hline
\cle{las costumbres} & les coutumes & \esp{Respeto las costumbres de mis vecinos.} \\
\hline
\cle{la etnia} & l'ethnie & \esp{Camerún tiene más de doscientas etnias.} \\
\hline
\cle{la nacionalidad} & la nationalité & \esp{¿Cuál es tu nacionalidad?} \\
\hline
\cle{la raíz} (pl. \esp{las raíces}) & la racine & \esp{No debemos olvidar nuestras raíces.} \\
\hline
\end{tabularx}
\end{center}

\begin{exemplebox}[Exemples traduits]
\esp{Las raíces culturales nos dan identidad.} « Les racines culturelles nous donnent une identité. »

\esp{En mi escuela hay alumnos de muchas etnias diferentes.} « Dans mon école, il y a des élèves de nombreuses ethnies différentes. »

\esp{Cada región de Camerún tiene sus tradiciones y sus costumbres.} « Chaque région du Cameroun a ses traditions et ses coutumes. »
\end{exemplebox}

\begin{attention}[Faux ami / Piège orthographique : « la raíz »]
\esp{La raíz} signifie bien « la racine », mais attention au pluriel : \esp{raíz} devient \esp{raíces} (le \esp{z} devient \esp{c} devant \esp{e} pour conserver le son [θ] / [s]). On écrit donc : \esp{una raíz} / \esp{muchas raíces}. N'écrivez jamais \esp{*las raízes} (avec \esp{z}).
\end{attention}

\courssub{Champ lexical 2 : la tolerancia y el respeto}

La tolérance, c'est accepter que l'autre soit différent ; le respect, c'est le traiter avec considération.

\begin{center}
\begin{tabularx}{\linewidth}{|X|X|X|}
\hline
\rowcolor{popgreenL} Español & Français & Exemple \\
\hline
\cle{la tolerancia} & la tolérance & \esp{La tolerancia es necesaria en la escuela.} \\
\hline
\cle{el respeto} & le respect & \esp{El respeto es la base de la convivencia.} \\
\hline
\cle{aceptar} & accepter & \esp{Acepto las opiniones de los demás.} \\
\hline
\cle{respetar} & respecter & \esp{Respetamos a todas las personas.} \\
\hline
\cle{escuchar} & écouter & \esp{Hay que escuchar antes de juzgar.} \\
\hline
\cle{comprender} & comprendre & \esp{Comprendo tu punto de vista.} \\
\hline
\cle{el diálogo} & le dialogue & \esp{El diálogo resuelve los conflictos.} \\
\hline
\cle{la paz} (pl. \esp{las paces}) & la paix & \esp{Vivimos en paz con nuestros vecinos.} \\
\hline
\end{tabularx}
\end{center}

\begin{exemplebox}[Exemple traduit]
\esp{Debemos respetar las diferencias culturales.} « Nous devons respecter les différences culturelles. »

Remarque : \esp{debemos} (verbe \esp{deber} + infinitif) exprime l'obligation morale, comme « nous devons » en français.
\end{exemplebox}

\courssub{Champ lexical 3 : la convivencia}

\esp{La convivencia} est un mot très fréquent en espagnol : il vient du verbe \esp{convivir} (\esp{con-} = « avec » + \esp{vivir} = « vivre »), donc littéralement « vivre avec ».

\begin{center}
\begin{tabularx}{\linewidth}{|X|X|X|}
\hline
\rowcolor{poporangeL} Español & Français & Exemple \\
\hline
\cle{la convivencia} & le vivre-ensemble & \esp{La convivencia es posible.} \\
\hline
\cle{la integración} & l'intégration & \esp{La integración de los alumnos nuevos es importante.} \\
\hline
\cle{la solidaridad} & la solidarité & \esp{La solidaridad une a las personas.} \\
\hline
\cle{la diversidad} & la diversité & \esp{La diversidad es una riqueza.} \\
\hline
\cle{la armonía} & l'harmonie & \esp{Vivimos en armonía.} \\
\hline
\cle{la comunidad} & la communauté & \esp{Nuestra comunidad es muy unida.} \\
\hline
\cle{compartir} & partager & \esp{Compartimos la comida y la alegría.} \\
\hline
\end{tabularx}
\end{center}

\begin{attention}[Faux ami : « compartir »]
\esp{Compartir} signifie « partager » (au sens de « mettre en commun »), pas « comparer ». « Comparer » se dit \esp{comparar}. \esp{Comparto mi merienda con mis amigos.} « Je partage mon goûter avec mes amis. »
\end{attention}

\courssub{Champ lexical 4 : los obstáculos}

Le vivre-ensemble a des ennemis : les préjugés, la discrimination, le rejet. Il faut savoir les nommer pour les combattre.

\begin{definition}[Prejuicio]
\esp{El prejuicio} : \emph{juicio u opinión formada antes de conocer a una persona o una cultura}. En français : un préjugé, c'est-à-dire un jugement formé avant de connaître. Exemple : \esp{Los prejuicios nacen de la ignorancia.} « Les préjugés naissent de l'ignorance. »
\end{definition}

\begin{center}
\begin{tabularx}{\linewidth}{|X|X|X|}
\hline
\rowcolor{poppinkL} Español & Français & Exemple \\
\hline
\cle{el prejuicio} & le préjugé & \esp{No debemos tener prejuicios.} \\
\hline
\cle{la discriminación} & la discrimination & \esp{La discriminación divide la sociedad.} \\
\hline
\cle{el racismo} & le racisme & \esp{Luchamos contra el racismo.} \\
\hline
\cle{el rechazo} & le rejet & \esp{El rechazo del otro es un error.} \\
\hline
\cle{la intolerancia} & l'intolérance & \esp{La intolerancia destruye la paz.} \\
\hline
\cle{excluir} & exclure & \esp{No hay que excluir a nadie.} \\
\hline
\cle{el conflicto} & le conflit & \esp{El diálogo evita los conflictos.} \\
\hline
\end{tabularx}
\end{center}

\begin{exemplebox}[Exemple traduit]
\esp{La discriminación divide la sociedad.} « La discrimination divise la société. »
\end{exemplebox}

\begin{popfigure}
\begin{center}
\begin{tikzpicture}
\node[draw=popgreen, fill=popgreenL, rounded corners, thick, align=center, minimum width=5.5cm, minimum height=4.6cm] (pos) at (0,0) {};
\node[draw=popgreen, fill=popgreen, text=white, rounded corners, align=center, minimum width=5.5cm] at (0,2.0) {\textbf{Valores positivos}};
\node[align=center, text=popdark, font=\small] at (0,0.1) {la tolerancia\\[2pt] el respeto\\[2pt] la integración\\[2pt] la solidaridad\\[2pt] el diálogo};
\node[draw=poppink, fill=poppinkL, rounded corners, thick, align=center, minimum width=5.5cm, minimum height=4.6cm] (neg) at (8,0) {};
\node[draw=poppink, fill=poppink, text=white, rounded corners, align=center, minimum width=5.5cm] at (8,2.0) {\textbf{Obstáculos}};
\node[align=center, text=popdark, font=\small] at (8,0.1) {el prejuicio\\[2pt] la discriminación\\[2pt] el rechazo\\[2pt] la intolerancia\\[2pt] el conflicto};
\draw[-{Stealth[length=3mm]}, very thick, poporange] (2.9,0) -- (5.1,0) node[midway, above, text=popdark]{\small luchar};
\end{tikzpicture}
\end{center}
\legende{Les valeurs du vivre-ensemble face aux obstacles : il faut « luchar » (lutter) pour faire passer la société de la colonne rose à la colonne verte.}
\end{popfigure}

\courssub{Verbes et adjectifs associés}

\textbf{Les verbes du vivre-ensemble}

\begin{center}
\begin{tabularx}{\linewidth}{|X|X|X|}
\hline
\rowcolor{popblueL} Verbo & Français & Exemple \\
\hline
\cle{convivir} & vivre ensemble & \esp{Convivimos con personas de otras culturas.} \\
\hline
\cle{integrarse} & s'intégrer & \esp{María se integró rápidamente en la clase.} \\
\hline
\cle{diferenciarse} & se différencier & \esp{Cada cultura se diferencia por sus costumbres.} \\
\hline
\cle{discriminar} & discriminer & \esp{Es injusto discriminar a una persona.} \\
\hline
\cle{enriquecerse} & s'enrichir & \esp{Nos enriquecemos con la diversidad.} \\
\hline
\end{tabularx}
\end{center}

\begin{attention}[Verbes pronominaux : le pronom est obligatoire]
\esp{Integrarse}, \esp{diferenciarse} et \esp{enriquecerse} sont des verbes \emph{pronominaux} : ils se conjuguent avec \esp{me, te, se, nos, os, se}. On dit \esp{yo me integro}, jamais \esp{*yo integro} dans ce sens (s'intégrer). C'est comme le français « je \emph{m}'intègre » : le pronom réfléchi est obligatoire. \esp{Discriminar} n'est pas pronominal.
\end{attention}

\textbf{Les adjectifs}

\begin{center}
\begin{tabularx}{\linewidth}{|X|X|X|}
\hline
\rowcolor{popgreenL} Adjetivo & Français & Exemple \\
\hline
\cle{tolerante} / \cle{intolerante} & tolérant / intolérant & \esp{Mi profesor es muy tolerante.} \\
\hline
\cle{diferente} & différent & \esp{Somos diferentes pero iguales en derechos.} \\
\hline
\cle{diverso, -a} & divers, -e & \esp{Camerún es un país muy diverso.} / \esp{Una cultura diversa.} \\
\hline
\cle{multicultural} & multiculturel(le) & \esp{Yaoundé es una ciudad multicultural.} \\
\hline
\cle{respetuoso, -a} & respectueux/-se & \esp{Hay que ser respetuoso con todos.} \\
\hline
\end{tabularx}
\end{center}

\begin{attention}[« Diferente » et « Tolerante » : invariables en genre]
Les adjectifs \esp{diferente}, \esp{tolerante}, \esp{intolerante}, \esp{multicultural} s'écrivent \textbf{sans accent} et sont \textbf{invariables en genre} : on dit \esp{un país diferente} (masc.) et \esp{una cultura diferente} (fém.) — jamais \esp{*diferenta}. Ils ne varient qu'au pluriel : \esp{países diferentes}, \esp{culturas diferentes}. 
\emph{Attention :} \esp{diverso} et \esp{respetuoso} \textbf{varient en genre} (\esp{diverso/diversa}, \esp{respetuoso/respetuosa}). Ne les confondez pas.
\end{attention}

\courssub{La formation des contraires}

L'espagnol, comme le français, fabrique des contraires avec des préfixes ou des expressions négatives. C'est un excellent moyen d'enrichir vite son vocabulaire.

\begin{propriete}[Les contraires du lexique de la tolérance]
\begin{itemize}
\item Préfixe \esp{in-} (devant \esp{t-}, \esp{r-}, \esp{l-}, \esp{m-}, \esp{n-}, \esp{s-}, \esp{p-}...) : \esp{tolerancia} $\to$ \esp{intolerancia} ; \esp{tolerante} $\to$ \esp{intolerante} ; \esp{justo} $\to$ \esp{injusto}.
\item Expression \esp{falta de} + nom : \esp{respeto} $\to$ \esp{falta de respeto} (« manque de respect ») ; \esp{diálogo} $\to$ \esp{falta de diálogo} ; \esp{tolerancia} $\to$ \esp{falta de tolerancia}.
\item Préfixe \esp{des-} : \esp{acuerdo} $\to$ \esp{desacuerdo} (« désaccord ») ; \esp{igual} $\to$ \esp{desigual} (« inégal, différent ») ; \esp{honesto} $\to$ \esp{deshonesto}.
\end{itemize}
\end{propriete}

\begin{exemplebox}[Exemples traduits]
\esp{La intolerancia es el enemigo de la paz.} « L'intolérance est l'ennemie de la paix. »

\esp{La falta de respeto crea conflictos.} « Le manque de respect crée des conflits. »

\esp{Estoy en desacuerdo con esa opinión.} « Je suis en désaccord avec cette opinion. »
\end{exemplebox}

\begin{methode}[Mémoriser le lexique par familles de mots]
Pour retenir durablement ce vocabulaire, classez chaque mot en famille : \textbf{nom $\to$ verbe $\to$ adjectif}. Un seul effort de mémoire vous donne trois mots !
\begin{enumerate}
\item \esp{tolerancia} (nom) $\to$ \esp{tolerar} (verbe) $\to$ \esp{tolerante} (adjectif).
\item \esp{respeto} $\to$ \esp{respetar} $\to$ \esp{respetuoso}.
\item \esp{integración} $\to$ \esp{integrar(se)} $\to$ \esp{integrado}.
\item \esp{discriminación} $\to$ \esp{discriminar} $\to$ \esp{discriminatorio}.
\item \esp{diferencia} $\to$ \esp{diferenciar(se)} $\to$ \esp{diferente}.
\end{enumerate}
Faites des fiches ou une carte mentale avec ces familles, puis rédigez une phrase avec chaque famille : \esp{Una persona tolerante tolera la diferencia y practica la tolerancia.}
\end{methode}

\begin{exoresolu}[Classer et compléter]
\textbf{Énoncé.} 1) Classez les mots suivants en deux colonnes « valores » et « obstáculos » : \esp{solidaridad, racismo, diálogo, prejuicio, armonía, rechazo}. 2) Complétez avec le mot de la famille indiquée (nom, verbe ou adjectif selon la syntaxe) : \esp{La \_\_\_\_\_\_ (tolerar) es necesaria para la convivencia.} ; \esp{Es importante \_\_\_\_\_\_ (integración) a los alumnos nuevos.}
\tcblower
\textbf{Solution.} 1) Valores : \esp{solidaridad, diálogo, armonía}. Obstáculos : \esp{racismo, prejuicio, rechazo}. 2) \esp{La tolerancia es necesaria para la convivencia} (on attend le nom, car l'article \esp{la} précède). \esp{Es importante integrar a los alumnos nuevos} (après \esp{es importante}, on utilise l'infinitif ; attention, \esp{integración} est le nom, le verbe est \esp{integrar}).
\end{exoresolu}

\begin{savaistu}[El Día Internacional de la Tolerancia]
Chaque année, le 16 novembre, l'Espagne et de nombreux pays hispanophones célèbrent le \esp{Día Internacional de la Tolerancia}, proclamé par l'UNESCO en 1995. Dans les écoles espagnoles, ce jour est l'occasion de débats, d'expositions et d'activités sur le respect de l'autre. Au Cameroun, pays de plus de deux cents groupes ethniques où cohabitent le français, l'anglais et des centaines de langues locales, la tolérance et la \esp{convivencia} sont aussi des valeurs essentielles de la vie quotidienne : elles font partie de ce que l'on appelle le « vivre-ensemble ».
\end{savaistu}

\begin{exercice}[Entraînement lexical]
1) Traduisez en espagnol : « Nous devons respecter les différences culturelles. » ; « L'intolérance divise la société. » ; « Le dialogue est la base du vivre-ensemble. »
2) Trouvez le contraire de : \esp{tolerancia}, \esp{respeto}, \esp{acuerdo}, \esp{igual}.
3) Complétez avec le mot approprié (nom, verbe ou adjectif de la famille) : \esp{Mi vecino es muy \_\_\_\_\_\_ (tolerancia), siempre \_\_\_\_\_\_ (tolerar) mis errores.} ; \esp{La \_\_\_\_\_\_ (integrar) de los refugiados es un derecho humano.}
4) Rédigez trois phrases sur votre quartier ou votre école en utilisant : \esp{convivir}, \esp{diversidad}, \esp{respetuoso}.
\end{exercice}

\begin{corrige}[Exercice : Entraînement lexical]
1) \esp{Debemos respetar las diferencias culturales.} ; \esp{La intolerancia divide la sociedad.} ; \esp{El diálogo es la base de la convivencia.}
2) \esp{tolerancia} $\to$ \esp{intolerancia} ; \esp{respeto} $\to$ \esp{falta de respeto} (ou \esp{irrespeto}, rare) ; \esp{acuerdo} $\to$ \esp{desacuerdo} ; \esp{igual} $\to$ \esp{desigual}.
3) \esp{Mi vecino es muy tolerante, siempre tolera mis errores.} ; \esp{La integración de los refugiados es un derecho humano.}
4) Exemple : \esp{En mi barrio conviven personas de muchas regiones. La diversidad es una riqueza. Debemos ser respetuosos con todos.}
\end{corrige}

\coursec{Exprimer l'accord et le désaccord}

Maintenant que le lexique est en place, nous pouvons l'utiliser pour interagir : donner son opinion, approuver ou contester celle d'autrui. C'est une compétence clé de l'expression orale et écrite (débat, commentaire, lettre d'opinion).

\courssub{Les structures de base : \esp{estar de acuerdo} / \esp{no estar de acuerdo}}

En espagnol, « être d'accord » ne se dit pas avec le verbe \esp{ser} mais avec \esp{estar} : \esp{estar de acuerdo}. C'est un état, une position prise à un moment donné, pas une caractéristique permanente.

\begin{center}
\begin{tabularx}{\linewidth}{|X|X|X|}
\hline
\rowcolor{popblueL} Structure & Français & Niveau de langue \\
\hline
\cle{Estoy de acuerdo.} & Je suis d'accord. & Neutre / Standard \\
\hline
\cle{No estoy de acuerdo.} & Je ne suis pas d'accord. & Neutre / Standard \\
\hline
\cle{Estoy totalmente de acuerdo.} & Je suis tout à fait d'accord. & Insistance \\
\hline
\cle{No estoy nada de acuerdo.} & Je ne suis pas du tout d'accord. & Insistance \\
\hline
\cle{Estoy de acuerdo en parte.} / \cle{Estoy parcialmente de acuerdo.} & Je suis d'accord en partie. / Partiellement. & Nuance \\
\hline
\end{tabularx}
\end{center}

\begin{attention}[Piège majeur : \esp{Ser} vs \esp{Estar} avec \esp{acuerdo}]
\textbf{Jamais} \esp{*Soy de acuerdo}. \textbf{Toujours} \esp{Estoy de acuerdo}.
\emph{Raison :} \esp{Acuerdo} exprime un état résultant d'une décision, d'une opinion du moment $\to$ \esp{estar}. \esp{Ser de acuerdo} n'existe pas en espagnol standard.
\emph{Exemple :} \esp{—¿Estás de acuerdo con la ley? —Sí, estoy de acuerdo.}
\end{attention}

\courssub{Autres expressions pour nuancer son opinion}

\begin{center}
\begin{tabularx}{\linewidth}{|X|X|X|}
\hline
\rowcolor{popgreenL} Expression & Français & Usage \\
\hline
\cle{Me parece bien / mal.} & Cela me semble bien / mal. & Opinion personnelle, poli. \\
\hline
\cle{Creo que sí / que no.} & Je crois que oui / que non. & Opinion subjective. \\
\hline
\cle{En mi opinión,} / \cle{Desde mi punto de vista,} & À mon avis, / De mon point de vue, & Introduction formelle. \\
\hline
\cle{Opino que...} & J'estime que... / Je pense que... & Verbe \esp{opinar}, + indicatif. \\
\hline
\cle{Es verdad, pero...} / \cle{Es cierto, sin embargo...} & C'est vrai, mais... / C'est certain, cependant... & Concession + opposition. \\
\hline
\cle{No lo veo claro.} & Je n'y vois pas clair / Je ne suis pas convaincu. & Doute, désaccord doux. \\
\hline
\cle{¡Ni hablar!} / \cle{¡De ninguna manera!} & Pas question ! / Absolument pas ! & Désaccord fort, familier / ferme. \\
\hline
\end{tabularx}
\end{center}

\courssub{Grammaire : Les connecteurs logiques pour argumenter}

Pour structurer un débat ou un commentaire, il faut articuler les idées.

\begin{center}
\begin{tabularx}{\linewidth}{|X|X|X|}
\hline
\rowcolor{poporangeL} Connecteur & Français & Fonction \\
\hline
\cle{Porque} / \cle{Como} / \cle{Puesto que} & Parce que / Comme / Puisque & Cause (indicatif) \\
\hline
\cle{Por lo tanto} / \cle{Por eso} / \cle{Así que} & Par conséquent / C'est pourquoi / Donc & Conséquence \\
\hline
\cle{Pero} / \cle{Sino} / \cle{Sin embargo} / \cle{No obstante} & Mais / Mais (après négation) / Cependant / Toutefois & Opposition / Concession \\
\hline
\cle{Además} / \cle{También} / \cle{Asimismo} & De plus / Aussi / Également & Addition \\
\hline
\cle{Por ejemplo} / \cle{Es decir} / \cle{O sea} & Par exemple / C'est-à-dire / Autrement dit & Illustration / Explication \\
\hline
\end{tabularx}
\end{center}

\begin{exemplebox}[Mini-dialogue : Débat sur la semaine de la tolérance]
\esp{— ¿Qué opinas de la semana de la tolerancia en el colegio?} \\
\esp{— \textbf{En mi opinión}, \textbf{es una excelente idea}. \textbf{Porque} permite conocer otras culturas.} \\
\esp{— \textbf{No estoy de acuerdo}. \textbf{Creo que} una semana no basta; \textbf{por lo tanto}, hay que trabajar la tolerancia todo el año.} \\
\esp{— \textbf{Es verdad, pero} es un buen comienzo. \textbf{Además}, los debates ayudan a reflexionar.} \\
\esp{— \textbf{Estoy de acuerdo en parte}. \textbf{Sin embargo}, si no hay seguimiento, \textbf{no sirve de nada}.}
\end{exemplebox}

\begin{methode}[Participer à un débat oral (Expression orale)]
\begin{enumerate}
\item \textbf{Écoutez} la thèse de l'autre (\esp{Escuchar activamente}).
\item \textbf{Reformulez} pour vérifier : \esp{¿Quieres decir que...?} (« Veux-tu dire que... ? »).
\item \textbf{Positionnez-vous} : \esp{Estoy de acuerdo / No estoy de acuerdo / Estoy de acuerdo en parte.}
\item \textbf{Argumentez} avec \esp{porque, ya que, puesto que} + indicatif (faits) ou \esp{es importante que} + \textbf{subjonctif} (valeur, nécessité).
\item \textbf{Concluez} : \esp{En conclusión, / Por todo ello, / Por tanto.}
\end{enumerate}
\end{methode}

\begin{exoresolu}[Transformer en espagnol]
\textbf{Énoncé.} Traduisez cet échange :
A : « À mon avis, la diversité culturelle est une richesse. »
B : « Je ne suis pas d'accord. Je pense que la diversité crée des problèmes. »
A : « Ce n'est pas vrai. Par exemple, dans mon école, nous vivons bien ensemble. »
B : « C'est vrai, mais il y a des préjugés. Donc, il faut plus de dialogue. »
\tcblower
\textbf{Solution.}
A : \esp{En mi opinión, la diversidad cultural es una riqueza.}
B : \esp{No estoy de acuerdo. Creo que la diversidad crea problemas.}
A : \esp{No es verdad. Por ejemplo, en mi escuela convivimos bien.}
B : \esp{Es cierto, pero hay prejuicios. Por lo tanto, hace falta más diálogo.} (\esp{Hace falta} = il faut, impersonnel).
\end{exoresolu}

\begin{exercice}[Expression de l'accord/désaccord]
1) Complétez avec \esp{estoy / no estoy / está / no está} : \esp{\_\_\_\_\_\_ de acuerdo con la discriminación.} ; \esp{Mi profesor \_\_\_\_\_\_ de acuerdo con la intolerancia.}
2) Répondez par \esp{Estoy de acuerdo} ou \esp{No estoy de acuerdo} et justifiez en une phrase (espagnol) :
   a) \esp{La discriminación enriquece a la sociedad.}
   b) \esp{El diálogo resuelve los conflictos.}
   c) \esp{Los prejuicios nacen de la ignorancia.}
3) Rédigez un mini-dialogue de 6 répliques (3 par personne) sur le thème : « Faut-il interdire les symboles religieux à l'école ? » Utilisez au moins 3 connecteurs logiques et 3 expressions d'opinion du tableau.
\end{exercice}

\begin{corrige}[Exercice : Expression de l'accord/désaccord]
1) \esp{No estoy de acuerdo con la discriminación.} (Je ne suis pas d'accord). \esp{Mi profesor no está de acuerdo con la intolerancia.} (Il n'est pas d'accord).
2) a) \esp{No estoy de acuerdo. La discriminación divide y destruye la convivencia.} b) \esp{Estoy de acuerdo. El diálogo permite entender al otro.} c) \esp{Estoy de acuerdo. Quien no conoce, juzga mal.}
3) Exemple :
A: \esp{En mi opinión, no se deben prohibir los símbolos religiosos. Por eso, la libertad es un derecho.}
B: \esp{No estoy de acuerdo. Pienso que la escuela debe ser laica. Por lo tanto, nada de símbolos.}
A: \esp{Es cierto, pero la laicidad no es prohibir, es respetar. Además, la diversidad enriquece.}
B: \esp{No lo veo claro. Sin embargo, acepto tu punto de vista.}
\end{corrige}

\coursec{Méthode du commentaire de texte et commentaire modèle}

Le commentaire de texte est un exercice d'examen incontournable (Baccalauréat). Il ne s'agit pas de résumer, mais d'\textbf{analyser} comment le texte construit son sens sur le thème imposé.

\begin{methode}[Étapes du commentaire de texte (Type Bac)]
\begin{enumerate}
\item \textbf{Lecture active (10 min) :} Identifiez le thème, la thèse, le ton, la structure (paragraphes), le vocabulaire clé, les connecteurs.
\item \textbf{Introduction (rédaction soignée) :}
   \begin{itemize}
   \item \emph{Accroche} : Citation, proverbe, fait de société lié au thème (identité, tolérance).
   \item \emph{Présentation du texte} : Titre (si connu), auteur (ici : document pédagogique original), source, date, thème général.
   \item \emph{Problématique} : Question à laquelle le texte répond (ex. : \emph{Comment le texte montre-t-il que la tolérance est la condition du vivre-ensemble ?}).
   \item \emph{Annonce du plan} : 2 ou 3 grandes parties logiques.
   \end{itemize}
\item \textbf{Développement (2 ou 3 parties, 1 paragraphe par idée forte) :}
   \begin{itemize}
   \item \textbf{Idée directrice} (phrase d'intro du paragraphe).
   \item \textbf{Citation courte} (intégrée dans la phrase : \esp{« ... »}).
   \item \textbf{Analyse} : Expliquez le sens, la portée, le vocabulaire, la grammaire (ex. : emploi du subjonctif, passé, connecteurs).
   \item \textbf{Lien} vers l'idée suivante.
   \end{itemize}
\item \textbf{Conclusion (synthèse + ouverture) :}
   \begin{itemize}
   \item Reprise de la thèse du texte (réponse à la problématique).
   \item Ouverture : lien avec l'actualité, autre texte, expérience personnelle (Cameroun, Guinée équatoriale), citation.
   \end{itemize}
\item \textbf{Relecture (5 min) :} Orthographe, accents, accords, cohérence, français correct.
\end{enumerate}
\end{methode}

\begin{textebox}[Commentaire modèle : « Vivir juntos en la diferencia »]
\textbf{Introduction}
« Un solo dedo no puede coger una piedra. » Ce proverbe africain, cité en conclusion du texte \emph{Vivir juntos en la diferencia}, résume à lui seul la thèse centrale : le vivre-ensemble ne se décrète pas, il se construit patiemment par le dialogue et le respect mutuel. Ce document, ancré dans la réalité du quartier Briqueterie à Yaoundé, met en scène l'arrivée de María, jeune Équato-Guinéenne, et la transformation des préjugés en intégration. Nous nous demanderons : \emph{De quelle manière le texte illustre-t-il le passage de l'intolérance à la convivencia comme richesse collective ?} Nous verrons d'abord que le texte pose le diagnostic d'une société fracturée par les préjugés (I), puis qu'il propose l'école et le dialogue comme leviers d'une convivencia active (II), pour enfin montrer que la diversité, loin d'être une menace, est présentée comme une richesse partagée (III).

\textbf{I. Le diagnostic : une société menacée par l'intolérance et les préjugés}
Le texte s'ouvre sur un tableau idyllique : \esp{« conviven personas de diferentes regiones y culturas »}. Le verbe \esp{convivir} (présent d'habitude) ancre le vivre-ensemble dans le quotidien. Mais la conjonction \esp{« Sin embargo »} (paragraphe 4) brise cette harmonie : \esp{« la intolerancia todavía existe »}. Le nom \esp{intolerancia} (préfixe \esp{in-} marquant le contraire) et le verbe \esp{excluye} (présent de vérité générale) désignent l'obstacle majeur. L'exemple de María (\esp{« tenían prejuicios contra ella »}) illustre le mécanisme du \esp{prejuicio} : juger sans connaître (\esp{« porque no conocían su cultura »}). Le texte montre que l'exclusion (\esp{excluye}) divise (\esp{divide la sociedad}), créant un cercle vicieux.

\textbf{II. Les leviers : l'école, le dialogue et le respect mutuel}
Face à ce constat, le texte oppose une réponse institutionnelle et humaine. L'école organise \esp{« una semana de la tolerancia »} : cadre pédagogique structuré. Les activités (\esp{debates, canciones, platos}) mobilisent le lexique du partage (\esp{compartir}) et de la découverte. Le professeur incarne l'autorité morale : \esp{« Debemos respetar las diferencias culturales y rechazar la discriminación »}. L'obligation morale (\esp{deber} + infinitif) s'impose comme norme collective. Le processus d'intégration de María suit trois étapes clés : \esp{diálogo} $\to$ \esp{respeto mutuo} $\to$ \esp{integración}. Le verbe pronominal \esp{integrarse} (passé \esp{se integró}) marque l'appropriation subjective de ce processus par le sujet.

\textbf{III. La thèse : la diversité comme richesse, la convivencia comme construction}
La phrase-choc \esp{« La diversidad no es un problema, es una riqueza »} (antithèse \esp{problema/riqueza}) renverse la vision négative. La métaphore finale du doigt et de la pierre (\esp{« Un solo dedo no puede coger una piedra »}) scelle l'argumentation : la force est dans l'union (\esp{« Juntos somos más fuertes »}). Le texte ne présente pas la tolérance comme une tolérance passive (« supporter »), mais comme une \esp{convivencia} active (\esp{enriquece a todos}, \esp{participa en todas las fiestas}).

\textbf{Conclusion}
En définitive, ce texte démontre que la \esp{convivencia} est un chantier permanent, non un acquis. Il ancre cette universalité dans le local (Yaoundé, Guinée équatoriale), faisant écho à la devise du Cameroun : « Paix – Travail – Patrie ». Ouverture : dans un monde globalisé où les replis identitaires menacent, la leçon de Briqueterie rappelle que, comme le disait Machado, \esp{« Caminante, no hay camino, se hace camino al andar »} : le chemin de la tolérance se fait en marchant ensemble.
\end{textebox}

\coursec{Production guidée : le vivre-ensemble au Cameroun}

\courssub{Consigne de rédaction (Type Bac / Évaluation continue)}

\textbf{Sujet :} « Vous êtes délégué(e) de votre classe au Lycée de Briqueterie. Le proviseur vous demande de rédiger un article pour le journal scolaire à l'occasion de la \emph{Semaine de la Tolérance}. Titre imposé : \esp{« Construir la convivencia en la diferencia »}. Longueur : 120–150 mots. »

\textbf{Contraintes :}
\begin{itemize}
\item Utilisez au moins 8 mots du lexique étudié (surlignez-les dans votre copie).
\item Utilisez au moins 3 connecteurs logiques (\esp{porque, por lo tanto, sin embargo, además...}).
\item Exprimez votre opinion personnelle avec \esp{Estoy de acuerdo / En mi opinión / Creo que...}.
\item Structurez : Introduction (contexte) — 2 arguments — Conclusion (appel à l'action).
\end{itemize}

\begin{methode}[Rédaction guidée : Plan type]
\begin{enumerate}
\item \textbf{Introduction (20-30 mots) :} Présentez l'événement (Semaine de la Tolérance), le lieu (votre lycée/quartier), la thèse : la diversité est une chance.
   \textit{Amorce :} \esp{Con motivo de la Semana de la Tolerancia, nuestro liceo celebra la diversidad cultural...}
\item \textbf{Argument 1 : Le diagnostic / Le défi (40-50 mots) :} Les préjugés, l'intolérance existent encore. Citez un exemple concret (langues, origines, religions). Utilisez \esp{prejuicio, intolerancia, excluir, discriminación}.
   \textit{Amorce :} \esp{Sin embargo, no debemos ignorar que la intolerancia y los prejuicios persisten...}
\item \textbf{Argument 2 : La solution / L'action (40-50 mots) :} Le dialogue, l'école, le respect, l'intégration. Montrez ce qui se passe concrètement (débats, repas partagés, sport). Utilisez \esp{diálogo, respeto, integración, convivencia, enriquecerse, solidaridad}.
   \textit{Amorce :} \esp{Por eso, el diálogo y el respeto mutuo son la base de la convivencia...}
\item \textbf{Conclusion (20-30 mots) :} Appel à l'engagement de tous. Proverbe ou phrase forte. Utilisez \esp{Estoy convencido de que / Juntos podemos / Debemos}.
   \textit{Amorce :} \esp{En conclusión, estoy convencido de que la convivencia se construye cada día. ¡Juntos somos más fuertes!}
\end{enumerate}
\end{methode}

\begin{experience}[Atelier d'écriture collaborative (Classe)]
\begin{enumerate}
\item \textbf{Étape 1 (Individuel, 10 min) :} Chaque élève rédige son brouillon en respectant le plan et les contraintes.
\item \textbf{Étape 2 (Binôme, 10 min) :} Échange de copies. Relecture croisée avec grille : \textit{Lexique (8 mots ?) — Connecteurs (3 ?) — Opinion marquée ? — Structure respectée ? — Orthographe/Accents ?} Annotation au crayon vert.
\item \textbf{Étape 3 (Groupe, 15 min) :} Mise en commun au tableau : construction d'une « version idéale » collective en piquant les meilleures phrases de chaque binôme. Le professeur guide la reformulation (subjonctif ? accords ?).
\item \textbf{Étape 4 (Individuel, 5 min) :} Réécriture finale propre sur la copie d'examen.
\end{enumerate}
\end{experience}

\begin{exemplebox}[Production modèle (≈ 140 mots)]
\esp{Con motivo de la Semana de la Tolerancia, nuestro liceo de Briqueteria celebra la diversidad cultural como una riqueza. En mi opinión, la convivencia no es un sueño, es una tarea diaria.}

\esp{Sin embargo, no debemos ignorar que la intolerancia y los prejuicios persisten. A veces, algunos alumnos excluyen a otros por su etnia, su lengua o su religión. La discriminación divide la sociedad y destruye la armonía.}

\esp{Por eso, el diálogo y el respeto mutuo son la base de la convivencia. Durante esta semana, organizamos debates, cantamos canciones y compartimos platos típicos de todas las regiones. La integración de los compañeros nuevos, como María de Guinea Ecuatorial, enriquece a toda la comunidad. Además, la solidaridad nos une más que las diferencias nos separan.}

\esp{En conclusión, estoy convencido de que la tolerancia se aprende practicándola. Como dice el proverbio: « Un solo dedo no puede coger una piedra. » ¡Juntos somos más fuertes! ¡Viva la convivencia!}
\end{exemplebox}

\begin{corrige}[Grille d'auto-évaluation pour l'élève]
Cochez si vous avez respecté :
\begin{itemize}
\item[$\square$] Titre exact : \esp{Construir la convivencia en la diferencia}
\item[$\square$] 120–150 mots (comptez-les)
\item[$\square$] 8 mots de lexique surlignés : \esp{tolerancia, convivencia, diversidad, prejuicio, intolerancia, excluir, discriminación, diálogo, respeto, integración, solidaridad, enriquecer, compartir, armonía, comunidad...}
\item[$\square$] 3 connecteurs : \esp{sin embargo, por eso, además, en conclusión, porque...}
\item[$\square$] 1 marqueur d'opinion : \esp{en mi opinión, estoy convencido de que, creo que...}
\item[$\square$] Structure en 4 paragraphes distincts (saut de ligne).
\item[$\square$] Accords (genre/nombre) et accents vérifiés (\esp{raíces, intolerancia, convivencia, diversidad, étnia}).
\item[$\square$] Ponctuation espagnole : \esp{¿ ? ¡ !} si questions/exclamations directes.
\end{itemize}
\end{corrige}

\begin{savaistu}[Cameroun et Guinée équatoriale : un espace hispanophone partagé]
Le Cameroun et la Guinée équatoriale partagent une frontière et une histoire. La Guinée équatoriale est le seul pays africain où l'espagnol est langue officielle (avec le français et le portugais). De nombreux Équato-Guinéens étudient ou travaillent au Cameroun (Yaoundé, Douala), et des Camerounais y vont pour le commerce ou les études. Le quartier Briqueterie à Yaoundé est un symbole vivant de cette \esp{convivencia} : on y parle espagnol, français, ewondo, fulfuldé, pidgin... Apprendre l'espagnol au Cameroun, c'est donc aussi se doter d'un outil pour mieux vivre avec ses voisins immédiats et comprendre une culture sœur.
\end{savaistu}

\coursec{Exprimer l'accord et le désaccord}

Dans la section précédente, nous avons appris le vocabulaire essentiel du vivre-ensemble : \esp{identidad}, \esp{tolerancia}, \esp{respeto}, \esp{convivencia}, \esp{prejuicio}, \esp{discriminación}. Mais connaître des mots ne suffit pas : un citoyen doit aussi savoir \cle{débattre}, c'est-à-dire dire ce qu'il pense, réagir aux idées des autres et justifier son point de vue. Dans cette section, nous allons apprendre à exprimer l'\cle{accord} et le \cle{désaccord} en espagnol, de manière simple, nuancée ou argumentée. C'est une compétence clé de l'expression orale et écrite au programme de Première, et elle sera très utile pour le commentaire de texte et la production écrite du Baccalauréat.

\courssub{Observation : un débat en classe}

Lisons d'abord ce dialogue original. Deux élèves d'une classe de Première à Yaoundé débattent après la lecture du texte « Vivir juntos en la diferencia ». Observez les expressions en gras : elles servent à donner son avis, à être d'accord ou pas d'accord.

\begin{textebox}[Un debate en clase de español]
\textbf{Profesora:} Hoy vamos a debatir sobre la convivencia. En mi opinión, la diversidad cultural es una riqueza. \textbf{¿Qué piensas tú}, Amina?

\textbf{Amina:} \textbf{Estoy totalmente de acuerdo} contigo, profesora. En el Camerún hay más de doscientas lenguas, y eso es una riqueza enorme. \textbf{Pienso lo mismo}: cada cultura aporta algo diferente.

\textbf{Profesora:} Muy bien. ¿Y tú, Jean? \textbf{¿Estás de acuerdo} con Amina?

\textbf{Jean:} \textbf{Estoy de acuerdo en parte}. \textbf{Es cierto que} la diversidad es una riqueza, \textbf{sin embargo}, a veces las diferencias crean conflictos. \textbf{No creo que} la convivencia \textbf{sea} siempre fácil.

\textbf{Amina:} \textbf{No estoy de acuerdo}, Jean. \textbf{No tienes razón} cuando dices que las diferencias crean conflictos. Los conflictos nacen de los prejuicios, no de las diferencias.

\textbf{Jean:} \textbf{Puede ser, pero}... mira las noticias: hay discriminación en muchos países.

\textbf{Amina:} \textbf{Es verdad que} existe la discriminación, \textbf{pero no creo que sea} inevitable. Con educación y respeto, podemos cambiar las cosas.

\textbf{Profesora:} Excelente debate. ¿Y cuál es tu conclusión, Jean?

\textbf{Jean:} Bueno... \textbf{tienes razón}, Amina. \textbf{Comparto tu opinión}: la tolerancia se aprende, y la escuela es el mejor lugar para aprenderla.

\textbf{Profesora:} \textbf{Exactamente}. La convivencia empieza aquí, en esta clase, cuando escuchamos y respetamos las opiniones de los demás.
\end{textebox}

\textbf{Glossaire :} \esp{debate} = débat ; \esp{en mi opinión} = à mon avis ; \esp{aporta} = apporte ; \esp{nacen de} = naissent de ; \esp{los demás} = les autres ; \esp{se aprende} = s'apprend ; \esp{empieza} = commence.

\textbf{Questions de compréhension :}
\begin{enumerate}
\item ¿Qué piensa Amina de la diversidad cultural?
\item ¿Está Jean totalmente de acuerdo con Amina al principio?
\item Según Amina, ¿de dónde nacen los conflictos?
\item ¿Cambia Jean de opinión al final del debate?
\end{enumerate}

\textbf{Réponses attendues :} 1. Piensa que es una riqueza enorme. 2. No, está de acuerdo solo en parte. 3. Según Amina, nacen de los prejuicios, no de las diferencias. 4. Sí, al final comparte la opinión de Amina.

\courssub{Exprimer l'accord}

\textbf{A. L'accord simple}\par

Pour dire simplement que l'on est d'accord, l'espagnol dispose de plusieurs expressions très fréquentes à l'oral :

\begin{center}
\begin{tabularx}{\linewidth}{|X|X|X|}
\hline
\rowcolor{popblueL} Español & Français & Remarque \\
\hline
\esp{Estoy de acuerdo.} & Je suis d'accord. & expression de base \\
\hline
\esp{Estoy de acuerdo contigo.} & Je suis d'accord avec toi. & \esp{contigo} = avec toi \\
\hline
\esp{Tienes razón.} & Tu as raison. & \esp{tener razón}, pas « être raison » \\
\hline
\esp{Es verdad.} & C'est vrai. & \\
\hline
\esp{Claro. / Claro que sí.} & Bien sûr. / Évidemment. & très oral \\
\hline
\esp{Exactamente. / Exacto.} & Exactement. / Exact. & \\
\hline
\esp{¡Por supuesto!} & Bien sûr ! & accord enthousiaste \\
\hline
\end{tabularx}
\end{center}

\begin{exemplebox}[Exemples traduits]
\esp{—La tolerancia es fundamental en la escuela. —Estoy de acuerdo contigo.} « —La tolérance est fondamentale à l'école. —Je suis d'accord avec toi. »

\esp{—El respeto empieza en casa. —Tienes razón, es verdad.} « —Le respect commence à la maison. —Tu as raison, c'est vrai. »

\esp{—¿Debemos respetar todas las culturas? —¡Claro que sí!} « —Devons-nous respecter toutes les cultures ? —Bien sûr que oui ! »
\end{exemplebox}

\begin{propriete}[La construction de « estar de acuerdo »]
« Être d'accord » se dit \cle{estar de acuerdo}. On utilise le verbe \esp{estar} (et jamais \esp{ser}) parce qu'être d'accord est un \textbf{état d'opinion}, quelque chose qui peut changer. La locution se conjugue comme un verbe ordinaire :

\begin{center}
\begin{tabular}{|l|l|}
\hline
\rowcolor{popgreenL} Personne & Forme \\
\hline
yo & \esp{estoy de acuerdo} \\
\hline
tú & \esp{estás de acuerdo} \\
\hline
él / ella / usted & \esp{está de acuerdo} \\
\hline
nosotros & \esp{estamos de acuerdo} \\
\hline
vosotros & \esp{estáis de acuerdo} \\
\hline
ellos / ustedes & \esp{están de acuerdo} \\
\hline
\end{tabular}
\end{center}

Avec une personne : \esp{estar de acuerdo} + \textbf{con} + pronom : \esp{conmigo} (avec moi), \esp{contigo} (avec toi), \esp{con él}, \esp{con ella}, \esp{con nosotros}, \esp{con ustedes}. Attention : \esp{conmigo} et \esp{contigo} sont des formes spéciales en un seul mot ; pour toutes les autres personnes, on utilise la forme normale (\esp{con él}, \esp{con nosotros}...).
\end{propriete}

\begin{attention}[Erreur fréquente : « *soy de acuerdo »]
Les francophones traduisent mot à mot « je suis d'accord » par \esp{*soy de acuerdo}. C'est une erreur grave : on dit toujours \esp{\textbf{estoy} de acuerdo}. Retenez : une opinion est un état passager, donc \esp{estar}. De même : \esp{*eres de acuerdo} est faux ; on dit \esp{estás de acuerdo}. Autre piège : \esp{*tienes la razón} avec article est incorrect ; on dit \esp{tienes razón}, sans article.
\end{attention}

\textbf{B. L'accord argumenté}\par

Un bon débatteur ne dit pas seulement « d'accord » : il \textbf{justifie}. Pour cela, on ajoute une explication introduite par \esp{porque} (parce que) :

\begin{itemize}
\item \esp{Estoy de acuerdo contigo porque la diversidad es una riqueza.} « Je suis d'accord avec toi parce que la diversité est une richesse. »
\item \esp{Pienso lo mismo.} « Je pense la même chose. »
\item \esp{Comparto tu opinión porque todos somos iguales.} « Je partage ton opinion parce que nous sommes tous égaux. »
\item \esp{Opino igual que tú.} « Je suis du même avis que toi. »
\end{itemize}

Remarquez la structure : \textbf{expression d'accord + porque + justification}. C'est la structure de base de tout débat en espagnol.

\courssub{Exprimer le désaccord}

\textbf{A. Le désaccord simple}\par

\begin{center}
\begin{tabularx}{\linewidth}{|X|X|X|}
\hline
\rowcolor{popblueL} Español & Français & Remarque \\
\hline
\esp{No estoy de acuerdo.} & Je ne suis pas d'accord. & négation devant le verbe \\
\hline
\esp{No tienes razón.} & Tu n'as pas raison. & direct, un peu sec \\
\hline
\esp{Al contrario.} & Au contraire. & \\
\hline
\esp{No es verdad.} & Ce n'est pas vrai. & \\
\hline
\esp{¡Qué va! / ¡Para nada!} & Pas du tout ! & très oral \\
\hline
\end{tabularx}
\end{center}

\begin{exemplebox}[Exemples traduits]
\esp{—Los prejuicios no son peligrosos. —No estoy de acuerdo: al contrario, son muy peligrosos.} « —Les préjugés ne sont pas dangereux. —Je ne suis pas d'accord : au contraire, ils sont très dangereux. »

\esp{—Aquí no existe la discriminación. —No es verdad, existe en muchas partes.} « —Ici, la discrimination n'existe pas. —Ce n'est pas vrai, elle existe dans beaucoup d'endroits. »
\end{exemplebox}

\textbf{B. Le désaccord nuancé : la politesse du débat}\par

Dans un débat respectueux (valeur essentielle du vivre-ensemble !), on évite de blesser l'autre. L'espagnol, comme le français, utilise des formules \textbf{nuancées} : on reconnaît d'abord une part de vérité, puis on objecte.

\begin{itemize}
\item \esp{Puede ser, pero...} « C'est possible, mais... »
\item \esp{Es cierto que la convivencia es difícil, sin embargo, no es imposible.} « Il est vrai que le vivre-ensemble est difficile, cependant il n'est pas impossible. »
\item \esp{Es verdad que hay conflictos, pero también hay solidaridad.} « Il est vrai qu'il y a des conflits, mais il y a aussi de la solidarité. »
\item \esp{Entiendo tu opinión, pero no la comparto.} « Je comprends ton opinion, mais je ne la partage pas. »
\item \esp{Estoy de acuerdo en parte.} « Je suis d'accord en partie. »
\end{itemize}

Structure à retenir : \esp{Es cierto que... / Es verdad que...} + \textbf{indicatif} (on constate un fait), puis \esp{pero / sin embargo} + objection.

\courssub{Le désaccord avec le subjonctif : « no creo que... »}

Observons deux phrases du dialogue :

\esp{Creo que la convivencia es posible.} « Je crois que le vivre-ensemble est possible. » (indicatif : \esp{es})

\esp{No creo que la discriminación sea inevitable.} « Je ne crois pas que la discrimination soit inévitable. » (subjonctif : \esp{sea})

\begin{propriete}[« Creer que » + indicatif, « no creer que » + subjonctif]
\begin{itemize}
\item Après \esp{creo que / pienso que / me parece que} (opinion \textbf{affirmative}), le verbe suivant est à l'\textbf{indicatif} : \esp{Creo que la tolerancia \textbf{es} importante.}
\item Après \esp{no creo que / no pienso que / no me parece que} (opinion \textbf{négative}, donc doute), le verbe suivant est au \textbf{subjonctif} : \esp{No creo que \textbf{sea} justo.}
\end{itemize}
Le subjonctif présent des verbes fréquents : \esp{ser} $\rightarrow$ \esp{sea} ; \esp{estar} $\rightarrow$ \esp{esté} ; \esp{haber} $\rightarrow$ \esp{haya} ; \esp{tener} $\rightarrow$ \esp{tenga} ; \esp{poder} $\rightarrow$ \esp{pueda} ; verbes en \esp{-ar} $\rightarrow$ \esp{-e} (\esp{hable}) ; verbes en \esp{-er/-ir} $\rightarrow$ \esp{-a} (\esp{viva}, \esp{escriba}).
\end{propriete}

\begin{exemplebox}[Exemples traduits]
\esp{No creo que la discriminación sea inevitable.} « Je ne crois pas que la discrimination soit inévitable. »

\esp{No pienso que los prejuicios desaparezcan fácilmente.} « Je ne pense pas que les préjugés disparaissent facilement. »

\esp{No creo que haya conflictos en nuestra clase.} « Je ne crois pas qu'il y ait des conflits dans notre classe. »

\esp{—¿Crees que la integración es fácil? —No creo que sea fácil, pero es necesaria.} « —Crois-tu que l'intégration est facile ? —Je ne crois pas qu'elle soit facile, mais elle est nécessaire. »
\end{exemplebox}

\begin{attention}[Piège : l'indicatif après « no creo que »]
Erreur typique du francophone : \esp{*No creo que la discriminación \textbf{es} inevitable}. En français, « je ne crois pas qu'elle \emph{est} » est possible à l'oral, mais en espagnol la négation de l'opinion exige \textbf{toujours} le subjonctif : \esp{no creo que \textbf{sea}}. Astuce : dès que vous voyez \esp{no creo que}, demandez-vous : « quel est le subjonctif de ce verbe ? »
\end{attention}

\courssub{Demander l'avis de l'autre}

Un débat est un échange : il faut aussi savoir \textbf{interroger} son partenaire.

\begin{center}
\begin{tabularx}{\linewidth}{|X|X|}
\hline
\rowcolor{popblueL} Español & Français \\
\hline
\esp{¿Qué piensas tú?} & Qu'en penses-tu ? \\
\hline
\esp{¿Qué opinas?} & Quelle est ton opinion ? \\
\hline
\esp{¿Estás de acuerdo?} & Es-tu d'accord ? \\
\hline
\esp{¿Cuál es tu opinión?} & Quelle est ton opinion ? \\
\hline
\esp{¿Qué te parece?} & Qu'est-ce que cela te semble ? \\
\hline
\esp{¿Y tú, qué dices?} & Et toi, qu'en dis-tu ? \\
\hline
\end{tabularx}
\end{center}

Attention à la forme de \esp{parecer} : \esp{¿Qué \textbf{te} parece?} (littéralement « que te semble-t-il ? »), avec un pronom complément devant le verbe, comme \esp{gustar}.

\courssub{L'échelle de l'accord : du « oui » total au « non » total}

Toutes les opinions ne sont pas blanches ou noires. Voici l'échelle complète, de l'accord total au désaccord total :

\begin{popfigure}
\begin{center}
\begin{tikzpicture}[font=\small]
\draw[-{Stealth}, line width=1.2pt, popmuted] (0,0) -- (10.5,0);
\fill[popgreen] (0.6,0) circle (3pt);
\node[above=6pt, align=center, text=popgreen] at (0.6,0) {\esp{totalmente}\\\esp{de acuerdo}};
\node[below=6pt, align=center, text=popmuted] at (0.6,0) {tout à fait\\d'accord};
\fill[popgreen!60!popblue] (3,0) circle (3pt);
\node[above=6pt, align=center, text=popgreen!60!popblue] at (3,0) {\esp{de acuerdo}};
\node[below=6pt, align=center, text=popmuted] at (3,0) {d'accord};
\fill[popgold] (5.25,0) circle (3pt);
\node[above=6pt, align=center, text=popgold] at (5.25,0) {\esp{en parte}};
\node[below=6pt, align=center, text=popmuted] at (5.25,0) {en partie};
\fill[poporange] (7.6,0) circle (3pt);
\node[above=6pt, align=center, text=poporange] at (7.6,0) {\esp{no muy}\\\esp{de acuerdo}};
\node[below=6pt, align=center, text=popmuted] at (7.6,0) {pas vraiment\\d'accord};
\fill[popink] (9.9,0) circle (3pt);
\node[above=6pt, align=center, text=popink] at (9.9,0) {\esp{totalmente}\\\esp{en desacuerdo}};
\node[below=6pt, align=center, text=popmuted] at (9.9,0) {totalement\\en désaccord};
\end{tikzpicture}
\end{center}
\legende{L'échelle de l'accord : cinq degrés pour nuancer son opinion dans un débat.}
\end{popfigure}

Exemples : \esp{Estoy totalmente de acuerdo con la profesora.} « Je suis tout à fait d'accord avec la professeure. » ; \esp{Estoy de acuerdo en parte.} « Je suis d'accord en partie. » ; \esp{Estoy totalmente en desacuerdo con esa idea.} « Je suis totalement en désaccord avec cette idée. » Remarquez : pour le désaccord total, on peut dire \esp{estar en desacuerdo} ou \esp{no estar de acuerdo}.

\courssub{Comparaison français / espagnol}

\begin{center}
\begin{tabularx}{\linewidth}{|X|X|X|}
\hline
\rowcolor{popblueL} Français & Español & Remarque \\
\hline
être d'accord & \esp{estar de acuerdo} & jamais \esp{ser} ! \\
\hline
avoir raison & \esp{tener razón} & \esp{tener}, pas \esp{ser} \\
\hline
je pense la même chose & \esp{pienso lo mismo} & \esp{lo mismo}, invariable \\
\hline
il est vrai que + indicatif & \esp{es verdad que} + indicatif & même construction \\
\hline
je ne crois pas qu'il soit & \esp{no creo que sea} & subjonctif dans les deux langues \\
\hline
au contraire & \esp{al contrario} & \esp{al}, pas « à la » \\
\hline
en partie & \esp{en parte} & pas de \esp{*en parte de} \\
\hline
\end{tabularx}
\end{center}

\begin{methode}[Comment débattre en classe : quatre étapes]
\begin{enumerate}
\item \textbf{Écouter} attentivement l'opinion de l'autre (et noter les mots-clés : \esp{pienso que...}, \esp{en mi opinión...}).
\item \textbf{Réagir} avec une formule adaptée au degré de votre avis : \esp{estoy de acuerdo}, \esp{estoy de acuerdo en parte}, \esp{no estoy de acuerdo}, \esp{no creo que...} + subjonctif.
\item \textbf{Justifier} toujours avec \esp{porque} : une opinion sans argument ne vaut rien dans un débat.
\item \textbf{Conclure poliment} : \esp{entiendo tu opinión}, \esp{gracias por tu punto de vista}, \esp{comparto tu idea final}. Le respect (\esp{respeto}) est la règle d'or du vivre-ensemble, même dans le désaccord.
\end{enumerate}
\end{methode}

\begin{exoresolu}[Réagir et justifier]
Énoncé : Réagissez à chaque opinion avec la formule indiquée, puis justifiez avec \esp{porque}.

1. \esp{La diversidad cultural es una riqueza.} (accord total)

2. \esp{Los conflictos entre culturas son inevitables.} (désaccord + \esp{no creo que})

3. \esp{La escuela debe enseñar la tolerancia.} (accord argumenté)

4. \esp{Las diferencias siempre crean problemas.} (désaccord nuancé : \esp{es cierto que... pero...})

\tcblower
Solution détaillée :

1. \esp{Estoy totalmente de acuerdo contigo porque cada cultura aporta algo único a la sociedad.} « Je suis tout à fait d'accord avec toi parce que chaque culture apporte quelque chose d'unique à la société. » (accord total + \esp{porque} + argument)

2. \esp{No estoy de acuerdo. No creo que los conflictos sean inevitables porque con respeto y educación podemos evitarlos.} « Je ne suis pas d'accord. Je ne crois pas que les conflits soient inévitables parce qu'avec le respect et l'éducation nous pouvons les éviter. » (attention : \esp{sean}, subjonctif de \esp{ser}, et non \esp{*son})

3. \esp{Comparto tu opinión porque la escuela es el lugar donde aprendemos a vivir juntos.} « Je partage ton opinion parce que l'école est le lieu où nous apprenons à vivre ensemble. » (accord argumenté)

4. \esp{Es cierto que las diferencias a veces crean problemas, pero también enriquecen nuestra vida.} « Il est vrai que les différences créent parfois des problèmes, mais elles enrichissent aussi notre vie. » (concession + \esp{pero} + objection : après \esp{es cierto que}, indicatif \esp{crean})
\end{exoresolu}

\begin{attention}[Récapitulatif des pièges des francophones]
\begin{itemize}
\item \esp{*Soy de acuerdo} $\rightarrow$ \esp{\textbf{Estoy} de acuerdo} (état d'opinion = \esp{estar}).
\item \esp{*Tienes la razón} $\rightarrow$ \esp{Tienes razón} (sans article).
\item \esp{*No creo que es justo} $\rightarrow$ \esp{No creo que \textbf{sea} justo} (subjonctif obligatoire).
\item \esp{*Estoy de acuerdo con tú} $\rightarrow$ \esp{Estoy de acuerdo \textbf{contigo}} (forme spéciale après \esp{con}).
\item Ne traduisez pas « je suis d'accord \emph{sur} ce point » par \esp{*de acuerdo sobre} : dites \esp{estoy de acuerdo en este punto} ou \esp{en parte}.
\end{itemize}
\end{attention}

\begin{exercice}[À vous de débattre]
1. Conjuguez : \esp{estar de acuerdo} aux six personnes du présent.

2. Traduisez en espagnol : a) « Je suis d'accord avec toi parce que la tolérance est importante. » b) « Tu as raison : le respect commence à l'école. » c) « Je ne crois pas que les préjugés soient une fatalité. » d) « Il est vrai que le vivre-ensemble est difficile, cependant il est possible. »

3. Réagissez à cette opinion avec un désaccord nuancé : \esp{En el Camerún, la convivencia es imposible porque hay demasiadas lenguas diferentes.}
\end{exercice}

\begin{corrige}[Exercice : À vous de débattre]
1. \esp{estoy, estás, está, estamos, estáis, están de acuerdo}.

2. a) \esp{Estoy de acuerdo contigo porque la tolerancia es importante.} b) \esp{Tienes razón: el respeto empieza en la escuela.} c) \esp{No creo que los prejuicios sean una fatalidad.} (subjonctif \esp{sean}) d) \esp{Es cierto que la convivencia es difícil, sin embargo, es posible.}

3. Réponse possible : \esp{No estoy de acuerdo. Es verdad que hay muchas lenguas en el Camerún, pero no creo que la convivencia sea imposible: al contrario, esta diversidad es una riqueza porque nos enseña a respetar a los demás.}
\end{corrige}

\begin{aretenir}[L'essentiel de la section]
\begin{itemize}
\item « Être d'accord » = \esp{\textbf{estar} de acuerdo} (jamais \esp{ser}) ; « avoir raison » = \esp{tener razón}.
\item Accord : \esp{estoy de acuerdo, tienes razón, es verdad, claro, exactamente, pienso lo mismo, comparto tu opinión}.
\item Désaccord : \esp{no estoy de acuerdo, al contrario, puede ser pero..., es cierto que... sin embargo...}.
\item \esp{No creo que} + \textbf{subjonctif} : \esp{no creo que sea justo}.
\item Demander l'avis : \esp{¿qué piensas tú? ¿estás de acuerdo? ¿cuál es tu opinión?}.
\item Dans un débat : écouter, réagir, justifier avec \esp{porque}, conclure poliment — c'est cela, la \esp{convivencia}.
\end{itemize}
\end{aretenir}

\coursec{Méthode du commentaire de texte et commentaire modèle}

Dans la section précédente, tu as appris à exprimer ton accord ou ton désaccord avec des formules comme \esp{estoy de acuerdo} ou \esp{no estoy de acuerdo}. C'est exactement cette compétence que tu vas maintenant mobiliser dans l'exercice roi de la classe de Première A : le \cle{commentaire de texte} (\esp{comentario de texto}). En effet, commenter un texte, c'est d'abord le comprendre, puis dire si l'on partage ou non les idées de l'auteur. Voyons comment procéder, étape par étape, puis lisons ensemble un commentaire modèle rédigé sur notre texte support \emph{Vivir juntos en la diferencia}.

\courssub{Qu'est-ce qu'un commentaire de texte ?}

\begin{definition}[Le comentario de texto]
Le \cle{commentaire de texte} est un exercice écrit dans lequel l'élève \textbf{présente} un texte, \textbf{analyse} ses idées et ses procédés, puis donne son \textbf{opinion personnelle}. En espagnol : \esp{un comentario de texto es un ejercicio en el que presentamos, analizamos y valoramos un texto}. Il ne s'agit ni de résumer, ni de traduire : il s'agit d'\emph{expliquer} et de \emph{juger}.
\end{definition}

Avant de commenter, relisons le texte support de la leçon.

\begin{textebox}[Vivir juntos en la diferencia (texte support)]
En un país como el Camerún, donde conviven más de doscientos grupos étnicos, la convivencia es un desafío diario. Sin embargo, esta diversidad no es un problema : es una riqueza.

La tolerancia no significa aceptar todo sin reflexionar, sino respetar al otro aunque piense de manera diferente. Los prejuicios nacen de la ignorancia : cuando no conocemos al vecino, inventamos ideas falsas sobre él. En cambio, el diálogo y la educación construyen puentes entre las culturas.

En la escuela, los alumnos aprenden juntos, juegan juntos y celebran juntos. El aula es un laboratorio de la convivencia : allí se aprende el respeto, la escucha y la solidaridad. Un alumno bamiléké, un alumno beti y un alumno fulbé pueden ser amigos si descubren lo que tienen en común.

La identidad de cada uno es importante, pero no debe convertirse en una muralla. Ser camerunés no significa renunciar a su cultura : significa compartirla con los demás. Como dice un proverbio africano, una sola mano no puede aplaudir.

En conclusión, vivir juntos en la diferencia exige esfuerzo, pero es el único camino hacia la paz. Debemos educar en el respeto desde la infancia, porque el futuro de nuestra sociedad depende de nuestra capacidad de convivir.
\end{textebox}

\textbf{Glossaire :} \esp{convivir} = vivre ensemble ; \esp{desafío} = défi ; \esp{sin embargo} = cependant ; \esp{prejuicio} = préjugé ; \esp{ignorancia} = ignorance ; \esp{puente} = pont ; \esp{aula} = salle de classe ; \esp{muralla} = muraille ; \esp{renunciar} = renoncer ; \esp{aplaudir} = applaudir ; \esp{esfuerzo} = effort.

\begin{attention}[Deux subtilités du texte à remarquer]
\begin{itemize}
\item \esp{respetar al otro aunque piense de manera diferente} : après \esp{aunque}, on emploie le \textbf{subjonctif} (\esp{piense}) quand l'idée est hypothétique ou générale (« même s'il pense autrement »). Avec l'indicatif (\esp{aunque piensa}), on affirme un fait réel.
\item \esp{el aula} : ce mot est \textbf{féminin} malgré l'article \esp{el} (on dit \esp{el aula} pour éviter la rencontre de deux « a », comme \esp{el agua}), mais l'accord reste féminin : \esp{el aula es grande}.
\end{itemize}
\end{attention}

\courssub{Les trois étapes du commentaire}

\begin{methode}[Rédiger un commentaire de texte en trois étapes]
\begin{enumerate}
\item \textbf{Introduction} (\esp{introducción}) : présente le texte (titre, thème, genre) et annonce la problématique (\esp{problemática}), c'est-à-dire la question centrale du texte. Utilise : \esp{El texto trata de...} / \esp{El autor defiende la idea de que...}.
\item \textbf{Développement} (\esp{desarrollo}) : analyse la thèse de l'auteur, ses arguments, le champ lexical (\esp{campo léxico}) et les procédés (exemples, comparaisons, proverbes). Cite au moins deux passages du texte entre guillemets.
\item \textbf{Conclusion} (\esp{conclusión}) : fais le bilan (\esp{balance}) et donne ton opinion personnelle avec \esp{estoy de acuerdo / no estoy de acuerdo}, puis ouvre sur une question : \esp{¿Y tú, qué piensas?}
\end{enumerate}
\end{methode}

\begin{popfigure}
\begin{center}
\begin{tikzpicture}[node distance=0.9cm]
\node[draw=popblue, fill=popblueL, rounded corners, text width=9.5cm, align=center, font=\bfseries] (intro) {1. INTRODUCCIÓN\\ \normalfont tema + problemática};
\node[draw=poporange, fill=poporangeL, rounded corners, text width=9.5cm, align=center, font=\bfseries, below=of intro] (des) {2. DESARROLLO\\ \normalfont tesis + argumentos + campo léxico + citas};
\node[draw=popgreen, fill=popgreenL, rounded corners, text width=9.5cm, align=center, font=\bfseries, below=of des] (conc) {3. CONCLUSIÓN\\ \normalfont balance + opinión personal};
\draw[-{Stealth}, very thick, popblue] (intro) -- (des);
\draw[-{Stealth}, very thick, poporange] (des) -- (conc);
\end{tikzpicture}
\end{center}
\legende{La structure en trois blocs du commentaire de texte}
\end{popfigure}

\courssub{Étape 1 : l'introduction}

L'introduction est courte : deux ou trois phrases suffisent. Elle doit contenir le \textbf{thème} du texte et la \textbf{problématique}.

\begin{exemplebox}[Modèles d'introduction]
\esp{El texto trata de la tolerancia y la convivencia en una sociedad multicultural. El autor defiende la idea de que la diversidad es una riqueza. La problemática es la siguiente : ¿cómo vivir juntos en la diferencia?}

« Le texte traite de la tolérance et du vivre-ensemble dans une société multiculturelle. L'auteur défend l'idée que la diversité est une richesse. La problématique est la suivante : comment vivre ensemble dans la différence ? »
\end{exemplebox}

\begin{attention}[Ne confonds pas thème et problématique]
Le \textbf{thème} est le sujet général (\esp{la convivencia}) ; la \textbf{problématique} est la question précise que le texte pose (\esp{¿cómo convivir respetando las diferencias?}). Une bonne problématique commence souvent par \esp{¿cómo...?}, \esp{¿por qué...?} ou \esp{¿es posible...?} N'oublie pas le point d'interrogation ouvrant \esp{¿} en espagnol, et l'accent sur les mots interrogatifs : \esp{cómo}, \esp{por qué}, \esp{qué}.
\end{attention}

\courssub{Étape 2 : le développement}

Le développement est la partie la plus longue (deux paragraphes environ). Tu y analyses quatre éléments :

\begin{itemize}
\item la \textbf{thèse} (\esp{tesis}) : l'idée principale défendue par l'auteur ;
\item les \textbf{arguments} (\esp{argumentos}) : les raisons et exemples qui soutiennent la thèse ;
\item le \textbf{champ lexical} (\esp{campo léxico}) : les familles de mots du thème (\esp{tolerancia, respeto, diálogo, prejuicio...}) ;
\item les \textbf{procédés} (\esp{procedimientos}) : comparaisons, métaphores, proverbes, oppositions (\esp{en cambio}, \esp{sin embargo}).
\end{itemize}

\begin{exemplebox}[Phrases-outils du développement]
\esp{Según el autor, los prejuicios nacen de la ignorancia.} = « Selon l'auteur, les préjugés naissent de l'ignorance. »

\esp{El autor afirma que el aula es un laboratorio de la convivencia.} = « L'auteur affirme que la salle de classe est un laboratoire du vivre-ensemble. »

\esp{Para el autor, la identidad no debe convertirse en una muralla.} = « Pour l'auteur, l'identité ne doit pas se transformer en muraille. »
\end{exemplebox}

\begin{propriete}[Les verbes pour rapporter les idées de l'auteur]
Pour présenter la pensée de l'auteur sans la répéter mot à mot, utilise ces verbes à l'indicatif présent : \esp{el autor afirma que} (affirme), \esp{defiende que} (défend), \esp{explica que} (explique), \esp{propone} (propose), \esp{advierte que} (avertit), \esp{concluye que} (conclut). Attention à la diphtongaison : \esp{defender} donne \esp{defiende}, \esp{advertir} donne \esp{advierte}.
\end{propriete}

\courssub{Étape 3 : la conclusion}

La conclusion fait le bilan et donne ton avis. C'est ici que tu réutilises les formules d'accord et de désaccord de la section précédente.

\begin{exemplebox}[Modèles de conclusion]
\esp{En conclusión, debemos educar en el respeto desde la infancia.} = « En conclusion, nous devons éduquer au respect dès l'enfance. »

\esp{Estoy de acuerdo con el autor porque la diversidad enriquece nuestra sociedad.} = « Je suis d'accord avec l'auteur parce que la diversité enrichit notre société. »

\esp{¿Y tú, qué piensas de la convivencia en tu barrio?} = « Et toi, que penses-tu du vivre-ensemble dans ton quartier ? »
\end{exemplebox}

\begin{attention}[L'accord de « estar de acuerdo »]
La formule \esp{estar de acuerdo} est \textbf{invariable} : on dit \esp{estoy de acuerdo}, \esp{ella está de acuerdo}, jamais \esp{estoy de acuerda}. En revanche, on accorde le sujet du verbe : \esp{nosotros estamos de acuerdo}. Pour nuancer : \esp{estoy totalmente de acuerdo} (tout à fait d'accord), \esp{estoy parcialmente de acuerdo} (partiellement d'accord), \esp{no estoy de acuerdo en absoluto} (pas d'accord du tout).
\end{attention}

\begin{center}
\begin{tabularx}{\linewidth}{|X|X|X|}
\hline
\rowcolor{popblueL} Étape & Phrase-outil (español) & Traduction \\
\hline
Introduction & El texto trata de... & Le texte traite de... \\
\hline
Introduction & El autor defiende la idea de que... & L'auteur défend l'idée que... \\
\hline
Développement & Según el autor,... & Selon l'auteur,... \\
\hline
Développement & El autor afirma que... & L'auteur affirme que... \\
\hline
Développement & Por ejemplo / en cambio & Par exemple / en revanche \\
\hline
Conclusion & En conclusión,... & En conclusion,... \\
\hline
Conclusion & Estoy de acuerdo porque... & Je suis d'accord parce que... \\
\hline
\end{tabularx}
\end{center}

\begin{attention}[Le piège n° 1 : traduire au lieu de commenter]
Beaucoup d'élèves francophones se contentent de \textbf{traduire} ou de \textbf{recopier} le texte. Erreur grave ! Le commentaire \textbf{analyse} et \textbf{donne un avis} ; il ne répète pas. Interdit : recopier trois phrases du texte sans les expliquer. Obligatoire : après chaque citation, ajoute ton explication avec \esp{esto significa que...} (« cela signifie que... ») ou \esp{el autor quiere decir que...} (« l'auteur veut dire que... »).
\end{attention}

\courssub{Commentaire modèle rédigé}

\begin{textebox}[Comentario modelo : Vivir juntos en la diferencia]
\textbf{Introducción.} El texto trata de la tolerancia y la convivencia en una sociedad multicultural como el Camerún. El autor defiende la idea de que la diversidad es una riqueza y no un problema. La problemática es la siguiente : ¿cómo vivir juntos respetando las diferencias?

\textbf{Desarrollo.} En primer lugar, el autor presenta su tesis : la convivencia es un desafío diario, pero la diversidad es una riqueza. Según el autor, los prejuicios nacen de la ignorancia : cuando no conocemos al otro, inventamos ideas falsas. Para combatirlos, propone el diálogo y la educación, que « construyen puentes entre las culturas ». El campo léxico del texto es muy rico : tolerancia, respeto, solidaridad, identidad, prejuicio.

En segundo lugar, el autor utiliza ejemplos concretos y procedimientos eficaces. Afirma que « el aula es un laboratorio de la convivencia » : esta metáfora muestra que la escuela enseña el respeto. También cita un proverbio africano : « una sola mano no puede aplaudir », que significa que necesitamos a los demás. La imagen de la « muralla » advierte contra el encierro en la propia identidad.

\textbf{Conclusión.} En conclusión, el texto nos invita a educar en el respeto desde la infancia. Estoy de acuerdo con el autor porque en mi ciudad, Yaoundé, personas de culturas diferentes viven y trabajan juntas cada día. ¿Y tú, qué piensas de la convivencia en tu escuela?
\end{textebox}

\textbf{Glossaire :} \esp{en primer lugar} = d'abord ; \esp{combatir} = combattre (\esp{combatirlos} = les combattre) ; \esp{advertir} = avertir ; \esp{encierro} = enfermement ; \esp{juntos} = ensemble ; \esp{procedimientos} = procédés.

\begin{methode}[Comment lire le commentaire modèle]
\begin{enumerate}
\item Repère les \textbf{trois blocs} : introduction, développement (deux paragraphes), conclusion.
\item Souligne les \textbf{phrases-outils} : \esp{el texto trata de...}, \esp{según el autor...}, \esp{en conclusión...}, \esp{estoy de acuerdo porque...}.
\item Entoure les \textbf{deux citations} entre guillemets et vérifie que chacune est suivie d'une explication (\esp{esta metáfora muestra que...}, \esp{que significa que...}).
\item Observe la \textbf{conclusion} : bilan, opinion personnelle justifiée par un exemple camerounais, puis question d'ouverture au lecteur.
\end{enumerate}
\end{methode}

\courssub{Les critères d'évaluation au Bac}

\begin{propriete}[Ce que l'examinateur attend]
Au Bac A, le commentaire fait environ \textbf{10 à 15 lignes} et doit \textbf{citer au moins deux passages du texte} entre guillemets. Les quatre critères d'évaluation sont :
\begin{enumerate}
\item la \textbf{fidélité au texte} : tu analyses les idées de l'auteur, pas les tiennes d'abord ;
\item la \textbf{richesse du lexique} : tu utilises le vocabulaire du thème (\esp{tolerancia, convivencia, prejuicio, respeto}) ;
\item la \textbf{correction grammaticale} : accords, conjugaisons, accents ;
\item l'\textbf{opinion personnelle} : tu termines par ton avis argumenté.
\end{enumerate}
\end{propriete}

\begin{savaistu}[Le commentaire au Bac A]
Les sujets du Bac A proposent souvent un \textbf{texte de société} (tolérance, éducation, environnement, jeunesse) suivi de \textbf{questions de compréhension} et d'une \textbf{expression personnelle} de 5 à 8 lignes. S'entraîner au commentaire, c'est donc préparer directement l'épreuve : les phrases-outils que tu apprends ici (\esp{el texto trata de...}, \esp{según el autor...}, \esp{en conclusión...}) peuvent être réutilisées dans presque tous les sujets.
\end{savaistu}

\begin{exoresolu}[Identifier les éléments d'un commentaire]
Énoncé : dans la phrase suivante, extraite d'un commentaire, identifie (a) la présentation du thème, (b) la thèse de l'auteur, (c) l'opinion personnelle de l'élève : \esp{El texto trata de la integración de los jóvenes ; el autor defiende la idea de que la escuela une a los alumnos ; en mi opinión, tiene razón porque en mi clase hay amigos de todas las regiones.}
\tcblower
Solution : (a) \esp{El texto trata de la integración de los jóvenes} : c'est la présentation du thème (formule \esp{trata de}). (b) \esp{el autor defiende la idea de que la escuela une a los alumnos} : c'est la thèse (formule \esp{defiende la idea de que}). (c) \esp{en mi opinión, tiene razón porque...} : c'est l'opinion personnelle, justifiée par un exemple concret (\esp{en mi clase hay amigos de todas las regiones}). Remarque : l'élève donne une raison avec \esp{porque}, ce qui est indispensable pour obtenir tous les points du critère « opinion personnelle ». Note aussi l'expression \esp{tener razón} (« avoir raison ») : \esp{el autor tiene razón}, sans article.
\end{exoresolu}

\begin{exercice}[À ton tour]
Rédige l'introduction d'un commentaire (2 ou 3 phrases) sur un texte dont le thème est \esp{el respeto entre los jóvenes} et dont la thèse est : \esp{el respeto se aprende en la familia y en la escuela}. Utilise obligatoirement \esp{el texto trata de...} et \esp{el autor defiende la idea de que...}.
\end{exercice}

\begin{corrige}[Exercice]
Modèle de réponse : \esp{El texto trata del respeto entre los jóvenes. El autor defiende la idea de que el respeto se aprende en la familia y en la escuela. La problemática es la siguiente : ¿dónde y cómo aprenden los jóvenes a respetar a los demás?} — Vérifie la présence des deux formules obligatoires, de la problématique sous forme de question (avec \esp{¿} ouvrant), et des accents (\esp{jóvenes}, \esp{problemática}, \esp{cómo}, \esp{dónde}). Remarque la contraction \esp{trata del} (\esp{de} + \esp{el} = \esp{del}).
\end{corrige}

\begin{aretenir}[L'essentiel de la méthode]
Un commentaire de texte se construit en trois blocs : une \textbf{introduction} (thème + problématique), un \textbf{développement} (thèse + arguments + champ lexical + deux citations expliquées) et une \textbf{conclusion} (bilan + opinion personnelle). Garde en mémoire les phrases-outils : \esp{el texto trata de...}, \esp{según el autor...}, \esp{en conclusión...}. Et surtout : commenter, c'est analyser et donner son avis, jamais traduire ou recopier le texte.
\end{aretenir}

\coursec{Production guidée : le vivre-ensemble au Cameroun}

Dans la section précédente, nous avons appris à \emph{commenter} un texte sur la tolérance : observer, comprendre, analyser. Il est temps maintenant de passer de lecteur à \emph{auteur}. Le savoir-faire visé ici est double : exprimer son accord ou son désaccord à l'oral, et rédiger quelques phrases personnelles sur le vivre-ensemble au Cameroun. C'est exactement le type de tâche demandée à l'évaluation harmonisée et, plus tard, au Baccalauréat : une courte production écrite, guidée, personnelle et correcte.

\courssub{Brainstorming collectif : la convivencia au Cameroun}

Avant d'écrire, on réfléchit à voix haute. Le professeur écrit au tableau la question : \esp{¿Dónde vemos la convivencia en Camerún?} (« Où voyons-nous le vivre-ensemble au Cameroun ? »). Chaque élève propose une idée, en espagnol simple, que la classe valide ou discute avec \esp{Estoy de acuerdo} ou \esp{No estoy de acuerdo}.

\begin{experience}[Brainstorming en classe entière]
\begin{enumerate}
\item Le professeur écrit au tableau quatre domaines : \esp{la escuela}, \esp{el barrio}, \esp{las fiestas}, \esp{la familia}.
\item Par groupes de quatre, les élèves notent en deux minutes un exemple de \esp{convivencia} pour chaque domaine.
\item Chaque groupe présente une idée à la classe. Les autres réagissent : \esp{Estoy de acuerdo porque...} / \esp{No estoy de acuerdo porque...}.
\item Le professeur garde au tableau les meilleures idées : elles serviront de matière à la rédaction.
\end{enumerate}
\end{experience}

Exemples d'idées attendues (le professeur les reformule si besoin) :

\begin{itemize}
\item \esp{En mi escuela hay alumnos de muchas regiones y hablamos juntos.} « Dans mon école, il y a des élèves de nombreuses régions et nous parlons ensemble. »
\item \esp{En mi barrio, los vecinos se ayudan durante las fiestas.} « Dans mon quartier, les voisins s'entraident pendant les fêtes. »
\item \esp{Hay matrimonios entre personas de culturas diferentes.} « Il y a des mariages entre personnes de cultures différentes. »
\item \esp{En el mercado de Mokolo, los comerciantes hablan varias lenguas con los clientes.} « Au marché Mokolo, les commerçants parlent plusieurs langues avec les clients. »
\item \esp{Los cristianos y los musulmanes celebran sus fiestas en paz.} « Les chrétiens et les musulmans célèbrent leurs fêtes en paix. »
\end{itemize}

\begin{savaistu}[Le Cameroun, « l'Afrique en miniature »]
On dit souvent que le Cameroun es \esp{« África en miniatura »} : environ 250 groupes ethniques, deux langues officielles (le français et l'anglais), des langues nationales comme l'ewondo, le duala ou le fulfuldé, et plusieurs religions qui cohabitent. Cette diversité fait du thème de la \esp{convivencia} un sujet très naturel pour une production écrite camerounaise : chaque élève a un exemple concret sous les yeux, dans sa classe, son quartier ou sa famille.
\end{savaistu}

\courssub{Le plan guidé en quatre étapes}

Une bonne petite production n'est pas une suite de phrases au hasard : elle suit un \emph{plan}. Pour cette leçon, le plan comporte quatre étapes, faciles à mémoriser.

\begin{popfigure}
\begin{center}
\begin{tikzpicture}[node distance=0.55cm]
\node[draw=popblue, fill=popblueL, rounded corners, thick, minimum width=3.1cm, minimum height=1.1cm, align=center] (d) {\textbf{1. Diversidad}\\ \footnotesize presentar};
\node[draw=poporange, fill=poporangeL, rounded corners, thick, minimum width=3.1cm, minimum height=1.1cm, align=center, right=of d] (e) {\textbf{2. Ejemplo}\\ \footnotesize de tolerancia};
\node[draw=popgreen, fill=popgreenL, rounded corners, thick, minimum width=3.1cm, minimum height=1.1cm, align=center, right=of e] (o) {\textbf{3. Opinión}\\ \footnotesize personal};
\node[draw=poppurple, fill=poppurpleL, rounded corners, thick, minimum width=3.1cm, minimum height=1.1cm, align=center, right=of o] (c) {\textbf{4. Conclusión}\\ \footnotesize final};
\draw[-{Stealth}, very thick, popdark] (d) -- (e);
\draw[-{Stealth}, very thick, popdark] (e) -- (o);
\draw[-{Stealth}, very thick, popdark] (o) -- (c);
\end{tikzpicture}
\end{center}
\legende{Frise du plan de production : Diversidad, Ejemplo, Opinión, Conclusión.}
\end{popfigure}

\begin{enumerate}
\item \textbf{Présenter la diversité camerounaise} : une ou deux phrases avec \esp{hay}, \esp{en Camerún hay...}, \esp{mi país tiene...}.
\item \textbf{Donner un exemple de tolérance} : un fait concret (école, quartier, fête, mariage mixte) tiré du brainstorming.
\item \textbf{Exprimer son opinion} : avec \esp{en mi opinión}, \esp{pienso que}, \esp{estoy de acuerdo con...}, justifiée par \esp{porque}.
\item \textbf{Conclure} : une phrase d'ouverture ou d'appel, avec \esp{por eso} ou \esp{en conclusión}.
\end{enumerate}

\courssub{Les connecteurs obligatoires du paragraphe d'opinion}

Observons d'abord deux phrases modèles, que tu dois pouvoir réutiliser telles quelles ou presque.

\begin{exemplebox}[Deux phrases clés à retenir]
\esp{En mi opinión, la diversidad de Camerún es una riqueza.} « À mon avis, la diversité du Cameroun est une richesse. »

\esp{Debemos convivir en paz y respetar a los demás.} « Nous devons vivre ensemble en paix et respecter les autres. »
\end{exemplebox}

Remarque la structure de la deuxième phrase : \esp{debemos} + deux infinitifs (\esp{convivir}, \esp{respetar}). C'est un modèle très rentable : un seul verbe conjugué, puis des infinitifs, donc moins de risques d'erreur.

\begin{methode}[Rédiger un paragraphe d'opinion en cinq connecteurs]
\begin{enumerate}
\item \textbf{Ouvrir} avec \esp{primero} pour annoncer le premier argument ou la présentation : \esp{Primero, Camerún es un país muy diverso.}
\item \textbf{Ajouter} avec \esp{además} un exemple ou un argument supplémentaire : \esp{Además, en mi barrio los vecinos se ayudan.}
\item \textbf{Nouancer} avec \esp{sin embargo} si tu veux mentionner une difficulté (préjugés, conflits) : \esp{Sin embargo, a veces hay prejuicios.}
\item \textbf{Tirer la conséquence} avec \esp{por eso} : \esp{Por eso, debemos respetar a los demás.}
\item \textbf{Conclure} avec \esp{en conclusión} : \esp{En conclusión, la tolerancia es necesaria para vivir en paz.}
\end{enumerate}
Ces cinq connecteurs sont \emph{obligatoires} dans le paragraphe d'opinion complet ; pour la production courte de 5 à 8 phrases, on en exige au minimum deux.
\end{methode}

\begin{center}
\begin{tabularx}{\linewidth}{|X|X|X|}
\hline
\rowcolor{popblueL} \textbf{Connecteur} & \textbf{Fonction} & \textbf{Français} \\
\hline
\esp{primero} & annoncer le premier point & d'abord, premièrement \\
\hline
\esp{además} & ajouter une idée & de plus, en outre \\
\hline
\esp{sin embargo} & opposer, nuancer & cependant, pourtant \\
\hline
\esp{por eso} & exprimer la conséquence & c'est pourquoi, donc \\
\hline
\esp{en conclusión} & conclure & en conclusion \\
\hline
\esp{en mi opinión} & donner son avis & à mon avis \\
\hline
\esp{porque} & justifier & parce que \\
\hline
\end{tabularx}
\end{center}

\begin{attention}[\esp{Los demás} est invariable]
Pour dire « les autres » en général, l'espagnol utilise \esp{los demás}, qui ne change \emph{jamais} : \esp{Hay que respetar a los demás.} « Il faut respecter les autres. » Ne dis pas \esp{*los otros} dans ce sens général : \esp{los otros} ne s'emploie que pour « les autres » parmi un groupe précis déjà nommé (\esp{los otros alumnos de mi clase} « les autres élèves de ma classe »). Autre piège : \esp{demás} prend toujours l'accent, et on écrit \esp{respetar \textbf{a} los demás} avec la préposition \esp{a} (le « a » personnel devant une personne), que le français ne traduit pas.
\end{attention}

\begin{propriete}[\esp{Estar de acuerdo} : un accord figé]
Dans l'expression \esp{estar de acuerdo}, le mot \esp{acuerdo} est \emph{invariable} : on dit \esp{Estoy de acuerdo}, \esp{Estamos de acuerdo}, jamais \esp{*estoy de acuerda}. Seul le verbe \esp{estar} se conjugue : \esp{¿Estás de acuerdo conmigo?} « Es-tu d'accord avec moi ? » Pour renforcer : \esp{Estoy totalmente de acuerdo.} « Je suis tout à fait d'accord. » Pour nuancer : \esp{Estoy de acuerdo en parte.} « Je suis d'accord en partie. »
\end{propriete}

\courssub{Un texte modèle avant de rédiger}

Avant de te lancer, lis ce court texte original : c'est exactement le genre de production attendue, avec le plan en quatre étapes, le lexique de la leçon et les connecteurs.

\begin{textebox}[La convivencia en mi país]
Primero, Camerún es un país muy diverso: hay muchas etnias, lenguas y religiones diferentes. En mi escuela de Yaoundé, por ejemplo, hay alumnos de todas las regiones del país y hablamos juntos en francés, en inglés y en nuestras lenguas.

Además, en mi barrio la convivencia es una realidad: durante las fiestas, los vecinos cristianos y musulmanes se visitan y comparten la comida. También hay matrimonios entre personas de culturas diferentes, y eso es una riqueza.

Sin embargo, a veces existen prejuicios y discriminación entre las personas. En mi opinión, la diversidad de Camerún es una riqueza y no un problema, porque cada cultura aporta algo importante a la sociedad. Estoy de acuerdo con la idea de que la tolerancia es necesaria para la paz.

Por eso, debemos convivir en paz y respetar a los demás. En conclusión, la identidad de cada persona es valiosa, pero la convivencia es el valor más importante de nuestro país.
\end{textebox}

\textbf{Glossaire} : \esp{etnia} « ethnie » ; \esp{se visitan} « se rendent visite » (verbe pronominal réciproque) ; \esp{comparten} « partagent » (\esp{compartir}) ; \esp{aporta} « apporte » (\esp{aportar}) ; \esp{valiosa} « précieuse, de valeur » ; \esp{la paz} « la paix ».

\textbf{Questions de repérage} (réponds en espagnol, en une phrase) :
\begin{enumerate}
\item ¿Qué ejemplo de convivencia da el autor en su barrio ?
\item ¿Cuál es la opinión del autor sobre la diversidad ?
\item ¿Qué conectores usa el autor para concluir ?
\end{enumerate}

\begin{corrige}[Questions de repérage]
1. \esp{En su barrio, los vecinos cristianos y musulmanes se visitan y comparten la comida durante las fiestas.} — 2. \esp{En su opinión, la diversidad de Camerún es una riqueza y no un problema.} — 3. \esp{Usa « por eso » y « en conclusión ».}
\end{corrige}

\courssub{La tâche de rédaction}

\begin{exercice}[Rédaction guidée : 5 à 8 phrases]
Rédige un court texte en espagnol sur le vivre-ensemble au Cameroun, en suivant le plan en quatre étapes (diversidad, ejemplo, opinión, conclusión).

Contraintes obligatoires :
\begin{itemize}
\item entre 5 et 8 phrases ;
\item au moins 5 mots du lexique de la leçon : \esp{identidad}, \esp{tolerancia}, \esp{respeto}, \esp{convivencia}, \esp{prejuicio}, \esp{diferencia}, \esp{integración}, \esp{discriminación} ;
\item au moins 1 expression d'accord ou de désaccord : \esp{estoy de acuerdo (con...)} ou \esp{no estoy de acuerdo (con...)} ;
\item au moins 2 connecteurs parmi : \esp{primero}, \esp{además}, \esp{sin embargo}, \esp{por eso}, \esp{en conclusión}.
\end{itemize}
\end{exercice}

\begin{exoresolu}[Production corrigée d'un élève]
Énoncé : voici la production d'un élève ; repère si toutes les contraintes sont respectées, puis corrige les deux erreurs.

\esp{Primero, Camerún es un país diverso. En mi escuela hay alumnos de muchas culturas. Estoy de acuerdo con la tolerancia. Debemos respetar a los otros. En conclusión, la convivencia son importante.}

\tcblower

\textbf{Vérification des contraintes} : 5 phrases (correct) ; lexique : \esp{tolerancia}, \esp{convivencia} (seulement 2 mots, il en faut 5 : insuffisant) ; expression d'accord : \esp{estoy de acuerdo con la tolerancia} (correct) ; connecteurs : \esp{primero}, \esp{en conclusión} (correct).

\textbf{Erreur 1} : \esp{*los otros} doit devenir \esp{los demás} (« les autres » en général est invariable).

\textbf{Erreur 2} : \esp{*la convivencia son importante} : le sujet \esp{la convivencia} est singulier, donc \esp{la convivencia \textbf{es} importante}.

\textbf{Version améliorée} (avec 5 mots du lexique) : \esp{Primero, Camerún es un país diverso, con muchas identidades diferentes. En mi escuela hay alumnos de muchas culturas y la integración es buena. Estoy de acuerdo con la tolerancia y el respeto entre todos. Debemos respetar a los demás y luchar contra los prejuicios. En conclusión, la convivencia es importante para nuestro país.}
\end{exoresolu}

\courssub{Relecture en binôme et correction collective}

On ne rend jamais une production sans la relire. En binôme, chaque élève lit le texte de son voisin et coche une grille simple.

\begin{methode}[Grille de relecture en binôme]
\begin{enumerate}
\item \textbf{Accords} : le nom et l'adjectif s'accordent-ils ? \esp{un país divers\textbf{o}}, \esp{las culturas diferente\textbf{s}}.
\item \textbf{Verbes} : le verbe est-il conjugué à la bonne personne ? \esp{hay} (il y a) est invariable ; \esp{debemos} = « nous devons ».
\item \textbf{Lexique} : y a-t-il au moins 5 mots de la leçon ? Souligne-les.
\item \textbf{Connecteurs et opinion} : y a-t-il 2 connecteurs et 1 expression d'accord ou de désaccord ? Encadre-les.
\item \textbf{Orthographe} : accents (\esp{opinión}, \esp{además}, \esp{demás}), \esp{ñ}, point d'exclamation ouvrant \esp{¡} si besoin.
\end{enumerate}
Ensuite, le professeur choisit une production anonyme, l'écrit au tableau, et la classe la corrige collectivement : chaque erreur est repérée par un élève, expliquée, puis corrigée.
\end{methode}

\begin{propriete}[Les critères d'une production réussie au Bac]
Une production écrite réussie au Baccalauréat contient au minimum :
\begin{itemize}
\item \textbf{3 mots du lexique} de la leçon (ici : \esp{tolerancia}, \esp{convivencia}, \esp{respeto}, etc.) ;
\item \textbf{2 connecteurs} logiques (\esp{además}, \esp{por eso}, \esp{sin embargo}...) ;
\item \textbf{1 opinion personnelle justifiée} : \esp{en mi opinión... porque...}.
\end{itemize}
Sans opinion justifiée, la copie reste descriptive et perd des points ; sans connecteurs, les phrases sont juxtaposées et le texte est jugé pauvre.
\end{propriete}

\begin{aretenir}[Complément utile au Bac : le paragraphe d'opinion en quatre mouvements]
Pour le Bac, apprends par cœur le squelette suivant, que tu remplis avec tes idées :

\esp{En mi opinión, ... porque ... Además, ... Por eso, ...}

Exemple complet : \esp{En mi opinión, la diversidad de Camerún es una riqueza porque cada cultura aporta algo importante. Además, la convivencia hace la sociedad más fuerte. Por eso, debemos respetar a los demás y convivir en paz.} « À mon avis, la diversité du Cameroun est une richesse parce que chaque culture apporte quelque chose d'important. De plus, le vivre-ensemble rend la société plus forte. C'est pourquoi nous devons respecter les autres et vivre ensemble en paix. »
\end{aretenir}

\begin{exercice}[Pour aller plus loin]
Transforme ce squelette en paragraphe complet : \esp{En mi opinión, la tolerancia es importante porque ... Además, ... Sin embargo, ... Por eso, ...} Utilise au moins deux mots du lexique de la leçon et le mot \esp{los demás}.
\end{exercice}

\begin{corrige}[Pour aller plus loin]
Proposition : \esp{En mi opinión, la tolerancia es importante porque sin ella no hay convivencia posible. Además, el respeto a los demás hace la vida en comunidad más fácil. Sin embargo, todavía existen prejuicios en nuestra sociedad. Por eso, debemos combatir la discriminación y convivir en paz.} — Lexique utilisé : \esp{tolerancia}, \esp{convivencia}, \esp{respeto}, \esp{prejuicios}, \esp{discriminación} ; \esp{los demás} correctement employé ; quatre connecteurs.
\end{corrige}

Tu sais désormais passer du commentaire à la production : observer la diversité qui t'entoure, la formuler en espagnol avec un plan clair, des connecteurs et une opinion justifiée. C'est cette compétence — dire ce que l'on pense, correctement et avec des mots précis — qui fera la différence à l'évaluation comme au Baccalauréat.

\newpage

\coursec{Activité d'intégration}

\courssub{Objectifs de l'activité}

À l'issue de cette activité d'intégration, l'élève sera capable de :
\begin{itemize}
\item réutiliser dans une communication orale authentique le lexique du vivre-ensemble : \esp{identidad}, \esp{tolerancia}, \esp{respeto}, \esp{convivencia}, \esp{prejuicio}, \esp{diferencia}, \esp{integración}, \esp{discriminación} ;
\item exprimer clairement son accord ou son désaccord à l'oral avec les structures étudiées (\esp{estoy de acuerdo porque...}, \esp{no estoy de acuerdo, sin embargo...}) ;
\item écouter les arguments des camarades, les comprendre et y réagir de manière respectueuse ;
\item rédiger un paragraphe court et cohérent (6 à 8 lignes) sur le vivre-ensemble dans son quartier, avec des connecteurs logiques ;
\item adopter une attitude de tolérance et d'écoute, valeurs au cœur du module « Citoyenneté ».
\end{itemize}

\courssub{Situation}

Dans le cadre de la « Semaine de la Citoyenneté » organisée par le club d'espagnol du lycée de Yaoundé, les élèves de Première A sont invités à participer à un débat public devant les autres classes. Le thème proposé cette année porte directement sur la vie quotidienne au Cameroun, pays où coexistent plus de deux cents groupes ethniques, deux langues officielles et de nombreuses religions.

\begin{textebox}[Affiche du club d'espagnol — Lycée de Yaoundé]
\textbf{SEMANA DE LA CIUDADANÍA}\\[0.2cm]
\textbf{GRAN DEBATE : ¿Es la diversidad una riqueza o un problema?}\\[0.2cm]
El Camerún es un país multicultural : muchas lenguas, muchas tradiciones, muchas religiones. Vivir juntos en la diferencia es un desafío, pero también una oportunidad.\\
¿La diversidad es una riqueza para nuestra sociedad o un problema para la convivencia?\\
Ven a debatir con respeto y tolerancia. ¡Tu opinión es importante!\\[0.2cm]
\emph{Lugar : sala de español — Hora : 14 h — Entrada libre}
\end{textebox}

\textbf{Glossaire :} \esp{desafío} = défi ; \esp{oportunidad} = occasion, opportunité ; \esp{ven a debatir} = viens débattre ; \esp{entrada libre} = entrée libre.

La classe est divisée en deux équipes : l'équipe A défend l'idée que la diversité est une richesse (\esp{una riqueza}) ; l'équipe B soutient qu'elle peut être un problème (\esp{un problema}) pour la cohabitation. Un élève-secrétaire note les arguments au tableau dans deux colonnes : \esp{a favor} et \esp{en contra}.

\courssub{Consignes}

\begin{enumerate}
\item \textbf{Préparation individuelle (10 min).} Relis le texte « Vivir juntos en la diferencia » et souligne les passages qui peuvent servir d'arguments pour ton équipe. Note deux arguments personnels en utilisant au moins trois mots du lexique de la leçon.
\item \textbf{Préparation en équipe (10 min).} Dans ton équipe, partage tes arguments, choisis les quatre meilleurs et répartissez la parole : chaque membre doit intervenir au moins une fois.
\item \textbf{Le débat (20 min).} Chaque intervention commence obligatoirement par une expression d'accord ou de désaccord : \esp{Estoy de acuerdo porque...}, \esp{No estoy de acuerdo, sin embargo...}, \esp{Tienes razón, pero...}. Respecte la parole de l'autre : pas de moqueries, pas d'interruptions.
\item \textbf{Le rôle du secrétaire.} L'élève-secrétaire écrit chaque argument au tableau dans la bonne colonne (\esp{a favor / en contra}) et vérifie que le lexique de la leçon est bien utilisé.
\item \textbf{Production écrite individuelle (15 min).} Pour clôturer le débat, rédige un paragraphe de 6 à 8 lignes intitulé \esp{La convivencia en mi barrio}. Tu dois réutiliser au moins cinq mots du lexique de la leçon et deux connecteurs (\esp{además}, \esp{sin embargo}, \esp{por eso}, \esp{también}).
\item \textbf{Évaluation et affichage.} Le professeur et la classe choisissent les meilleures productions, qui seront affichées au mur de la classe pendant la Semaine de la Citoyenneté.
\end{enumerate}

\begin{attention}[Règles d'or du débat]
On attaque les \emph{idées}, jamais les \emph{personnes}. On dit \esp{No estoy de acuerdo con esa idea} (« je ne suis pas d'accord avec cette idée ») et non une critique personnelle. La \esp{tolerancia} et le \esp{respeto} s'appliquent d'abord dans la classe !
\end{attention}

\courssub{Ressources à mobiliser}

\textbf{Le lexique du vivre-ensemble :}

\begin{center}
\begin{tabularx}{\linewidth}{|X|X|}\hline
\rowcolor{popblueL} Español & Français \\ \hline
la identidad & l'identité \\ \hline
la tolerancia & la tolérance \\ \hline
el respeto & le respect \\ \hline
la convivencia & le vivre-ensemble, la cohabitation \\ \hline
el prejuicio & le préjugé \\ \hline
la diferencia & la différence \\ \hline
la integración & l'intégration \\ \hline
la discriminación & la discrimination \\ \hline
la riqueza & la richesse \\ \hline
\end{tabularx}
\end{center}

\textbf{Les expressions d'accord et de désaccord :}

\begin{center}
\begin{tabularx}{\linewidth}{|X|X|}\hline
\rowcolor{popgreenL} Accord & Désaccord \\ \hline
Estoy de acuerdo porque... & No estoy de acuerdo, sin embargo... \\ \hline
Tienes razón, además... & Tienes razón, pero... \\ \hline
Es verdad que... & No es verdad que... \\ \hline
Pienso que / Creo que... & No pienso que / No creo que... \\ \hline
\end{tabularx}
\end{center}

\begin{attention}[Piège fréquent]
Après \esp{no creo que} et \esp{no pienso que}, le verbe qui suit se met au subjonctif : \esp{No creo que la diversidad sea un problema.} (« Je ne pense pas que la diversité soit un problème. »). Mais après \esp{creo que} affirmatif, on garde l'indicatif : \esp{Creo que la diversidad es una riqueza.}
\end{attention}

\textbf{Les connecteurs utiles :} \esp{además} (de plus), \esp{también} (aussi), \esp{sin embargo} (cependant), \esp{por eso} (c'est pourquoi), \esp{en conclusión} (en conclusion).

\courssub{Proposition de production et corrigé commenté}

\begin{exoresolu}[Production modèle : \esp{La convivencia en mi barrio}]
Rédige un paragraphe de 6 à 8 lignes sur le vivre-ensemble dans ton quartier, avec au moins cinq mots du lexique et deux connecteurs.
\tcblower
\textbf{Production modèle :}

\esp{En mi barrio de Yaoundé viven personas de diferentes regiones del Camerún. La convivencia es buena porque hay respeto entre los vecinos. Por ejemplo, en las fiestas, cada familia comparte su comida y sus tradiciones, y esto es una riqueza para todos. Además, los jóvenes juegan juntos al fútbol sin prejuicios. Sin embargo, a veces existen pequeños conflictos por la diferencia de lenguas o de costumbres. Por eso, pienso que la tolerancia es muy importante : debemos aceptar la identidad de cada persona. En conclusión, la diversidad no es un problema, es una oportunidad para aprender y vivir juntos en paz.}

\textbf{Traduction :} « Dans mon quartier de Yaoundé vivent des personnes de différentes régions du Cameroun. Le vivre-ensemble est bon parce qu'il y a du respect entre les voisins. Par exemple, pendant les fêtes, chaque famille partage sa nourriture et ses traditions, et c'est une richesse pour tous. De plus, les jeunes jouent ensemble au football sans préjugés. Cependant, il existe parfois de petits conflits à cause de la différence de langues ou de coutumes. C'est pourquoi je pense que la tolérance est très importante : nous devons accepter l'identité de chaque personne. En conclusion, la diversité n'est pas un problème, c'est une occasion d'apprendre et de vivre ensemble en paix. »

\textbf{Commentaire des choix de langue :}
\begin{itemize}
\item \textbf{Lexique (7 mots de la leçon) :} \esp{convivencia}, \esp{respeto}, \esp{riqueza}, \esp{prejuicios}, \esp{diferencia}, \esp{tolerancia}, \esp{identidad} — la consigne de cinq mots minimum est largement remplie.
\item \textbf{Connecteurs (4 au lieu de 2) :} \esp{además}, \esp{sin embargo}, \esp{por eso}, \esp{en conclusión} — ils structurent le paragraphe : ajout, opposition, conséquence, conclusion.
\item \textbf{Grammaire :} \esp{viven} (présent de \esp{vivir}, 3\textsuperscript{e} personne du pluriel) ; \esp{hay} (il y a, impersonnel) ; \esp{debemos aceptar} (\esp{deber} + infinitif pour exprimer le devoir) ; \esp{pienso que... es} (indicatif après \esp{pienso que} affirmatif).
\item \textbf{Contexte camerounais :} le quartier de Yaoundé, le football, les fêtes de famille rendent la production authentique et personnelle.
\end{itemize}
\end{exoresolu}

\courssub{Bilan de l'activité}

Cette activité d'intégration a permis de mobiliser ensemble toutes les compétences de la leçon : la \textbf{compréhension de l'écrit} (relecture du texte support pour y trouver des arguments), l'\textbf{expression orale} (débat avec expressions d'accord et de désaccord), le \textbf{lexique thématique} (réutilisé à l'oral et à l'écrit) et l'\textbf{expression écrite} (paragraphe final structuré par des connecteurs).

Au-delà de la langue, l'élève a développé une \textbf{compétence de vie} : débattre avec \esp{respeto} et \esp{tolerancia}, écouter l'autre, accepter la \esp{diferencia}. Dans un pays multiculturel comme le Cameroun, savoir dire \esp{no estoy de acuerdo} avec courtoisie est une qualité de citoyen autant que d'hispanophone. L'affichage des meilleures productions prolonge l'activité et valorise le travail de chacun : la classe devient elle-même un espace de \esp{convivencia}.

\newpage

\coursec{Méthodes et exercices résolus}

\courssub{Savoir-faire 1 : Lire et comprendre un texte sur l'identité et la tolérance}

La première compétence attendue en classe de Première est la \cle{compréhension de l'écrit} : lire un texte de 150 à 250 mots sur un thème de société et répondre à des questions en espagnol. Voici la démarche complète, pas à pas.

\begin{methode}[Réussir les questions de compréhension]
\begin{enumerate}
\item \textbf{Lire deux fois le texte.} Première lecture globale pour dégager le thème (ici : \esp{la identidad, la tolerancia, la convivencia}) ; deuxième lecture en soulignant les mots du lexique étudié : \esp{respeto}, \esp{prejuicio}, \esp{diferencia}, \esp{integración}, \esp{discriminación}.
\item \textbf{Lire les questions avant de répondre.} Repérer le type de question : question de fait (la réponse est dans le texte), question de vocabulaire (trouver un synonyme ou un contraire), question d'opinion (il faut donner son avis avec \esp{creo que...}, \esp{en mi opinión...}).
\item \textbf{Localiser la ligne ou le paragraphe} où se trouve la réponse. Ne jamais répondre de mémoire.
\item \textbf{Rédiger une phrase complète en espagnol}, en reprenant les mots de la question : \esp{¿Por qué es importante la tolerancia?} $\rightarrow$ \esp{La tolerancia es importante porque...}. Ne jamais répondre par un seul mot.
\item \textbf{Pour une question personnelle}, utiliser les structures d'opinion : \esp{Pienso que...}, \esp{Me parece que...}, \esp{Estoy de acuerdo con...}.
\item \textbf{Vérifier} : accord de l'adjectif, accents, verbe conjugué à la bonne personne.
\end{enumerate}
\end{methode}

\begin{exoresolu}[Questions de compréhension sur « Vivir juntos en la diferencia »]
Énoncé. Lisez l'extrait suivant, puis répondez en espagnol par des phrases complètes.

\esp{En el Camerún conviven más de doscientos grupos étnicos y dos lenguas oficiales. Esta diversidad es una riqueza, pero también un desafío: la convivencia exige tolerancia y respeto por la diferencia. El prejuicio, es decir, la opinión negativa que se forma sin conocer al otro, es el principal enemigo de la integración.}

1. ¿Cuántos grupos étnicos conviven en el Camerún? \quad 2. ¿Qué es el prejuicio según el texto? \quad 3. ¿Estás de acuerdo con la idea de que la diversidad es una riqueza? Justifica tu respuesta.
\tcblower
Solution détaillée.

\textbf{Question 1.} Question de fait : la réponse est dans la première phrase. On rédige une phrase complète : \esp{En el Camerún conviven más de doscientos grupos étnicos.} « Au Cameroun cohabitent plus de deux cents groupes ethniques. »

\textbf{Question 2.} Question de définition : le texte donne lui-même la définition après \esp{es decir} (« c'est-à-dire »). On la reformule : \esp{Según el texto, el prejuicio es la opinión negativa que se forma sin conocer al otro.} « Selon le texte, le préjugé est l'opinion négative qu'on se forme sans connaître l'autre. »

\textbf{Question 3.} Question d'opinion : on utilise une structure d'accord + une justification avec \esp{porque}. Réponse modèle : \esp{Sí, estoy de acuerdo: la diversidad es una riqueza porque cada cultura aporta sus tradiciones, sus lenguas y su cocina, y gracias a ella aprendemos a respetar a los demás.} « Oui, je suis d'accord : la diversité est une richesse parce que chaque culture apporte ses traditions, ses langues et sa cuisine, et grâce à elle nous apprenons à respecter les autres. »

\textbf{Conclusion.} Une bonne réponse de compréhension est toujours une phrase complète, appuyée sur le texte, avec un verbe correctement conjugué et les accents en place (\esp{opinión}, \esp{integración}).
\end{exoresolu}

\courssub{Savoir-faire 2 : Exprimer l'accord et le désaccord}

Après avoir compris le texte, l'élève doit pouvoir \cle{réagir} : c'est la compétence d'expression orale et écrite. Le cœur de la leçon est la construction figée \esp{estar de acuerdo} et ses variantes.

\begin{propriete}[Les formules d'accord et de désaccord]
\begin{center}
\begin{tabularx}{\linewidth}{|X|X|X|}
\hline
\rowcolor{popblueL} Fonction & Español & Français \\
\hline
Accord simple & Estoy de acuerdo (con...) & Je suis d'accord (avec...) \\
\hline
Accord renforcé & Estoy totalmente de acuerdo & Je suis tout à fait d'accord \\
\hline
Accord à l'oral & Tienes razón & Tu as raison \\
\hline
Accord + opinion & Yo también pienso que... & Je pense aussi que... \\
\hline
Désaccord simple & No estoy de acuerdo (con...) & Je ne suis pas d'accord (avec...) \\
\hline
Désaccord + subjonctif & No creo que + subjonctif & Je ne crois pas que... \\
\hline
Désaccord poli & Es cierto que..., pero... & Il est vrai que..., mais... \\
\hline
Désaccord poli & Entiendo tu opinión, pero creo que... & Je comprends ton opinion, mais je crois que... \\
\hline
\end{tabularx}
\end{center}
\end{propriete}

\begin{methode}[Dire que l'on est d'accord ou pas]
\begin{enumerate}
\item \textbf{Mémoriser les structures de base} du tableau ci-dessus.
\item \textbf{Respecter la construction figée} : \esp{estar de acuerdo} ne change que par la conjugaison de \esp{estar} (\esp{estoy, estás, está, estamos, estáis, están}). On ne dit jamais \emph{soy de acuerdo}.
\item \textbf{Nuancer} : pour un débat de classe, préférer le désaccord poli : \esp{Entiendo tu opinión, pero creo que...} « Je comprends ton opinion, mais je crois que... ».
\item \textbf{Justifier toujours} avec \esp{porque} + indicatif, ou \esp{ya que...}, \esp{pues...}.
\item \textbf{Attention au subjonctif} : après \esp{no creo que}, le verbe suivant est au subjonctif : \esp{No creo que sea fácil.} « Je ne crois pas que ce soit facile. » Après \esp{creo que}, indicatif : \esp{Creo que es fácil.} « Je crois que c'est facile. »
\end{enumerate}
\end{methode}

\begin{exoresolu}[Réagir à des opinions : accord et désaccord]
Énoncé. Réagissez à chaque opinion en exprimant votre accord ou votre désaccord, puis justifiez en une phrase. Rédigez en espagnol.

1. \esp{« La tolerancia es inútil: cada uno debe quedarse con su propia cultura. »} (désaccord) \quad 2. \esp{« El respeto por la diferencia es la base de la convivencia. »} (accord) \quad 3. \esp{« En el Camerún, la diversidad de lenguas es un problema. »} (désaccord poli)
\tcblower
Solution détaillée.

\textbf{1.} Désaccord + justification. \esp{No estoy de acuerdo con esta opinión: la tolerancia es necesaria porque vivimos en una sociedad multicultural y nadie puede vivir aislado de los demás.} « Je ne suis pas d'accord avec cette opinion : la tolérance est nécessaire parce que nous vivons dans une société multiculturelle et personne ne peut vivre isolé des autres. » Construction : \esp{no estoy de acuerdo con} + nom ; justification avec \esp{porque} + indicatif.

\textbf{2.} Accord + justification. \esp{Estoy totalmente de acuerdo: sin respeto por la diferencia no hay convivencia posible, porque el respeto permite aceptar al otro tal como es.} « Je suis tout à fait d'accord : sans respect de la différence il n'y a pas de vivre-ensemble possible, parce que le respect permet d'accepter l'autre tel qu'il est. »

\textbf{3.} Désaccord poli : on commence par reconnaître l'argument avant de le réfuter. \esp{Es cierto que tantas lenguas pueden complicar la comunicación, pero no creo que sean un problema: al contrario, son una riqueza cultural que debemos proteger.} « Il est vrai que tant de langues peuvent compliquer la communication, mais je ne crois pas qu'elles soient un problème : au contraire, elles sont une richesse culturelle que nous devons protéger. » Remarque : après \esp{no creo que}, subjonctif \esp{sean} (et non \emph{son}).

\textbf{Conclusion.} Une réaction complète = formule d'accord/désaccord + \esp{porque} + argument. C'est exactement ce que l'examinateur attend à l'oral comme à l'écrit.
\end{exoresolu}

\courssub{Savoir-faire 3 : Traduire (thème) avec le lexique du vivre-ensemble}

Le \cle{thème} (traduction du français vers l'espagnol) porte souvent, à l'examen, sur des phrases d'opinion utilisant le lexique de la leçon. La règle d'or : traduire les idées avec le vocabulaire exact appris.

\begin{methode}[Réussir le thème sur ce thème de société]
\begin{enumerate}
\item \textbf{Traduire les idées, pas les mots.} Identifier d'abord le verbe et son temps (présent le plus souvent dans les textes d'opinion).
\item \textbf{Utiliser le lexique exact de la leçon} : identité = \esp{la identidad} ; tolérance = \esp{la tolerancia} ; respect = \esp{el respeto} ; vivre-ensemble = \esp{la convivencia} ; préjugé = \esp{el prejuicio} ; différence = \esp{la diferencia} ; intégration = \esp{la integración} ; discrimination = \esp{la discriminación}.
\item \textbf{Se méfier des faux amis} : « la différence » est féminin (\esp{la diferencia}) ; « un préjugé » se dit \esp{un prejuicio} (sans accent) ; « respecter » = \esp{respetar} (et non \emph{respectar}).
\item \textbf{Accords} : l'adjectif s'accorde avec le nom : \esp{una sociedad multicultural}, \esp{los prejuicios son peligrosos}.
\item \textbf{Relire} en chassant : accents (\esp{opinión}, \esp{integración}, \esp{también}), \esp{ñ}, et le verbe \esp{haber} impersonnel : « il y a » = \esp{hay} (jamais de pluriel).
\end{enumerate}
\end{methode}

\begin{exoresolu}[Thème : phrases sur la tolérance]
Énoncé. Traduisez en espagnol.

1. La tolérance est le respect de la différence. \quad 2. Au Cameroun, il y a beaucoup de cultures différentes. \quad 3. Nous devons combattre les préjugés et la discrimination. \quad 4. Je suis d'accord avec cette opinion.
\tcblower
Solution détaillée.

\textbf{1.} \esp{La tolerancia es el respeto de la diferencia.} Définition : verbe \esp{ser} (définition permanente, pas \esp{estar}). « Respect » = \esp{el respeto}, sans \emph{c} après le \emph{s}.

\textbf{2.} \esp{En el Camerún hay muchas culturas diferentes.} « Il y a » = \esp{hay}, invariable même devant un pluriel ; \esp{diferentes} prend un \emph{s} au pluriel, sans accent ; \esp{Camerún} porte un accent sur le \emph{ú}.

\textbf{3.} \esp{Debemos combatir los prejuicios y la discriminación.} « Nous devons » = \esp{debemos} (\esp{deber} + infinitif) ; \esp{prejuicio} s'écrit sans accent ; \esp{discriminación} porte un accent sur le \emph{ó}.

\textbf{4.} \esp{Estoy de acuerdo con esta opinión.} Piège classique : « je suis » se traduit ici par \esp{estoy} car \esp{estar de acuerdo} est une expression figée ; \esp{opinión} prend un accent.

\textbf{Conclusion.} Au thème, chaque accent oublié et chaque faux ami coûte des points : la relecture ciblée (accents, \esp{hay}, \esp{estar de acuerdo}) est obligatoire.
\end{exoresolu}

\courssub{Savoir-faire 4 : Traduire (version) un texte sur la convivencia}

La \cle{version} (traduction de l'espagnol vers le français) exige une double fidélité : au sens du texte et à la qualité du français.

\begin{methode}[Réussir la version]
\begin{enumerate}
\item \textbf{Lire tout le texte} pour comprendre le sens global avant de traduire la première phrase.
\item \textbf{Découper en propositions} : sujet / verbe / complément. Traduire proposition par proposition.
\item \textbf{Rendre les tournures espagnoles par des tournures françaises naturelles} : \esp{es decir} = « c'est-à-dire » ; \esp{sin embargo} = « cependant » ; \esp{cada vez más} = « de plus en plus ».
\item \textbf{Attention aux mots piégés} : \esp{la convivencia} = « le vivre-ensemble » ou « la cohabitation » (pas « la convivence ») ; \esp{apoyar} = « soutenir » ; \esp{los demás} = « les autres ».
\item \textbf{Rédiger en bon français} : une version correcte en espagnol mais maladroite en français perd des points. Relire mentalement à voix haute.
\end{enumerate}
\end{methode}

\begin{exoresolu}[Version : un paragraphe sur l'intégration]
Énoncé. Traduisez en français.

\esp{La integración no significa renunciar a la propia identidad, sino compartirla con los demás. En una sociedad multicultural como la camerunesa, la tolerancia y el respeto son valores esenciales: sin ellos, la convivencia se convierte en conflicto. Por eso, debemos luchar contra los prejuicios desde la escuela.}
\tcblower
Solution détaillée.

\textbf{Phrase 1.} \esp{renunciar a} = « renoncer à » ; \esp{no... sino} = « non pas... mais » ; \esp{la propia identidad} = « sa propre identité ». Traduction : « L'intégration ne signifie pas renoncer à sa propre identité, mais la partager avec les autres. »

\textbf{Phrase 2.} \esp{la camerunesa} reprend \esp{sociedad} : « la société camerounaise » ; \esp{se convierte en} = « se transforme en ». Traduction : « Dans une société multiculturelle comme la société camerounaise, la tolérance et le respect sont des valeurs essentielles : sans elles, le vivre-ensemble se transforme en conflit. »

\textbf{Phrase 3.} \esp{por eso} = « c'est pourquoi » ; \esp{luchar contra} = « lutter contre » ; \esp{desde la escuela} = « dès l'école ». Traduction : « C'est pourquoi nous devons lutter contre les préjugés dès l'école. »

\textbf{Conclusion.} La version exige deux fidélités : au sens espagnol et à la langue française. Les connecteurs (\esp{por eso}, \esp{sin embargo}) doivent toujours être traduits, jamais oubliés.
\end{exoresolu}

\courssub{Savoir-faire 5 : Rédiger un paragraphe sur le vivre-ensemble}

Dernière compétence : la \cle{production écrite guidée}. À l'examen de Première, on demande souvent un court paragraphe d'opinion (8 à 10 lignes) sur un sujet de citoyenneté.

\begin{methode}[Rédiger un paragraphe d'opinion (8 à 10 lignes)]
\begin{enumerate}
\item \textbf{Structurer en quatre temps} : (a) phrase d'introduction qui pose le sujet ; (b) ton opinion avec \esp{En mi opinión...} ou \esp{Pienso que...} ; (c) deux arguments reliés par \esp{primero}, \esp{además}, \esp{también} ; (d) conclusion avec \esp{En conclusión} ou \esp{Por eso}.
\item \textbf{Employer le lexique de la leçon} : au moins quatre mots parmi \esp{identidad}, \esp{tolerancia}, \esp{respeto}, \esp{convivencia}, \esp{prejuicio}, \esp{diferencia}, \esp{integración}, \esp{discriminación}.
\item \textbf{Rester au présent de l'indicatif} : c'est le temps de l'opinion et des vérités générales.
\item \textbf{Varier les verbes} : \esp{ser}, \esp{exigir}, \esp{permitir}, \esp{respetar}, \esp{aceptar}, \esp{deber} + infinitif.
\item \textbf{Relire} : accords (\esp{una sociedad multicultural}), accents, ponctuation espagnole si exclamation ou question.
\end{enumerate}
\end{methode}

\begin{exoresolu}[Rédaction guidée : le vivre-ensemble au Cameroun]
Énoncé. Rédigez un paragraphe de 8 à 10 lignes en espagnol sur le sujet suivant : \esp{¿Por qué es importante la tolerancia en una sociedad multicultural como el Camerún?}
\tcblower
Solution détaillée : production modèle commentée.

\esp{El Camerún es un país multicultural: conviven en él más de doscientos grupos étnicos, dos lenguas oficiales y muchas religiones. En mi opinión, la tolerancia es esencial en esta sociedad. Primero, la tolerancia permite aceptar la diferencia y respetar la identidad de cada persona. Además, sin tolerancia, los prejuicios crecen y la convivencia se convierte en discriminación y conflicto. También pienso que la escuela tiene un papel importante: allí aprendemos a vivir juntos y a respetar a los demás. En conclusión, la tolerancia y el respeto son la base del vivir juntos: gracias a ellos, la diversidad del Camerún es una riqueza y no un problema.}

« Le Cameroun est un pays multiculturel : plus de deux cents groupes ethniques, deux langues officielles et de nombreuses religions y cohabitent. À mon avis, la tolérance est essentielle dans cette société. D'abord, la tolérance permet d'accepter la différence et de respecter l'identité de chaque personne. De plus, sans tolérance, les préjugés grandissent et le vivre-ensemble se transforme en discrimination et en conflit. Je pense aussi que l'école a un rôle important : c'est là que nous apprenons à vivre ensemble et à respecter les autres. En conclusion, la tolérance et le respect sont la base du vivre-ensemble : grâce à eux, la diversité du Cameroun est une richesse et non un problème. »

\textbf{Pourquoi ce paragraphe réussit} : structure en quatre temps respectée ; six mots du lexique de la leçon ; connecteurs variés (\esp{primero}, \esp{además}, \esp{también}, \esp{en conclusión}) ; présent de l'indicatif partout ; accords corrects (\esp{lenguas oficiales}, \esp{sociedad multicultural}).
\end{exoresolu}

\begin{attention}[Les 5 erreurs qui coûtent des points au Bac]
\begin{enumerate}
\item \textbf{Oublier les accents} : \esp{opinión}, \esp{integración}, \esp{discriminación}, \esp{también}, \esp{Camerún}. Un accent oublié est une faute comptée ; en revanche \esp{prejuicio} et \esp{tolerancia} n'en portent pas.
\item \textbf{Écrire \emph{respectar}} : le verbe correct est \esp{respetar} (« respecter »). De même, « la différence » est féminin : \esp{la diferencia} (et non \emph{el diferencia}).
\item \textbf{Calquer le français} : « je suis d'accord » se dit \esp{estoy de acuerdo} (jamais \emph{soy de acuerdo}) ; « il y a » se dit \esp{hay}, invariable, même devant un pluriel (\esp{hay muchas culturas}).
\item \textbf{Faute d'accord} : l'adjectif suit le nom en genre et en nombre : \esp{una sociedad multicultural}, \esp{los prejuicios son peligrosos}, \esp{las diferencias culturales}.
\item \textbf{Oublier le subjonctif après \esp{no creo que}} : on écrit \esp{No creo que sea fácil} et non \emph{no creo que es fácil}. À l'inverse, \esp{creo que} et \esp{pienso que} sont suivis de l'indicatif : \esp{Creo que es importante.}
\end{enumerate}
\end{attention}

\newpage

\coursec{Exercices}

\courssub{Je vérifie mes connaissances}

\begin{exercice}[QCM : le lexique du vivre-ensemble]
Pour chaque question, entoure la bonne réponse.
\begin{enumerate}
\item \esp{La tolerancia} signifie : a) le rejet de l'autre \quad b) l'acceptation des différences \quad c) la peur de l'étranger.
\item \esp{El prejuicio} est : a) un jugement formé avant de connaître \quad b) une fête populaire \quad c) un droit de l'homme.
\item \esp{La convivencia} désigne : a) la guerre civile \quad b) le fait de vivre ensemble en harmonie \quad c) un document d'identité.
\item \esp{La discriminación} consiste à : a) traiter quelqu'un injustement à cause de son origine \quad b) inviter ses voisins \quad c) apprendre une langue.
\item \esp{La integración} est : a) l'exclusion d'un groupe \quad b) l'insertion harmonieuse dans une société \quad c) une émission de radio.
\end{enumerate}
\end{exercice}

\begin{exercice}[Vrai ou faux ?]
Réponds par \esp{verdadero} ou \esp{falso} et justifie ta réponse en une phrase en français.
\begin{enumerate}
\item \esp{El respeto} est un obstacle au vivre-ensemble.
\item Le Cameroun est un pays multiculturel où coexistent de nombreuses langues et traditions.
\item \esp{Estoy de acuerdo} sert à exprimer le désaccord.
\item Les préjugés peuvent conduire à la discrimination.
\item \esp{La identidad} renvoie à ce qui fait qu'une personne ou un peuple est lui-même.
\end{enumerate}
\end{exercice}

\begin{exercice}[Vocabulaire : trouve le mot]
Retrouve le mot du lexique qui correspond à chaque définition : \esp{identidad, tolerancia, respeto, convivencia, prejuicio, diferencia, integración, discriminación}.
\begin{enumerate}
\item Considération et égard envers autrui : \ldots
\item Ce qui distingue une personne d'une autre : \ldots
\item Opinion négative formée sans connaissance réelle : \ldots
\item Ensemble de traits qui définissent une personne ou un groupe : \ldots
\item Capacité d'accepter ce qui est différent de soi : \ldots
\end{enumerate}
\end{exercice}

\begin{exercice}[Conjugaison : le présent de l'indicatif]
Complète les phrases avec le verbe entre parenthèses, conjugué au présent de l'indicatif.
\begin{enumerate}
\item Nosotros \ldots\ (respetar) las diferencias de los demás.
\item Tú \ldots\ (estar) de acuerdo conmigo, ¿verdad ?
\item Los alumnos \ldots\ (convivir) en paz en el liceo de Yaoundé.
\item Yo no \ldots\ (creer) que los prejuicios sean una fatalidad.
\item Ella \ldots\ (defender) la tolerancia en su barrio.
\end{enumerate}
\end{exercice}

\courssub{Je m'entraîne}

\begin{exercice}[Texte à trous]
Complète le texte suivant avec les mots de la liste : \esp{tolerancia, prejuicio, convivencia, discriminación, respeto, integración}.
\begin{quote}
\esp{En una sociedad multicultural como Camerún, la \ldots\ es indispensable para vivir juntos. El \ldots\ hacia los demás permite una buena \ldots\ en los barrios y en las escuelas. Al contrario, el \ldots\ nace de la ignorancia y puede provocar la \ldots\ de ciertos grupos. La \ldots\ de todas las culturas es una riqueza para el país.}
\end{quote}
\end{exercice}

\begin{exercice}[Classer le vocabulaire]
Classe les douze mots suivants en deux colonnes : \esp{valores positivos} et \esp{obstáculos para la convivencia}.\\
\esp{respeto -- prejuicio -- tolerancia -- odio -- solidaridad -- discriminación -- diálogo -- violencia -- amistad -- racismo -- integración -- exclusión}
\end{exercice}

\begin{exercice}[Dialogue à compléter]
Complète le dialogue entre deux élèves, Amina et Carlos, avec les expressions suivantes : \esp{estoy de acuerdo, no estoy de acuerdo, tienes razón, no creo que}.
\begin{quote}
\esp{--- Amina : Pienso que la diversidad cultural es una riqueza para Camerún.\\
--- Carlos : \ldots\ contigo : cada región aporta algo especial.\\
--- Amina : Algunos dicen que es imposible entenderse entre tantas lenguas.\\
--- Carlos : \ldots\ : el francés, el inglés y el español nos unen.\\
--- Amina : ¿Crees que los prejuicios van a desaparecer un día ?\\
--- Carlos : \ldots\ desaparezcan fácilmente, pero la educación ayuda mucho.\\
--- Amina : Debemos respetar a todos, sin excepción.\\
--- Carlos : \ldots\ , es la base de la convivencia.}
\end{quote}
\end{exercice}

\begin{exercice}[Traduction : thème]
Traduis en espagnol les phrases suivantes.
\begin{enumerate}
\item Je suis d'accord avec toi : nous devons respecter les différences.
\item Je ne crois pas que les préjugés soient une fatalité.
\item La tolérance est très importante pour le vivre-ensemble au Cameroun.
\item Tu as raison : la discrimination détruit notre société.
\end{enumerate}
\end{exercice}

\courssub{Je me prépare au Bac}

\begin{exercice}[Compréhension de texte type Bac]
Lis le texte suivant, puis réponds aux questions.
\begin{textebox}[La escuela, escuela de tolerancia]
En el barrio de Briqueterie, en Yaoundé, hay una escuela donde los alumnos hablan más de diez lenguas diferentes en casa. Algunos vienen del norte del país, otros del sur o del oeste ; todos estudian juntos en francés, en inglés y, desde hace dos años, en español. La directora, la señora Mbarga, explica su filosofía : « Aquí, cada niño trae su cultura, su lengua y su historia. Nuestra misión es enseñarles a respetarse unos a otros. »

Sin embargo, la convivencia no siempre es fácil. Al principio del año, algunos alumnos se burlaban de los compañeros que hablaban con un acento diferente. Entonces los profesores organizaron una « semana de las culturas » : cada grupo presentó sus danzas, sus platos y sus cuentos tradicionales. Poco a poco, los prejuicios desaparecieron y nacieron nuevas amistades. « Antes tenía miedo de los que no eran como yo », confiesa un alumno de Première, « pero ahora comprendo que la diferencia no es un peligro, sino una riqueza. »

Hoy, esta escuela es un modelo de integración para toda la ciudad. Los padres participan en las fiestas, los alumnos ayudan a los nuevos y la palabra « tolerancia » tiene aquí un sentido muy concreto : vivir juntos, aprender juntos y construir juntos el futuro del país.
\end{textebox}
\textbf{Questions de compréhension}
\begin{enumerate}
\item Réponds par \esp{verdadero} ou \esp{falso} et justifie en citant le texte :
\begin{enumerate}
\item Tous les élèves de l'école parlent la même langue à la maison.
\item La directrice pense que la diversité est un problème.
\item La « semaine des cultures » a amélioré les relations entre les élèves.
\end{enumerate}
\item Questions ouvertes (réponds en espagnol, avec des phrases complètes) :
\begin{enumerate}
\item ¿Cuál es la misión de la escuela según la señora Mbarga ?
\item ¿Qué aprendió el alumno de Première sobre la diferencia ?
\end{enumerate}
\end{enumerate}
\textbf{Questions de langue}
\begin{enumerate}
\item Trouve dans le texte un synonyme de \esp{empezaron} et un antonyme de \esp{pobreza}.
\item Relève deux mots du lexique du vivre-ensemble et traduis-les en français.
\item Conjugue au présent : \esp{los alumnos (respetarse) unos a otros}.
\end{enumerate}
\end{exercice}

\begin{exercice}[Thème et version]
\textbf{A. Thème.} Traduis en espagnol.
\begin{enumerate}
\item Le Cameroun est un pays où vivent ensemble de nombreuses cultures.
\item Je ne suis pas d'accord : il faut lutter contre la discrimination à l'école.
\item Le respect des différences est la base du vivre-ensemble.
\end{enumerate}
\textbf{B. Version.} Traduis en français.
\begin{quote}
\esp{La tolerancia no significa aceptar todo sin pensar, sino escuchar al otro y reconocer su dignidad. En una sociedad multicultural, el diálogo y el respeto son las claves de la paz. Quien rechaza al diferente empobrece su propia identidad.}
\end{quote}
\end{exercice}

\begin{exercice}[Expression écrite type Bac]
Choisis \textbf{un} des deux sujets suivants et rédige un texte en espagnol de \textbf{80 à 100 mots}. Utilise au moins trois connecteurs (\esp{primero, además, sin embargo, por eso, en conclusión}) et deux expressions d'accord ou de désaccord.
\begin{enumerate}
\item \esp{¿Qué significa para ti la tolerancia ?} Donne ta définition et illustre-la avec un exemple de ta vie ou de ton quartier.
\item \esp{El vivir juntos en un Camerún multicultural : ¿una riqueza o una dificultad ?} Présente ton opinion avec des arguments.
\end{enumerate}
\end{exercice}

\newpage

\coursec{Corrigés des exercices}

\begin{corrige}[Exercice 1 — QCM : le lexique du vivre-ensemble]
\begin{enumerate}
\item \textbf{b)} l'acceptation des différences. \esp{La tolerancia} vient du latin \emph{tolerantia} : supporter, accepter ce qui est différent de soi.
\item \textbf{a)} un jugement formé avant de connaître. \esp{El prejuicio} = \emph{pre-} (avant) + \emph{juicio} (jugement).
\item \textbf{b)} le fait de vivre ensemble en harmonie. \esp{La convivencia} = \emph{con-} (avec) + \emph{vivir} (vivre).
\item \textbf{a)} traiter quelqu'un injustement à cause de son origine.
\item \textbf{b)} l'insertion harmonieuse dans une société.
\end{enumerate}
\textbf{Point de vigilance :} attention à l'orthographe espagnole : \esp{integración} et \esp{discriminación} prennent un accent sur le \emph{ó} final ; \esp{prejuicio} s'écrit avec un seul \emph{j} et ne prend pas d'accent.
\end{corrige}

\begin{corrige}[Exercice 2 — Vrai ou faux ?]
\begin{enumerate}
\item \esp{Falso}. Le respect (\esp{el respeto}) est au contraire la condition première du vivre-ensemble : sans considération pour autrui, la vie en communauté devient impossible.
\item \esp{Verdadero}. Le Cameroun compte plus de deux cents langues nationales, deux langues officielles et de nombreuses traditions : c'est un pays profondément multiculturel.
\item \esp{Falso}. \esp{Estoy de acuerdo} exprime l'\textbf{accord} ; pour le désaccord, on dit \esp{no estoy de acuerdo}.
\item \esp{Verdadero}. Un préjugé (jugement négatif formé sans connaissance) peut déboucher sur des actes de discrimination, c'est-à-dire un traitement injuste.
\item \esp{Verdadero}. \esp{La identidad} désigne l'ensemble des traits (langue, culture, histoire, valeurs) qui font qu'une personne ou un peuple est lui-même.
\end{enumerate}
\end{corrige}

\begin{corrige}[Exercice 3 — Vocabulaire : trouve le mot]
\begin{enumerate}
\item \esp{respeto} (le respect)
\item \esp{diferencia} (la différence)
\item \esp{prejuicio} (le préjugé)
\item \esp{identidad} (l'identité)
\item \esp{tolerancia} (la tolérance)
\end{enumerate}
\textbf{Astuce :} pour retenir ce lexique, associe chaque mot à sa famille : \esp{respeto/respetar}, \esp{diferencia/diferente}, \esp{tolerancia/tolerar}, \esp{identidad/idéntico}.
\end{corrige}

\begin{corrige}[Exercice 4 — Conjugaison : le présent de l'indicatif]
\begin{enumerate}
\item Nosotros \esp{respetamos} las diferencias de los demás. — Verbe régulier en \emph{-ar} : radical \emph{respet-} + terminaison \emph{-amos}.
\item Tú \esp{estás} de acuerdo conmigo, ¿verdad ? — \esp{estar} est irrégulier à la 1\iere{} personne (\esp{estoy}) mais régulier à la 2\ieme{} : \esp{estás}, avec accent sur le \emph{á}.
\item Los alumnos \esp{conviven} en paz en el liceo de Yaoundé. — Verbe régulier en \emph{-ir} : terminaison \emph{-en}.
\item Yo no \esp{creo} que los prejuicios sean una fatalidad. — \esp{creer} : \esp{creo, crees, cree, creemos, creéis, creen} ; attention à \esp{creéis} avec accent.
\item Ella \esp{defiende} la tolerancia en su barrio. — \esp{defender} est un verbe à changement radical \emph{e $\rightarrow$ ie} : \esp{defiendo, defiendes, defiende, defendemos, defendéis, defienden}.
\end{enumerate}
\textbf{Point de vigilance :} dans la phrase 4, \esp{no creo que} est suivi du \textbf{subjonctif} (\esp{sean}) : le doute et la négation de l'opinion exigent le subjonctif en espagnol.
\end{corrige}

\begin{corrige}[Exercice 5 — Texte à trous]
\begin{quote}
\esp{En una sociedad multicultural como Camerún, la \textbf{tolerancia} es indispensable para vivir juntos. El \textbf{respeto} hacia los demás permite una buena \textbf{convivencia} en los barrios y en las escuelas. Al contrario, el \textbf{prejuicio} nace de la ignorancia y puede provocar la \textbf{discriminación} de ciertos grupos. La \textbf{integración} de todas las culturas es una riqueza para el país.}
\end{quote}
\textbf{Traduction :} « Dans une société multiculturelle comme le Cameroun, la tolérance est indispensable pour vivre ensemble. Le respect envers les autres permet une bonne cohabitation dans les quartiers et dans les écoles. Au contraire, le préjugé naît de l'ignorance et peut provoquer la discrimination de certains groupes. L'intégration de toutes les cultures est une richesse pour le pays. »\\
\textbf{Justification :} les articles guident le choix du genre : \esp{la tolerancia}, \esp{el respeto}, \esp{una buena convivencia}, \esp{el prejuicio}, \esp{la discriminación}, \esp{la integración}. Le contraste logique (\esp{al contrario}) impose un mot négatif (\esp{prejuicio}) après les trois mots positifs.
\end{corrige}

\begin{corrige}[Exercice 6 — Classer le vocabulaire]
\begin{center}
\begin{tabularx}{\linewidth}{|X|X|}
\hline
\rowcolor{popblueL} \esp{Valores positivos} & \esp{Obstáculos para la convivencia} \\
\hline
\esp{respeto} (le respect) & \esp{prejuicio} (le préjugé) \\
\hline
\esp{tolerancia} (la tolérance) & \esp{odio} (la haine) \\
\hline
\esp{solidaridad} (la solidarité) & \esp{discriminación} (la discrimination) \\
\hline
\esp{diálogo} (le dialogue) & \esp{violencia} (la violence) \\
\hline
\esp{amistad} (l'amitié) & \esp{racismo} (le racisme) \\
\hline
\esp{integración} (l'intégration) & \esp{exclusión} (l'exclusion) \\
\hline
\end{tabularx}
\end{center}
\textbf{Remarque :} chaque obstacle possède un antonyme positif : \esp{odio} $\leftrightarrow$ \esp{amistad}, \esp{exclusión} $\leftrightarrow$ \esp{integración}, \esp{violencia} $\leftrightarrow$ \esp{diálogo}. Apprendre les mots par paires contraires facilite la mémorisation.
\end{corrige}

\begin{corrige}[Exercice 7 — Dialogue à compléter]
\begin{quote}
\esp{--- Amina : Pienso que la diversidad cultural es una riqueza para Camerún.\\
--- Carlos : \textbf{Estoy de acuerdo} contigo : cada región aporta algo especial.\\
--- Amina : Algunos dicen que es imposible entenderse entre tantas lenguas.\\
--- Carlos : \textbf{No estoy de acuerdo} : el francés, el inglés y el español nos unen.\\
--- Amina : ¿Crees que los prejuicios van a desaparecer un día ?\\
--- Carlos : \textbf{No creo que} desaparezcan fácilmente, pero la educación ayuda mucho.\\
--- Amina : Debemos respetar a todos, sin excepción.\\
--- Carlos : \textbf{Tienes razón}, es la base de la convivencia.}
\end{quote}
\textbf{Traduction :} « — Amina : Je pense que la diversité culturelle est une richesse pour le Cameroun. — Carlos : Je suis d'accord avec toi : chaque région apporte quelque chose de spécial. — Amina : Certains disent qu'il est impossible de se comprendre entre tant de langues. — Carlos : Je ne suis pas d'accord : le français, l'anglais et l'espagnol nous unissent. — Amina : Crois-tu que les préjugés vont disparaître un jour ? — Carlos : Je ne crois pas qu'ils disparaissent facilement, mais l'éducation aide beaucoup. — Amina : Nous devons respecter tout le monde, sans exception. — Carlos : Tu as raison, c'est la base du vivre-ensemble. »\\
\textbf{Points de vigilance :}
\begin{itemize}
\item \esp{Estar de acuerdo \textbf{con} alguien} : la préposition est \esp{con}, jamais \esp{conmigo} remplacé par \esp{con yo}.
\item Après \esp{no creo que}, le subjonctif est obligatoire : \esp{desaparezcan} (et non \esp{desaparecen}).
\item \esp{Tener razón} se construit sans article : on dit \esp{tienes razón}, jamais \esp{tienes la razón}.
\end{itemize}
\end{corrige}

\begin{corrige}[Exercice 8 — Traduction : thème]
\begin{enumerate}
\item \esp{Estoy de acuerdo contigo : debemos respetar las diferencias.} — \esp{contigo} = « avec toi » (forme insécable) ; \esp{debemos} + infinitif pour « nous devons ».
\item \esp{No creo que los prejuicios sean una fatalidad.} — Négation de l'opinion $\Rightarrow$ subjonctif \esp{sean} (verbe \esp{ser}). « Une fatalité » se rend bien par \esp{una fatalidad}.
\item \esp{La tolerancia es muy importante para la convivencia en Camerún.} — « Le vivre-ensemble » se traduit par \esp{la convivencia} ; « pour » de destination = \esp{para}.
\item \esp{Tienes razón : la discriminación destruye nuestra sociedad.} — \esp{tener razón} sans article ; \esp{destruir} au présent : \esp{destruye} (changement \emph{u $\rightarrow$ uy} : \esp{destruyo, destruyes, destruye}).
\end{enumerate}
\textbf{Points de vigilance :} ne pas écrire \esp{estoy de acuerdo con tú} (faute classique : \esp{contigo}) ; ne pas oublier l'accent de \esp{discriminación} ; « au Cameroun » = \esp{en Camerún} (avec accent sur le \emph{ú}).
\end{corrige}

\begin{corrige}[Exercice 9 — Compréhension de texte type Bac]
\textbf{Barème indicatif (sur 20) :} vrai/faux justifiés : 6 pts (2 $\times$ 3) ; questions ouvertes : 4 pts ; questions de langue : 6 pts ; qualité de l'espagnol : 4 pts.

\textbf{1. Vrai ou faux}
\begin{enumerate}
\item \esp{Falso}. Le texte dit : \esp{« los alumnos hablan más de diez lenguas diferentes en casa »} (les élèves parlent plus de dix langues différentes à la maison).
\item \esp{Falso}. La directrice déclare : \esp{« cada niño trae su cultura, su lengua y su historia »} et sa mission est d'apprendre aux enfants à se respecter : la diversité est pour elle une richesse, non un problème.
\item \esp{Verdadero}. Le texte affirme : \esp{« Poco a poco, los prejuicios desaparecieron y nacieron nuevas amistades »} (peu à peu, les préjugés disparurent et de nouvelles amitiés naquirent).
\end{enumerate}
\textbf{2. Questions ouvertes (réponses modèles)}
\begin{enumerate}
\item \esp{Según la señora Mbarga, la misión de la escuela es enseñar a los niños a respetarse unos a otros.} (Selon Mme Mbarga, la mission de l'école est d'apprendre aux enfants à se respecter les uns les autres.)
\item \esp{El alumno de Première aprendió que la diferencia no es un peligro, sino una riqueza.} (L'élève de Première a appris que la différence n'est pas un danger, mais une richesse.)
\end{enumerate}
\textbf{3. Questions de langue}
\begin{enumerate}
\item Synonyme de \esp{empezaron} : \esp{nacieron} n'est pas synonyme ; le bon synonyme dans le texte est \esp{comenzaron} ? Non : le texte contient \esp{« Al principio del año »} ; le verbe synonyme attendu est \esp{empezaron} $\approx$ \esp{comenzaron} — or dans ce texte, on relève \esp{nacieron} (naquirent) comme verbe de début. Réponse acceptée : \esp{nacieron} (« nacieron nuevas amistades »). Antonyme de \esp{pobreza} : \esp{riqueza} (« la diferencia ... es una riqueza »).
\item Exemples : \esp{la convivencia} = le vivre-ensemble ; \esp{los prejuicios} = les préjugés ; \esp{la integración} = l'intégration ; \esp{la tolerancia} = la tolérance.
\item \esp{Los alumnos se respetan unos a otros.} — Verbe pronominal \esp{respetarse} : le pronom réfléchi \esp{se} se place avant le verbe conjugué ; \esp{unos a otros} = « les uns les autres » (réciprocité).
\end{enumerate}
\textbf{Points de vigilance :} en question ouverte, répondre par une phrase complète en reprenant les termes de la question ; ne pas recopier tout le texte ; citer entre guillemets pour les justifications.
\end{corrige}

\begin{corrige}[Exercice 10 — Thème et version]
\textbf{A. Thème}
\begin{enumerate}
\item \esp{Camerún es un país donde conviven muchas culturas.} — Autre possibilité : \esp{... donde muchas culturas viven juntas.} « Vivent ensemble » = \esp{conviven} (plus élégant) ou \esp{viven juntas}.
\item \esp{No estoy de acuerdo : hay que luchar contra la discriminación en la escuela.} — « Il faut » impersonnel = \esp{hay que} + infinitif ; « lutter contre » = \esp{luchar contra}.
\item \esp{El respeto de las diferencias es la base de la convivencia.} — « Le respect de » = \esp{el respeto de} (on peut aussi dire \esp{el respeto a las diferencias}) ; « la base » = \esp{la base}.
\end{enumerate}
\textbf{B. Version}\\
« La tolérance ne signifie pas tout accepter sans réfléchir, mais écouter l'autre et reconnaître sa dignité. Dans une société multiculturelle, le dialogue et le respect sont les clés de la paix. Celui qui rejette celui qui est différent appauvrit sa propre identité. »\\
\textbf{Points de vigilance :}
\begin{itemize}
\item \esp{no ... sino} = « ne ... pas ... mais » (correction d'une négation) : ne pas traduire \esp{sino} par « sinon ».
\item \esp{quien} = « celui qui » (et non « qui ? » sans accent interrogatif).
\item \esp{el diferente} = « celui qui est différent » (adjectif substantivé).
\item \esp{empobrecer} : verbe en \emph{-ecer} $\rightarrow$ \esp{empobrece} au présent ; en français : « appauvrit ».
\end{itemize}
\textbf{Barème indicatif :} thème 3 $\times$ 2 pts ; version 6 pts ; respect des accents et de la ponctuation espagnole 2 pts.
\end{corrige}

\begin{corrige}[Exercice 11 — Expression écrite type Bac]
\textbf{Barème indicatif (sur 20) :} respect du sujet et de la longueur (80--100 mots) : 4 pts ; connecteurs (au moins 3) : 3 pts ; expressions d'accord/désaccord (au moins 2) : 3 pts ; correction grammaticale : 6 pts ; richesse du vocabulaire : 4 pts.

\textbf{Proposition de rédaction modèle (sujet 2) :}
\begin{quote}
\esp{En mi opinión, vivir juntos en un Camerún multicultural es sobre todo una riqueza. Primero, la diversidad de las lenguas y de las tradiciones hace de nuestro país un país único en África. Además, estoy de acuerdo con los que dicen que cada región aporta algo especial : la música, los platos y las fiestas nos enriquecen a todos. Sin embargo, no estoy de acuerdo con la idea de que la convivencia sea siempre fácil : los prejuicios existen y pueden provocar la discriminación. Por eso, la escuela y la familia deben enseñar el respeto y el diálogo. En conclusión, la diversidad es una riqueza, pero debemos cultivarla cada día con tolerancia.}
\end{quote}
\textbf{Traduction :} « À mon avis, vivre ensemble dans un Cameroun multiculturel est avant tout une richesse. D'abord, la diversité des langues et des traditions fait de notre pays un pays unique en Afrique. De plus, je suis d'accord avec ceux qui disent que chaque région apporte quelque chose de spécial : la musique, les plats et les fêtes nous enrichissent tous. Cependant, je ne suis pas d'accord avec l'idée que la cohabitation soit toujours facile : les préjugés existent et peuvent provoquer la discrimination. C'est pourquoi l'école et la famille doivent enseigner le respect et le dialogue. En conclusion, la diversité est une richesse, mais nous devons la cultiver chaque jour avec tolérance. » (environ 100 mots)

\textbf{Points de vigilance :}
\begin{itemize}
\item Compter ses mots : entre 80 et 100 ; un texte trop court est pénalisé.
\item Après \esp{no estoy de acuerdo con la idea de que}, employer le subjonctif : \esp{sea}, et non \esp{es}.
\item \esp{En mi opinión} : accent sur le \emph{ó} ; \esp{África} : accent sur le \emph{Á}.
\item Varier les connecteurs et annoncer clairement sa position dès la première phrase.
\end{itemize}
\end{corrige}

\newpage

\coursec{Fiche bilan}

\begin{popfigure}
\begin{tikzpicture}[mindmap, grow cyclic, every node/.style={concept, text=white, align=center, font=\small}, root concept/.append style={concept color=popblue, font=\bfseries\small}, level 1 concept/.append style={level distance=3.2cm, sibling angle=60, font=\footnotesize}, level 2 concept/.append style={level distance=2.2cm, font=\scriptsize}]
\node[root concept, align=center]{Identidad,\\tolerancia y\\convivencia}
  child[concept color=poporange]{ node{Lexique clé} 
    child[concept color=poporange!70]{ node[align=center]{identidad,\\tolerancia, respeto} }
    child[concept color=poporange!70]{ node[align=center]{prejuicio,\\discriminación} }
  }
  child[concept color=popgreen]{ node[align=center]{Accord /\\désaccord}
    child[concept color=popgreen!70]{ node[align=center]{estoy de\\acuerdo} }
    child[concept color=popgreen!70]{ node[align=center]{no estoy\\de acuerdo} }
  }
  child[concept color=poppurple]{ node[align=center]{Valeurs du\\vivre-ensemble}
    child[concept color=poppurple!70]{ node[align=center]{respeto,\\integración,\\diálogo} }
  }
  child[concept color=poppink]{ node[align=center]{Texte support\\« Vivir juntos\\en la diferencia »} }
  child[concept color=popgold]{ node[align=center]{Méthode\\commentaire}
    child[concept color=popgold!70]{ node[align=center]{intro, analyse,\\conclusion} }
  }
  child[concept color=popblue!70]{ node[align=center]{Production\\le Cameroun\\multiculturel} };
\end{tikzpicture}
\legende{Carte mentale de la leçon : identité, tolérance et vivre-ensemble.}
\end{popfigure}

\begin{bilan}
\textbf{1. Lexique clé à connaître par cœur}
\begin{itemize}
  \item \esp{la identidad} : l'identité ; \esp{la tolerancia} : la tolérance ; \esp{el respeto} : le respect.
  \item \esp{la convivencia} : le vivre-ensemble, la cohabitation ; \esp{la diferencia} : la différence.
  \item \esp{el prejuicio} : le préjugé ; \esp{la discriminación} : la discrimination.
  \item \esp{la integración} : l'intégration ; \esp{la diversidad cultural} : la diversité culturelle.
  \item \esp{respetar a los demás} : respecter les autres ; \esp{aceptar las diferencias} : accepter les différences.
\end{itemize}

\textbf{2. Exprimer l'accord et le désaccord}
\begin{center}
\begin{tabularx}{\linewidth}{|X|X|X|}\hline
\rowcolor{popblueL} Fonction & Español & Français \\ \hline
Accord & \esp{Estoy de acuerdo (contigo).} & Je suis d'accord (avec toi). \\ \hline
Accord fort & \esp{Tienes (toda la) razón.} & Tu as (tout à fait) raison. \\ \hline
Accord & \esp{Es verdad. / Exacto.} & C'est vrai. / Exact. \\ \hline
Désaccord & \esp{No estoy de acuerdo.} & Je ne suis pas d'accord. \\ \hline
Désaccord nuancé & \esp{Pienso que no. / No creo.} & Je pense que non. / Je ne crois pas. \\ \hline
Opinion & \esp{Pienso que / Creo que / Me parece que...} & Je pense que / Je crois que / Il me semble que... \\ \hline
\end{tabularx}
\end{center}

\textbf{3. Points de grammaire mobilisés}
\begin{itemize}
  \item \cle{ser / estar} : \esp{estar de acuerdo} (et non \esp{ser}) ; \esp{La tolerancia \textbf{es} importante} (caractéristique permanente).
  \item \cle{Verbes d'opinion + indicatif} : \esp{creo que}, \esp{pienso que}, \esp{me parece que} + indicatif ; à la forme négative (\esp{no creo que...}), l'espagnol préfère le subjonctif : \esp{No creo que sea fácil.} « Je ne crois pas que ce soit facile. »
  \item \cle{Adjectifs d'accord} : l'adjectif s'accorde avec le nom : \esp{una sociedad tolerante}, \esp{los prejuicios sociales}.
  \item \cle{Present de l'indicatif} pour exprimer des vérités générales : \esp{Vivir juntos exige respeto.} « Vivre ensemble exige du respect. »
\end{itemize}

\textbf{4. Méthode express : commenter un texte}
\begin{enumerate}
  \item \textbf{Introduction} : présenter le thème et l'idée générale du texte.
  \item \textbf{Analyse} : dégager 2 ou 3 idées principales, relever le lexique du thème (ici : tolérance, respect, diversité).
  \item \textbf{Opinion personnelle} : donner son avis avec \esp{pienso que...} et justifier.
  \item \textbf{Conclusion} : résumer et ouvrir (ex. le cas du Cameroun, pays multiculturel).
\end{enumerate}

\textbf{5. Phrases modèles pour la production écrite}
\begin{itemize}
  \item \esp{El Camerún es un país multicultural con más de doscientas lenguas.} « Le Cameroun est un pays multiculturel avec plus de deux cents langues. »
  \item \esp{La convivencia exige tolerancia y respeto de las diferencias.} « Le vivre-ensemble exige tolérance et respect des différences. »
  \item \esp{Debemos luchar contra los prejuicios y la discriminación.} « Nous devons lutter contre les préjugés et la discrimination. »
\end{itemize}
\end{bilan}

\begin{aretenir}[Checklist avant le Bac]
\begin{itemize}
  \item $\square$ Je connais le lexique : \esp{identidad, tolerancia, respeto, convivencia, prejuicio, diferencia, integración, discriminación}.
  \item $\square$ Je sais exprimer l'accord : \esp{estoy de acuerdo}, \esp{tienes razón}, \esp{es verdad}.
  \item $\square$ Je sais exprimer le désaccord : \esp{no estoy de acuerdo}, \esp{pienso que no}.
  \item $\square$ Je n'oublie pas : on dit \esp{\textbf{estar} de acuerdo}, jamais \esp{ser de acuerdo}.
  \item $\square$ Je sais donner mon opinion avec \esp{pienso que}, \esp{creo que}, \esp{me parece que}.
  \item $\square$ Après \esp{no creo que}, je pense au subjonctif : \esp{no creo que sea...}.
  \item $\square$ Je fais accorder les adjectifs : \esp{una sociedad tolerante}, \esp{ciudadanos respetuosos}.
  \item $\square$ Je sais résumer les idées principales du texte « Vivir juntos en la diferencia ».
  \item $\square$ Je connais les 4 étapes du commentaire de texte (introduction, analyse, opinion, conclusion).
  \item $\square$ Je sais écrire 4 ou 5 phrases sur le vivre-ensemble au Cameroun.
  \item $\square$ Je vérifie les accents : \esp{identidad}, \esp{integración}, \esp{discriminación}, \esp{está}.
  \item $\square$ Je n'oublie pas les signes \esp{¿} et \esp{¡} en début de phrase interrogative ou exclamative.
\end{itemize}
\end{aretenir}

\findecours

\end{document}
